\subsection{Bài tập tắc nghiệm}
\subsubsection{Câu hỏi trắc nghiệm 4 phương án}
\begin{dang}{Nhận biết hệ thức giữa cạnh và góc trong tam giác vuông. Áp dụng hệ thức tính độ dài cạnh trong tam giác vuông}
\end{dang}
\Opensolutionfile{ans}[ans/ans9K1-12-TN-D1]
%%==========Câu 1
\begin{ex}%%[9H4N2-1]
	Cho tam giác $ABC$ vuông tại $B$. Hệ thức nào sau đây là đúng?
	\choice
	{$BC=AB.\tan C$}
	{$BC=AC.\cot C$}
	{$BC=AB.\cos C$}
	{\True $BC=AC.\cos C$}
	\loigiai{
	Xét tam giác $ABC$ vuông tại $A$ ta có: $BC=AC.\cos C$}
\end{ex}
%%==========Câu 2
\begin{ex}%%[9H4N2-1]
	\immini[thm]
	{
	Cho tam giác $ABC$ vuông tại $A$ có $AB=c$; $AC=b$; $BC=a$. Chọn khẳng định \textbf{sai}?
	\choice
	{$c=a.\sin C=a.\cos B$}
	{$a^2=b^2+c^2$}
	{$b=a.\sin B=a.\cos C$}
	{\True $a=c.\tan B=c.\cot C$}
	}
	{
	% TODO: \usepackage{graphicx} required
	\includegraphics[scale=0.6]{C:/texstudio-man/Anh/{FD6BD6A4-6B9B-409F-887E-9FF3111B2332}}
	}
	\loigiai{
	Xét tam giác $ABC$ vuông tại $A$ ta có: $a=c.\tan B=c.\cot C$}
\end{ex}
%%==========Câu 3
\begin{ex}%%[9H4N2-1]
	Cho tam giác $ABC$ vuông tại $A$ có $BC=a,AC=b$, $AB=c$, $\widehat{ABC}=50^\circ$. Chọn khẳng định đúng?
	\choice
	{$b=c.\sin 50^\circ$}
	{$b=a.\tan 50^\circ$}
	{$b=c.\cot 50^\circ$}
	{\True $c=b.\cot 50^\circ$}
	\loigiai{
	Xét tam giác $ABC$ vuông tại $A$ ta có: $c=b.\cot 50^\circ$}
\end{ex}
%%==========Câu 4
\begin{ex}%%[9H4N2-1]
	Cho tam giác $ABC$ vuông tại $A$ có $BC=a$, $AC=b$, $AB=c$, $\widehat{ABC}=30^\circ$. Chọn khẳng định \textbf{sai}?
	\choice
	{$c=a.\cos30^\circ$}
	{\True $b=a.\tan 30^\circ$}
	{$c=b.\cot 30^\circ$}
	{$b=a.\sin 30^\circ$}
	\loigiai{
	Xét tam giác $ABC$ vuông tại $A$ ta có: $b=a.\tan 30^\circ$}
\end{ex}
%%==========Câu 5
\begin{ex}%[9H4H2-1]
	Cho tam giác $ABC$ vuông tại $A$, $\widehat{B}=60^\circ$, $AB=5\mathrm{~cm}$. Khi đó độ dài cạnh $AC$ là
	\choice
	{$\dfrac{5}{2}\mathrm{~cm}$}
	{$\dfrac{5\sqrt{3}}{2}\mathrm{~cm}$}
	{$\dfrac{5\sqrt{3}}{3}\mathrm{~cm}$}
	{\True $5\sqrt{3}\mathrm{~cm}$}
	\loigiai{
	Xét tam giác $ABC$ vuông tại $A$ ta có: $AC=AB.\tan 60^\circ$ $=5.\tan 60^\circ$ $=5\sqrt{3}\mathrm{~cm}$}
\end{ex}
%%==========Câu 6
\begin{ex}%[9H4H2-1]
	Cho tam giác $ABC$ vuông tại $A$, có $AB=6\mathrm{~cm}$ và biết $\tan B=\dfrac{5}{12}$. Độ dài cạnh $BC$ bằng
	\choice
	{$7,2\mathrm{~cm}$}
	{$5,2\mathrm{~cm}$}
	{\True $6,5\mathrm{~cm}$}
	{$3,8\mathrm{~cm}$}
	\loigiai{
	Xét tam giác $ABC$ vuông tại $A$ ta có: $\tan\alpha=\dfrac{5}{12}$ Suy ra $\alpha\approx 22^\circ 37'$ \\
	$BC=\dfrac{AB}{\cos\alpha}$ $=\dfrac{6}{\cos22^\circ 37'}$ $\approx 6,5\mathrm{~cm}$}
\end{ex}
%%==========Câu 7
\begin{ex}%[9H4H2-1]
	Cho $\triangle ABC$ vuông cân tại $A$, biết $BC=10\mathrm{~cm}$ Độ dài cạnh $AB$ là
	\choice
	{$20\sqrt{2}\mathrm{~cm}$}
	{$5\mathrm{~cm}$}
	{$10\sqrt{2}\mathrm{~cm}$}
	{\True $5\sqrt{2}\mathrm{~cm}$}
	\loigiai{
	Ta có $\triangle ABC$ vuông cân tại $A$ suy ra $\widehat{B}=45^\circ$ \\
	Ta có $\sin B=\dfrac{AC}{BC}$ suy ra $\dfrac{\sqrt{2}}{2}=\dfrac{AC}{10}$ suy ra $AC=5\sqrt{2}\mathrm{~cm}$. Suy ra $AB=5\sqrt{2}\mathrm{~cm}$}
\end{ex}
%%==========Câu 8
\begin{ex}%[9H4V2-1]
\immini[thm]
{
	Cho tam giác $ABC$ vuông tại $A$, $AH\perp BC$ ($H$ thuộc $BC$). Cho biết $AB:AC=4:3$ và $BC=15\mathrm{~cm}$ Độ dài đoạn thẳng $AH$ là
	\choice
	{$4,5\mathrm{~cm}$}
	{\True $7,2\mathrm{~cm}$}
	{$5,4\mathrm{~cm}$}
	{$ 6,9\mathrm{~cm}$}
}
{
	% TODO: \usepackage{graphicx} required
	\includegraphics[scale=0.6]{C:/texstudio-man/Anh/{28678482-2D40-414D-BB8F-DF3C602EBE57}}
}
	\loigiai{
	Xét tam giác $ABC$ vuông tại $A$ ta có $\dfrac{AB}{AC}=\dfrac{4}{3}$ suy ra $\tan C=\dfrac{4}{3}$ suy ra $\cos C=\dfrac{3}{5}$; $\sin C=\dfrac{4}{5}$ \\
	$AC=BC.\cos C$ $=9\mathrm{~cm}$; $AH=AC.\sin C$ $=7,2\mathrm{~cm}$}
\end{ex}
%%==========Câu 9
\begin{ex}%[9H4V2-5]
	\immini[thm]
	{
	Cho $\triangle ABC$ vuông tại $A$ có $AB=21\mathrm{~cm}$, $\widehat{C}=40^\circ$. Độ dài đường phân giác $BD$ là (làm tròn kết quả đến hàng phần mười)
	\choice
	{$10,88\mathrm{~cm}$}
	{$7,71\mathrm{~cm}$}
	{$18,67\mathrm{~cm}$}
	{\True $23,2\mathrm{~cm}$}
	}
	{
	% TODO: \usepackage{graphicx} required
	\includegraphics[scale=0.5]{C:/texstudio-man/Anh/{11D1CDE1-1D2B-4BD2-9F01-03065361CA57}}
	}
	\loigiai{
	Xét tam giác $ABC$ vuông tại $A$ ta có $\widehat{ABC}=90^\circ-\widehat{C}=50^\circ$ suy ra $\widehat{ABD}=\dfrac{\widehat{ABC}}{2}=25^\circ$ \\
	$\cos\widehat{ABD}=\dfrac{AB}{BD}$ suy ra $BD=\dfrac{AB}{\cos\widehat{ABD}}=\dfrac{21}{\cos 25^\circ}$ $\approx 23,2\mathrm{~cm}$}
\end{ex}
%%==========Câu 10
\begin{ex}%[9H4V2-5]
	\immini[thm]
	{
	Bạn An đứng cách chân một toàn nhà $10m$ nhìn thấy đỉnh một tòa nhà dưới một góc nâng $40^\circ$. Hỏi khi bạn An lùi ra xa sao cho An nhìn tòa nhà dưới một góc $35^\circ$ thì khoảng cách giữa An và tòa nhà là (Kết quả làm tròn tới hàng đơn vị)
	\choice
	{$10m$}
	{$11m$}
	{\True $12m$}
	{$13m$}
	}
	{
	% TODO: \usepackage{graphicx} required
	\includegraphics[scale=0.5]{C:/texstudio-man/Anh/{F2BB491A-D24D-4A3F-9091-DA4A7BE085BC}}
	}
	\loigiai{
	Chiều cao của cây là $MN$ \\
	Xét tam giác $MNP$ vuông tại $N$ ta có $MN=MP.\tan P$ $=10.\tan 40^\circ$.\\
	Xét tam giác $MNQ$ vuông tại $N$ ta có\\
	$QN=MN.\cot Q=10.\tan 40^\circ.\cot 35^\circ=10.\tan 40^\circ.\tan 55^\circ\approx 12m$ \\
}
\end{ex}
\Closesolutionfile{ans}
\indapan{10}{ans/ans9K1-12-TN-D1}
\begin{dang}{Tính số đo góc khi biết độ dài các cạnh trong tam giác vuông }
\end{dang}
\Opensolutionfile{ans}[ans/ans9K1-12-TN-D2]
%%==========Câu 11
\begin{ex}%[9H4H2-1]
	Cho $\triangle ABC$ vuông tại $A$ có $AB=6\mathrm{~cm}$, $AC=8\mathrm{~cm}$. Số đo $\widehat{B}$ (làm tròn đến độ) là
	\choice
	{\True $53^\circ$}
	{$54^\circ$}
	{$53^\circ 8'$}
	{$53^\circ 7'$}
	\loigiai{
	Xét tam giác $ABC$ vuông tại $A$ ta có $\tan B=\dfrac{AC}{AB}=\dfrac{8}{6}=\dfrac{4}{3}$ suy ra $\widehat{C}\approx 53^\circ$}
\end{ex}
%%==========Câu 12
\begin{ex}%[9H4H2-1]
	Cho $\triangle ABC$ vuông tại $A$ có $AB=4\mathrm{~cm}$, $BC=8\mathrm{~cm}$. Số đo $\widehat{C}$ (làm tròn đến độ) là
	\choice
	{\True $30^\circ$}
	{$60^\circ$}
	{$36^\circ$}
	{$53^\circ$}
	\loigiai{
	Xét tam giác $ABC$ vuông tại $A$ ta có $\sin C=\dfrac{AB}{BC}=\dfrac{4}{8}=\dfrac{1}{2}$ suy ra $\widehat{C}=30^\circ$}
\end{ex}
%%==========Câu 13
\begin{ex}%[9H4H2-1]
	Cho $\triangle ABC$ vuông tại $A$ có $AC=21\mathrm{~cm}$, $BC=28\mathrm{~cm}$. Số đo $\widehat{C}$ (làm tròn đến độ) là
	\choice
	{$30^\circ$}
	{\True $41^\circ$}
	{$36^\circ$}
	{$63^\circ$}
	\loigiai{
	Xét tam giác $ABC$ vuông tại $A$ ta có $\cos C=\dfrac{AC}{BC}$ $=\dfrac{21}{28}$ $=\dfrac{3}{4}$ suy ra $\widehat{C}\approx 41^\circ$}
\end{ex}
%%==========Câu 14
\begin{ex}%[9H4H2-1]
	Cho $\triangle MNP$ vuông tại $M$ biết $MN=9\mathrm{~cm}$, $NP=15\mathrm{~cm}$. Số đo $\widehat{P}$ (làm tròn tới độ) là
	\choice
	{$65^\circ$}
	{$63^\circ$}
	{$25^\circ$}
	{\True $37^\circ$}
	\loigiai{
	Xét $\triangle MNP$ vuông tại $M$ ta có $\sin P=\dfrac{MN}{NP}=\dfrac{9}{15}$ suy ra $\widehat{P}\approx 37^\circ$}
\end{ex}
%%==========Câu 15
\begin{ex}%[9H4H2-1]
\immini[thm]
{
	Cho $\triangle ABC$ vuông tại $A$ đường cao $AH$ có $AB=5\mathrm{~cm}$, $AH=4\mathrm{~cm}$. Số đo $\widehat{B}$ (làm tròn đến độ) là
	\choice
	{\True $53^\circ$}
	{$54^\circ$}
	{$53^\circ 8'$}
	{$53^\circ 7'$}
}
{
	% TODO: \usepackage{graphicx} required
	\includegraphics[scale=0.6]{C:/texstudio-man/Anh/{532AB141-D8A0-4977-8D8C-A6AD3DF03E98}}
}
	\loigiai{
	Xét tam giác $ABC$ vuông tại $A$ ta có $\cos B=\dfrac{AH}{AB}=\dfrac{4}{5}$ suy ra $\widehat{B}\approx 53^\circ$}
\end{ex}
%%==========Câu 16
\begin{ex}%[9H4H2-1]
	Cho $\triangle ABC$ vuông tại $A$ đường cao $AH$ có $AB=10\mathrm{~cm}$, $AH=6\mathrm{~cm}$. Số đo $\widehat{C}$ (làm tròn đến độ) là
	\choice
	{\True $53^\circ 8'$}
	{$36^\circ 52'$}
	{$36^\circ$}
	{$36^\circ 51'$}
	\loigiai{
	Xét tam giác $ABH$ vuông tại $H$ ta có $\sin B=\dfrac{AH}{AB}$ $=\dfrac{6}{10}$ $=\dfrac{3}{5}$. Suy ra $\widehat{B}\approx 36^\circ 52'$. Suy ra $\widehat{C}\approx 53^\circ 8'$}
\end{ex}
%%==========Câu 17
\begin{ex}%[9H4H2-1]
	Cho $\triangle ABC$ vuông tại $A$ đường cao $AH$ có $BH=2\mathrm{~cm}$, $HC=6\mathrm{~cm}$. Số đo $\widehat{B}$ (làm tròn đến độ) là
	\choice
	{$30^\circ$}
	{\True $60^\circ$}
	{$36^\circ$}
	{$63^\circ$}
	\loigiai{
	Xét tam giác $ABC$ vuông tại $A$ ta có $AH^2=BH.HC$ $=2.6=12$ suy ra $AH=\sqrt{12}=2\sqrt{3}$ \\
	Xét tam giác $ABH$ vuông tại $H$ ta có: $\tan B=\dfrac{AH}{BH}$ $=\dfrac{2\sqrt{3}}{2}$ $=\sqrt{3}$ suy ra $\widehat{B}=60^\circ$}
\end{ex}
%%==========Câu 18
\begin{ex}%[9H4V2-1]
\immini[thm]
{
	Cho $\triangle ABC$ vuông tại $A$ đường phân giác $BD$ biết $AD=2,5\mathrm{~cm}$, $DC=5\mathrm{~cm}$. Số đo $\widehat{DBC}$ là
	\choice
	{\True $30^\circ$}
	{$60^\circ$}
	{$15^\circ$}
	{$20^\circ$}
}
{
	% TODO: \usepackage{graphicx} required
	\includegraphics[scale=0.6]{C:/texstudio-man/Anh/{C1B5CA1D-7E5F-4053-9F71-AB8D358C7C23}}
}
	\loigiai{
	Vì $BD$ là tia phân giác của $\widehat{ABD}$ nên $\dfrac{AD}{DC}=\dfrac{BA}{BC}$ hay $\dfrac{BA}{BC}=\dfrac{2.5}{5}=\dfrac{1}{2}$ \\
	Xét tam giác $ABC$ vuông tại $A$ ta có: $\cos\widehat{ABC}=\dfrac{AB}{BC}=\dfrac{1}{2}$ suy ra $\widehat{ABC}=60^\circ$ \\
	$\widehat{BDC}=\dfrac{1}{2}\widehat{ABC}=30^\circ$}
\end{ex}
%%==========Câu 19
\begin{ex}%[9H4V2-1]
	Cho $\triangle ABC$ có $\widehat{ABC}=60^\circ$. Hình chiếu của cạnh $AB$ trên cạnh $BC$ có độ dài bằng $4\mathrm{~cm}$. Độ dài $AC$ dài gấp đôi $AB$. Số đo $\widehat{ACB}$ (làm tròn đến phút) là
	\choice
	{$25^\circ$}
	{\True $25^\circ 40'$}
	{$25^\circ 39'$}
	{$ 26^\circ$}
	\loigiai{
% TODO: \usepackage{graphicx} required
\begin{center}
	\includegraphics[scale=0.6]{C:/texstudio-man/Anh/{018D1AAC-6A3A-4540-ACA4-611688DD42A4}}
\end{center}
	$\triangle ABC$ vuông tại $A$, nên ta có $\cos B=\dfrac{BH}{AB}$ suy ra $AB=\dfrac{BH}{\cos B}$ $=\dfrac{4}{\cos60^\circ}$ $=8\mathrm{~cm}$.\\
	$AH=BH.\tan B$ $=4.\tan 60^\circ$ $=4\sqrt{3}\mathrm{~cm}$.\\
	$AC=2.AB$ $=2.8$ $=16\mathrm{~cm}$.\\
	$\sin C=\dfrac{AH}{AC}$ $=\dfrac{4\sqrt{3}}{16}$ $=\dfrac{\sqrt{3}}{4}$ suy ra $\widehat{C}\approx 25^\circ 40'$}
\end{ex}
%%==========Câu 20
\begin{ex}%[9H4V2-1]
\immini[thm]
{
	Cho hình thang $ABCD$ có $\widehat{A}=\widehat{D}=90^\circ$, $DC=30\mathrm{~cm}$, $\widehat{B}=60^\circ$, $CA\perp CB$. Diện tích hình thang $ABCD$ là
	\choice
	{\True $350\sqrt{3}\mathrm{cm^2}$}
	{$250\sqrt{3}\mathrm{cm^2}$}
	{$50\sqrt{3}\mathrm{cm^2}$}
	{$700\sqrt{3}\mathrm{cm^2}$}
}
{
	% TODO: \usepackage{graphicx} required
	\includegraphics[scale=0.6]{C:/texstudio-man/Anh/{1B1CC3EA-DC56-4FB3-9F88-36444A1C8419}}
}
	\loigiai{
	Kẻ $CH\perp AB$ tại $H$. Ta có $AHCD$ là hình chữ nhật nên $AH=DC=30$; $\widehat{CAB}=90^\circ-\widehat{B}$ $=90^\circ-60^\circ=30^\circ$ \\
	Xét $\triangle ACH$ vuông tại $H$ ta có $CH=AH.\tan 30^\circ$ $=30.\tan 30^\circ$ $=10\sqrt{3}\mathrm{~cm}$ \\
	Xét $\triangle BCH$ vuông tại $H$ ta có $BH=CH.\cot 60^\circ$ $=10\sqrt{3}.\tan 30^\circ$ $=10\mathrm{~cm}$ \\
	$AB=AH+BH$ $=30+10$ $=40\mathrm{~cm}$ \\
	$S_{ABCD}=\dfrac{1}{2}.CH.(AB+DC)$ $=\dfrac{1}{2}.10\sqrt{3}.(30+40)$ $=350\sqrt{3}\mathrm{cm^2}$ \\
	}
\end{ex}
\Closesolutionfile{ans}
\indapan{10}{ans/ans9K1-12-TN-D2}
\begin{dang}{Giải tam giác vuông}
\end{dang}
\Opensolutionfile{ans}[ans/ans9K1-12-TN-D3]
%%==========Câu 21
\begin{ex}%[9H4H2-1]
	Cho $\triangle ABC$ vuông tại $A$, có $BC=9\mathrm{~cm}$ và $\widehat{C}=53^\circ$. Số đo cạnh $AB$ (kết quả làm tròn đến hàng phần chục) là
	\choice
	{$7,1\mathrm{~cm}$}
	{$5,4\mathrm{~cm}$}
	{\True $5,2\mathrm{~cm}$}
	{$12,6\mathrm{~cm}$}
	\loigiai{
	Xét $\triangle ABC$ vuông tại $A$, nên ta có $\sin C=\dfrac{AB}{BC}$ suy ra $AB=BC.\sin C$ $=9.\sin 53^\circ$ $\approx 5,2\mathrm{~cm}$}
\end{ex}
%%==========Câu 22
\begin{ex}%[9H4H2-1]
	Cho $\triangle ABC$ vuông tại $A$ có $AB=9\mathrm{~cm}$, $AC=12\mathrm{~cm}$. Số đo $\widehat{C}$ (làm tròn đến độ) là
	\choice
	{$53^\circ$}
	{$90^\circ$}
	{$47^\circ$}
	{\True $37^\circ$}
	\loigiai{
	Xét $\triangle ABC$ vuông tại $A$, ta có $\tan C=\dfrac{AB}{AC}=\dfrac{9}{12}$ suy ra $\widehat{C}\approx 37^\circ$}
\end{ex}
%%==========Câu 23
\begin{ex}%[9H4N2-6]
	Cho $\triangle ABC$ vuông tại $B$. Biết $AC=16\mathrm{~cm}$, $AB=12\mathrm{~cm}$. Giải tam giác $\triangle ABC$ ta cần tính các yếu tố là
	\choice
	{Cạnh $BC$, góc $B$, góc $C$}
	{\True Cạnh $BC$, góc $A$, góc $C$}
	{Cạnh $AB$, góc $A$, góc $C$}
	{Cạnh $AC$, góc $A$, góc $C$}
	\loigiai{
	Cạnh $BC$, góc $\widehat{A};\widehat{C}$}
\end{ex}
%%==========Câu 24
\begin{ex}%[9H4H2-1]
\immini[thm]
{
	Cho hình vẽ. Độ dài $x$ (làm tròn đến hàng phần trăm) là
	\choice
	{$0,87m$}
	{$0,18m$}
	{$0,53m$}
	{\True $1,46\mathrm{~cm}$}
}
{
	% TODO: \usepackage{graphicx} required
	\includegraphics[scale=0.5]{C:/texstudio-man/Anh/{B0F2748F-569B-4EF5-A470-671A4F9288DA}}
}
	\loigiai{
	Ta có $\cos70^\circ=\dfrac{0,5}{x}$ Suy ra $x=\dfrac{0,5}{\cos70^\circ}\approx 1,46m$}
\end{ex}
%%==========Câu 25
\begin{ex}%[9H4H2-2]
Cho $\triangle ABC$ vuông tại $A$ biết $AB=12\mathrm{~cm}$, $BC=1,5\mathrm{~dm}$ Độ dài $AC$ và số đo $\widehat{B}$ là (kết quả số đo góc làm tròn đến phút)
\choice
{$AC=8\mathrm{~cm}$; $\widehat{B}\approx 36^\circ 52'$}
{\True $AC=9\mathrm{~cm}$; $\widehat{B}\approx 36^\circ 52'$}
{$AC=9\mathrm{~cm}$; $\widehat{B}\approx 37^\circ 52'$}
{$AC=9\mathrm{~cm}$; $\widehat{B}\approx 36^\circ 55'$}
\loigiai{
Đổi $1,5\mathrm{~dm}=15\mathrm{~cm}$ \\
Xét $\triangle ABC$ vuông tại $A$, nên ta có $AC=\sqrt{BC^2-AB^2}$ $=\sqrt{15^2-12^2}$ $=9m$ \\
$\sin B=\dfrac{AC}{BC}=\dfrac{9}{15}$ Suy ra $\widehat{B}\approx 36^\circ 52'$}
\end{ex}
%%==========Câu 26
\begin{ex}%[9H4H2-1]
	Cho $\triangle ABC$ vuông tại $A$ biết $AC=10\mathrm{~cm}$, $\widehat{C}=30^\circ$. Độ dài cạnh $AB$ là
	\choice
	{$AB=\dfrac{5\sqrt{3}}{3}\mathrm{~cm}$}
	{$AB=5\mathrm{~cm}$}
	{$AB=10\mathrm{~cm}$}
	{\True $AB=\dfrac{10\sqrt{3}}{3}\mathrm{~cm}$}
	\loigiai{
	Xét $\triangle ABC$ vuông tại $A$, nên ta có $AB=AC.\tan 30^\circ$ $=10.\tan 30^\circ$ $=\dfrac{10\sqrt{3}}{3}\mathrm{~cm}$}
\end{ex}
%%==========Câu 27
\begin{ex}%[9H4H2-1]
	Cho $\triangle ABC$ vuông cân tại $A$ biết $BC=\sqrt{2}\mathrm{~dm}$. Độ dài đường cao $AH$ là
	\choice
	{$1\mathrm{~dm}$}
	{$\sqrt{2}\mathrm{~dm}$}
	{$\dfrac{\sqrt{2}}{4}\mathrm{~dm}$}
	{\True $\dfrac{\sqrt{2}}{2}\mathrm{~dm}$}
	\loigiai{
	Kẻ đường cao $AH$ vì $\triangle ABC$ vuông cân tại $A$ nên $BH=\dfrac{BC}{2}=\dfrac{\sqrt{2}}{2}$, $\widehat{B}=45^\circ$ \\
	$AH=BH.\tan 45^\circ$ $=\dfrac{\sqrt{2}}{2}.\tan 45^\circ$ $=\dfrac{\sqrt{2}}{2}\mathrm{~dm}$}
\end{ex}
%%==========Câu 28
\begin{ex}%[9H4V2-1]
\immini[thm]
{
	Độ dài $x$ và $ y$ trong hình vẽ sau lần lượt là
	\choice
	{\True $x=\dfrac{28\sqrt{3}}{3}\mathrm{~cm}$; $y=28\sqrt{2}\mathrm{~cm}$}
	{$x=\dfrac{13}{3}\mathrm{~cm}$; $y=2\sqrt{2}\mathrm{~cm}$}
	{$x=47\mathrm{~cm}$; $y=56\mathrm{~cm}$}
	{$x=47\mathrm{~cm}$; $y=55\mathrm{~cm}$}
}
{
	% TODO: \usepackage{graphicx} required
	\includegraphics[scale=0.45]{C:/texstudio-man/Anh/{31172B2F-5C51-4B8A-82B0-DB224A762766}}
}
	\loigiai{
	Xét $\triangle MNP$ vuông tại $N$ ta có $PN=MN.\tan 30^\circ$ suy ra $x=28.\tan 30^\circ$ $=\dfrac{28\sqrt{3}}{3}\mathrm{~cm}$ \\
	Xét $\triangle MNQ$ vuông tại $N$ ta có $\cos\widehat{MNQ}=\dfrac{NM}{MQ}$ suy ra $y=MQ$ $=\dfrac{28}{\cos45^\circ}$ $=28\sqrt{2}\mathrm{~cm}$}
\end{ex}
%%==========Câu 29
\begin{ex}%[9H4V2-1]
\immini[thm]
{
	Cho tam giác $ABC$ vuông tại $A$ có $BC=8\mathrm{~cm}$, $\cos C=\dfrac{\sqrt{2}}{2}$. Gọi $D$ là trung điểm $AC$ kẻ $DH\perp BC$. Độ dài $HC$ là
	\choice
	{$3\sqrt{10}\mathrm{~cm}$}
	{\True $1\mathrm{~cm}$}
	{$2\sqrt{10}\mathrm{~cm}$}
	{$4\sqrt{5}\mathrm{~cm}$}
}
{
	% TODO: \usepackage{graphicx} required
	\includegraphics[scale=0.6]{C:/texstudio-man/Anh/{A22C6733-942F-4C6F-8464-D39D6CF4A71D}}
}
	\loigiai{
	Xét tam giác $ABC$ vuông tại $A$ ta có $AC=BC.\cos C=8.\dfrac{\sqrt{2}}{4}=2\sqrt{2}\mathrm{~cm}$.\\
	$D$ là trung điểm $AC$ nên $DC=\dfrac{AC}{2}=\dfrac{2\sqrt{2}}{2}=\sqrt{2}\mathrm{~cm}$.\\
	Xét tam giác $DHC$ vuông tại $H$ ta có $HC=DC.\cos C=\sqrt{2}.\dfrac{\sqrt{2}}{2}=1\mathrm{~cm}$}
\end{ex}
%%==========Câu 30
\begin{ex}%[9H4V2-1]
\immini[thm]
{
Cho $\triangle ABC$ vuông tại $A$. Biết $AC=16\mathrm{~cm}$, $AB=12\mathrm{~cm}$. Các đường phân giác trong và ngoài của $\widehat{B}$ cắt $AC$ tại $D,E$. Độ dài cạnh $ED$ là
\choice
{$28\mathrm{~cm}$}
{\True $30\mathrm{~cm}$}
{$32\mathrm{~cm}$}
{$34\mathrm{~cm}$}
}
{
% TODO: \usepackage{graphicx} required
\includegraphics[scale=0.45]{C:/texstudio-man/Anh/{B41FEFE2-0362-4E20-8168-BD337AA947F8}}
}
\loigiai{
Xét $\triangle ABC$ vuông tại $A$, nên ta có $BC=\sqrt{AB^2+AC^2}=\sqrt{16^2+12^2}=20\mathrm{~cm}$.\\
Vì $BD$ là tia phân giác trong của $\widehat{B}$ nên $\dfrac{AD}{AB}=\dfrac{DC}{BC}$ $=>\dfrac{AD}{12}=\dfrac{DC}{20}=\dfrac{AD+DC}{12+20}=\dfrac{16}{32}=\dfrac{1}{2}$ nên $AD=6\mathrm{~cm},DC=10\mathrm{~cm}$.\\
$BD$ và $BE$ là phân giác trong và phân giác ngoài của góc B nên $BD\perp BE$ \\
Xét $\triangle BED$ vuông tại $B$ có $AE=\dfrac{AB^2}{AD}=\dfrac{12^2}{6}=24$; $ED=EA+AD=24+6=30\mathrm{~cm}$.\\
}
\end{ex}
\Closesolutionfile{ans}
\indapan{10}{ans/ans9K1-12-TN-D3}
\begin{dang}{Tính độ dài cạnh, số đo góc trong tam giác không vuông}
\end{dang}
\Opensolutionfile{ans}[ans/ans9K1-12-TN-D4]
%%==========Câu 31
\begin{ex}%[9H4H2-1]
\immini[thm]
{
	Cho hình vẽ sau. Hệ thức nào dưới đây đúng?
	\choice
	{\True $HC=BC.\sin B$}
	{$HC=BC.\cos B$}
	{$HC=BC.\tan B$}
	{$HC=BC.\cot B$}
}
{
	% TODO: \usepackage{graphicx} required
	\includegraphics[scale=0.6]{C:/texstudio-man/Anh/{C31DF424-2478-4F42-B682-F040E462E9E4}}
}
	\loigiai{
	Xét tam giác vuông $BHC\left(\widehat{BHC}=90^\circ\right)$, ta có: $HC=BC.\sin B$}
\end{ex}
%%==========Câu 32
\begin{ex}%[9H4H2-1]
	\immini[thm]
	{
	Cho hình vẽ sau. Hệ thức nào dưới đây đúng?
	\choice
	{$AB=AD.\cos B$}
	{$AB=AD.\sin B$}
	{\True $AB=\dfrac{AD}{\sin B}$}
	{$AB=\dfrac{AD}{\cos B}$}
	}
	{
	% TODO: \usepackage{graphicx} required
	\includegraphics[scale=0.6]{C:/texstudio-man/Anh/{872B6C8F-3C8C-4C8C-A2CF-6901F91D3C19}}
	}
	\loigiai{
	Xét tam giác vuông $ABD\left(\widehat{A}=90^\circ\right)$, ta có: $AD=AB.\sin B$ hay $AB=\dfrac{AD}{\sin B}$}
\end{ex}
%%==========Câu 33
\begin{ex}%[9H4H2-1]
\immini[thm]
{
	Cho hình vẽ sau. Hệ thức nào sau đây đúng?
	\choice
	{$AC=BC.\cos B$}
	{$BD=AD.\sin B$}
	{$BD=\dfrac{AD}{\sin B}$}
	{\True $BD=\dfrac{AD}{\tan B}$}
}
{
	% TODO: \usepackage{graphicx} required
	\includegraphics[scale=0.6]{C:/texstudio-man/Anh/{3BEBAA26-E934-4A5F-9C63-62EEED93F7A1}}
}
	\loigiai{
	Xét tam giác vuông $ABD\left(\widehat{D}=90^\circ\right)$, ta có: $AD=BD.\tan B$. Hay $BD=\dfrac{AD}{\tan B}$}
\end{ex}
%%==========Câu 34
\begin{ex}%[9H4H2-1]
\immini[thm]
{
	Cho hình vẽ sau. Hệ thức nào dưới đây đúng?
	\choice
	{$BH=BC.\tan B$}
	{$BH=HC.\cos B$}
	{$BC=\dfrac{BH}{\tan B}$}
	{\True $BD=AB.\cos B$}
}
{
	% TODO: \usepackage{graphicx} required
	\includegraphics[scale=0.6]{C:/texstudio-man/Anh/{702DF16B-52D2-45FD-BEAC-4DDAC9A082B0}}
}
	\loigiai{
	Xét $\triangle HBC$ vuông tại $H$ ta có:\\
	$\tan B=\dfrac{HC}{BH}$ hay $BH=\dfrac{HC}{\tan B}$. Suy ra A và C sai.\\
	$\cos B=\dfrac{BH}{BC}$ hay $BH=BC.\cos B$. Suy ra B sai\\
	Xét $\triangle ABD$ vuông tại $D$ ta có: $\cos B=\dfrac{BD}{AB}$ hay $BD=AB.\cos B$. Suy ra D đúng}
\end{ex}
%%==========Câu 35
\begin{ex}%[9H4V2-3]
Cho tam giác $ABC$ có $AB=4\mathrm{~cm}$, $AC=3,5\mathrm{~cm}$, $CH\perp AB$, $\widehat{A}=40^\circ$
Diện tích tam giác $ABC$ (kết quả làm tròn đến chữ số thập phân thứ nhất) là \ldots.
\choice
{$5,0\mathrm{cm^2}$}
{\True $4,4\mathrm{cm^2}$}
{$3,0\mathrm{cm^2}$}
{$3,5\mathrm{cm^2}$}
\loigiai{
Xét tam giác $HAC$ vuông tại $H$, ta có: $CH=AC.\sin\widehat{HAC}=3,5.\sin 40^\circ\approx 2,2\mathrm{~cm}$ \\
Diện tích của tam giác $ABC$ là: $S=\dfrac{1}{2}AB.CH\approx\dfrac{1}{2}.4.2,2\approx 4,4\mathrm{cm^2}$}
\end{ex}
%%==========Câu 36
\begin{ex}%[9H4H2-3]
Cho tam giác $ABC$ có $\widehat{B}=78^\circ$, đường cao $AH=10\mathrm{~cm}$ Độ dài của cạnh $AB$ của tam giác $ABC$ (kết quả làm tròn đến hàng phần mười) bằng.
\choice
{$13,1\mathrm{~cm}$}
{$\text{7,7cm}$}
{$\text{6,4cm}$}
{\True $\text{10,2cm}$}
\loigiai{
Xét tam giác $HAB$ vuông tại $H$, ta có: $\sin B=\dfrac{AH}{AB}$ hay $AB=\dfrac{AH}{\sin B}=\dfrac{10}{\sin 78^\circ}\approx 10,2\mathrm{~cm}$}
\end{ex}
%%==========Câu 37
\begin{ex}%[9H4H2-1]
\immini[thm]
{
Cho tam giác $DEF$ có $DE=7\mathrm{~cm}$, $\widehat{D}=40^\circ$, $\widehat{F}=58^\circ$. Kẻ đường cao $EI$ của tam giác đó. Độ dài của đường cao $EI$ (kết quả làm tròn đến chữ số thập phân thứ nhất) bằng \ldots.
\choice
{\True $4,5\mathrm{~cm}$}
{$\text{5,4cm}$}
{$\text{5,9cm}$}
{$\text{5,6cm}$}
}
{
% TODO: \usepackage{graphicx} required
\includegraphics[scale=0.6]{C:/texstudio-man/Anh/{358A31B8-9801-4E72-A670-F939F12C34FC}}
}
\loigiai{
Xét tam giác $DEI$ vuông tại $I$, ta có: $EI=ED.\sin D=7.\sin 40^\circ\approx 4,5\mathrm{~cm}$}
\end{ex}
%%==========Câu 38
\begin{ex}%[9H4H2-3]
	Cho tam giác $ABC$ có $AB=16\mathrm{~cm}$, $AC=14\mathrm{~cm}$ và $\widehat{B}=60^\circ$. Độ dài của $BC$ là \ldots.
	\choice
	{\True $10\mathrm{~cm}$}
	{$\text{11}\mathrm{~cm}$}
	{$9\mathrm{~cm}$}
	{$\text{12}\mathrm{~cm}$}
	\loigiai{
% TODO: \usepackage{graphicx} required
\begin{center}
	\includegraphics[scale=0.6]{C:/texstudio-man/Anh/{1DC963EE-2354-432D-81A7-8580828BF8D6}}
\end{center}
	Kẻ đường cao $AH$.\\
	+) Xét tam giác $ABH$ vuông tại $H$, ta có: $BH=AB.\cos B=AB.\cos 60^\circ=16\cdot\dfrac{1}{2}=8\mathrm{~cm}$.\\
	$AH=AB.\sin B=AB.\sin 60^\circ=16\cdot\dfrac{\sqrt{3}}{2}=8\sqrt{3}\mathrm{~cm}$.\\
	+) Xét tam giác $ACH$ vuông tại $H$, ta có: $HC^2=AC^2-AH^2=14^2-\left(8\sqrt{3}\right)^2=4$ \\
	Suy ra $HC=2\mathrm{~cm}$ \\
	Vậy $BC=BH+HC=8+2=10\mathrm{~cm}$}
\end{ex}
%%==========Câu 39
\begin{ex}%[9H4V2-3]
	Cho tam giác $ABC$ có $\widehat{B}=60^\circ$, $\widehat{C}=55^\circ$, $AC=3,5\mathrm{~cm}$. Diện tích tam giác $ABC$ (làm tròn đến hàng đơn vị) là \ldots.
	\choice
	{$\text{4c}\mathrm{m^2}$}
	{\True $5\mathrm{cm^2}$}
	{$7\mathrm{cm^2}$}
	{$\text{8c}\mathrm{m^2}$}
	\loigiai{
% TODO: \usepackage{graphicx} required
\begin{center}
	\includegraphics[scale=0.6]{C:/texstudio-man/Anh/{58BE9E69-D3F3-4BB8-A4C8-B9698B723EFD}}
\end{center}
	Kẻ đường cao $AD$\\
	+) Xét tam giác $ACD$ vuông tại $D$, ta có: $AD=AC.\sin C=3,5.\sin 50^\circ\approx 2,68\mathrm{~cm}$ \\
	$CD=AC.\cos C=3,5.\cos 50^\circ\approx 2,25\mathrm{~cm}$\\
	+) Xét tam giác $ABD$ vuông tại $D$, ta có: $\tan B=\dfrac{AD}{BD}$ hay $BD=\dfrac{AD}{\tan B}\approx\dfrac{2,68}{\tan 60^\circ}\approx 1,55\mathrm{~cm}$.\\
	Suy ra $BC=BD+DC\approx 1,55+2,55\approx 3,8\mathrm{~cm}$ \\
	Vậy $S_{ABC}=\dfrac{1}{2}\cdot AD\cdot BC\approx\dfrac{1}{2}\cdot 1,55\cdot 3,8\approx 5\mathrm{cm^2}$}
\end{ex}
%%==========Câu 40
\begin{ex}%[9H4V2-3]
	Tứ giác $ABCD$ có hai đường chéo cắt nhau tại $O$. Biết $\widehat{AOD}=70^\circ$; $AC=5,3\mathrm{~cm}$; $BD=4,0\mathrm{~cm}$. Diện tích tứ giác $ABCD$ (làm tròn đến số thập phân thứ nhất) là \ldots.
	\choice
	{$9,9\mathrm{cm^2}$}
	{$8,9\mathrm{cm^2}$}
	{$10,1\mathrm{cm^2}$}
	{\True $10,0\mathrm{cm^2}$}
	\loigiai{
% TODO: \usepackage{graphicx} required
\begin{center}
	\includegraphics[scale=0.6]{C:/texstudio-man/Anh/{B7F56D7D-2D1E-437F-9C26-9D3E3D325274}}
\end{center}
	Kẻ $BH\perp AC$ tại $H$ và $DK\perp AC$ tại $K$.\\
	Xét tam giác $HOB$ vuông tại H ta có $BH=OB.\sin O_1$ \\
	Xét tam giác $KOD$ vuông tại $K$ ta có $DK=OD.\sin O_2$ \\
	Mà $\widehat{O_1}=\widehat{O_2}$ (đối đỉnh)\\
	$\Rightarrow{S_{ABCD}}=S_{ABC}+S_{ADC}=\dfrac{1}{2}AC\left(BH+DK\right)=\dfrac{1}{2}AC\left(OB+OD\right)\sin O_1=\dfrac{1}{2}\cdot AC\cdot BD\cdot\sin O_1$ \\
	$=\dfrac{1}{2}5,34.\sin 70^\circ\approx 10,0\mathrm{cm^2}$.\\
	Vậy $S_{ABCD}\approx 10,0\mathrm{cm^2}$ \\
	}
\end{ex}
\Closesolutionfile{ans}
\indapan{10}{ans/ans9K1-12-TN-D4}
\begin{dang}{Chứng tỏ được hệ thức liên hệ giữa các cạnh}
\end{dang}
\Opensolutionfile{ans}[ans/ans9K1-12-TN-D5]
%%==========Câu 41
\begin{ex}%[9H4H2-4]
	Cho tam giác $ABC$ vuông tại $A$, có $\widehat{C}=30^\circ$ và $BC=a$, ta có:
	\choice
	{\True $AB=\dfrac{a}{2}$}
	{$AB=\dfrac{\sqrt{3}}{2}$}
	{$AB=\dfrac{a\sqrt{3}}{2}$}
	{$AB=\dfrac{1}{2}$}
	\loigiai{
	Xét $\triangle ABC$ vuông tại $A$, ta có: $AB=BC.\sin C=a.\sin 30^\circ=a\cdot\dfrac{1}{2}=\dfrac{a}{2}$}
\end{ex}
%%==========Câu 42
\begin{ex}%[9H4H2-4]
	Cho tam giác $ABC$ vuông tại $A$, có $\widehat{C}=30^\circ$ và $BC=a$, ta có:
	\choice
	{$AC=\dfrac{a}{2}$}
	{$AC=\dfrac{\sqrt{3}}{2}$}
	{\True $AC=\dfrac{a\sqrt{3}}{2}$}
	{$AC=\dfrac{1}{2}$}
	\loigiai{
	Xét $\triangle ABC$ vuông tại $A$, ta có: $AC=BC.\cos C=a\cdot\dfrac{\sqrt{3}}{2}=\dfrac{a\sqrt{3}}{2}$}
\end{ex}
%%==========Câu 43
\begin{ex}%[9H4H2-6]
	Cho tam giác $ABC$ vuông tại $A$, ta có:
	\choice
	{\True $\dfrac{AC}{AB}=\dfrac{\sin B}{\sin C}$}
	{$\dfrac{AB}{AC}=\dfrac{\sin B}{\sin C}$}
	{$\dfrac{AC}{AB}=\dfrac{\cos B}{\cos C}$}
	{$\dfrac{AB}{BC}=\dfrac{\cos B}{\cos C}$}
	\loigiai{
	Xét $\triangle ABC$ vuông tại $A$, ta có: $\sin B=\dfrac{AC}{BC}$, $\sin C=\dfrac{AB}{BC}$, $\cos B=\dfrac{AB}{BC}$, $\cos C=\dfrac{AC}{BC}$ \\
	Suy ra: $\dfrac{\sin B}{\sin C}=\dfrac{AC}{BC}:\dfrac{AB}{BC}=\dfrac{AC}{AB}$; $\dfrac{\cos B}{\cos C}=\dfrac{AB}{BC}:\dfrac{AC}{BC}=\dfrac{AB}{AC}$}
\end{ex}
%%==========Câu 44
\begin{ex}%[9H4H2-6]
	Cho tam giác $ABC$ vuông tại $A$ có $\widehat{B}=\alpha$.
	Trong các khẳng định sau, khẳng định nào \textbf{sai}?
	\choice
	{$\tan\alpha=\dfrac{\sin\alpha}{\cos\alpha}$}
	{\True $\tan\alpha=\dfrac{\cos\alpha}{\sin\alpha}$}
	{$\cot\alpha=\dfrac{\cos\alpha}{\sin\alpha}$}
	{$\tan\alpha.\cot\alpha=1$}
	\loigiai{Xét $\triangle ABC$ vuông tại $A$, ta có:\\ $\tan\alpha=\dfrac{AC}{AB}=\dfrac{AC}{BC}\cdot\dfrac{BC}{AB}=\dfrac{AC}{BC}:\dfrac{AB}{BC}=\dfrac{\sin\alpha}{\cos\alpha}$ \\
	(Vì $\sin\alpha=\dfrac{AC}{BC}$, $\cos\alpha=\dfrac{AB}{BC}$) (A đúng, B sai)\\
	Mà $\cot\alpha=\dfrac{1}{\tan\alpha}$ nên $\cot\alpha=\dfrac{\cos\alpha}{\sin\alpha}$ (C đúng)\\
	$\tan\alpha.\cot\alpha=\dfrac{AC}{AB}\cdot\dfrac{AB}{AC}=1$ (D đúng)}
\end{ex}
%%==========Câu 45
\begin{ex}%[9H4H2-4]
	Cho tam giác $ABC$ vuông tại $A$, đường cao $AH$, biết $BC=a$ và $\widehat{B}=\beta$.
	Độ dài đoạn $AH$ tính theo $a$ và $\beta$ là
	\choice
	{\True $AH=a.\sin\beta.\cos \beta$}
	{$AH=a.\sin\beta$}
	{$AH=a.\cos \beta$}
	{$AH=a.\tan \beta$}
	\loigiai{
	Xét $\triangle ABC$ vuông tại $A$, ta có: $AB=BC.\cos B=a.\cos \beta$ (1)\\
	Xét $\triangle AHB$ vuông tại $H$, ta có: $AH=AB.\sin B=AB.\sin\beta$ (2)\\
	Từ (1) và (2) ta có: $AH=a.\sin\beta.\cos \beta$}
\end{ex}
%%==========Câu 46
\begin{ex}%[9H4H2-4]
	Cho tam giác $ABC$, hai đường cao $BH$, $CK$. Trong các khẳng định sau, khẳng định nào đúng?
	\choice
	{Nếu $AB>AC$ thì $BH<CK$}
	{Nếu $AB<AC$ thì $BH>CK$}
	{\True Nếu $AB>AC$ thì $BH>CK$}
	{Cả A, B và C đều đúng}
	\loigiai{
% TODO: \usepackage{graphicx} required
\begin{center}
	\includegraphics[scale=0.6]{C:/texstudio-man/Anh/{AE57A515-F1CB-4C8A-8363-335ED2E363E8}}
\end{center}
	Giả sử $AB>AC$.\\
	Xét tam giác $AHB$ vuông tại $H$, ta có: $BH=AB.\sin A$ (1)\\
	Xét tam giác $AKC$ vuông tại $K$, ta có: $CK=AC.\sin A$ (2)\\
	Từ (1) và (2) suy ra: $\dfrac{BH}{CK}=\dfrac{AB.\sin A}{AC.\sin A}=\dfrac{AB}{AC}>1$ (Vì $\sin A>0$ và $AB>AC$)\\
	Do đó $BH>CK$.\\
	Vậy nếu $AB>AC$ thì $BH>CK$}
\end{ex}
%%==========Câu 47
\begin{ex}%[9H4H2-4]
	Cho tam giác đều $ABC$ có $AB=2a$. Kẻ đường cao $AH$.
	Trong các khẳng định sau, khẳng định nào \textbf{sai}?
	\choice
	{$BH=a$}
	{$AH=a\sqrt{3}$}
	{$BC=2a$}
	{\True $AC=a$}
	\loigiai{
	Xét $\triangle ABC$ đều, ta có $AH$ là đường cao nên $AH$ cũng là đường trung tuyến.\\
	Suy ra $BH=\dfrac{BC}{2}=\dfrac{2a}{2}=a$ (A đúng)\\
	Xét $\triangle ABH$ vuông tại $H$, ta có: $AH^2=AB^2-BH^2=(2a)^2-a^2=3a^2$ hay $AH=a\sqrt{3}$ (B đúng)\\
	Xét $\triangle ABC$ đều, ta có: $AB=BC=AC=2a$ (C đúng, D sai)}
\end{ex}
%%==========Câu 48
\begin{ex}%[9H4V2-4]
	Cho tam giác $ABC$ vuông cân tại $A$ $\left(AB=AC=a\right)$. Phân giác của góc $B$ cắt $AC$ tại $D$. Trong các khẳng định sau, khẳng định nào đúng?
	\choice
	{\True $AD=a.\cos 22,5^\circ$; $DC=a-a.\cos 22,5^\circ$}
	{$AD=a.\sin 22,5^\circ$; $DC=a-a.\sin 22,5^\circ$}
	{$AD=a.\tan 22,5^\circ$; $DC=a-a.\tan 22,5^\circ$}
	{$AD=a.\cot 22,5^\circ$; $DC=a-a.\cot 22,5^\circ$}
	\loigiai{
% TODO: \usepackage{graphicx} required
\begin{center}
	\includegraphics[scale=0.6]{C:/texstudio-man/Anh/{F743D5F0-4141-43A5-9B8F-FBEE4E98F09E}}
\end{center}
	Xét $\triangle ABC$ vuông cân tại $A$, ta có: $\widehat{B}=\widehat{C}=45^\circ$.\\
	Vì $BD$ là tia phân giác của góc $B$, nên: $\widehat{ABD}=\widehat{DBC}=\dfrac{1}{2}\widehat{ABC}=\dfrac{45^\circ}{2}=22,5^\circ$ \\
	Xét $\triangle ABD$ vuông tại $A$, ta có: $AD=AB.\tan ABD=a.\tan 22,5^\circ$ \\
	Ta có $AD+DC=AC$, suy ra: $DC=AC-AD=a-a.\tan 22,5^\circ$}
\end{ex}
%%==========Câu 49
\begin{ex}%[9H4V2-4]
	Cho tam giác $ABC$ vuông tại $A$, đường cao $AH$. Đặt $BC=a$. Trong các khẳng định sau, khẳng định nào \textbf{sai}?
	\choice
	{$AH=a.\sin B.\cos B$}
	{$BH=a.\cos^2B$}
	{$CH=a.\sin^2B$}
	{\True $BC.BH=a.\sin^2B.\cos^2B$}
	\loigiai{
% TODO: \usepackage{graphicx} required
\begin{center}
	\includegraphics[scale=0.6]{C:/texstudio-man/Anh/{CD47E82E-97F3-40C7-B637-C89DEBFEC997}}
\end{center}
	Xét tam giác $ABC$ vuông tại $A$, ta có: $AC=BC.\sin B=a.\sin B$ \\
	$AB=BC.\cos B=a.\cos B$\\
	Xét tam giác $AHB$ vuông tại $H$, $\widehat{HAC}=\widehat{B}$, ta có:\\
	$AH=AC.\cos B=a.\sin B.\cos B$ (A đúng)\\
	$BH=AB.\cos B=a.\cos B.\cos B=a.\cos^2B$ (B đúng)\\
	Xét tam giác $AHC$ vuông tại $H$, ta có:\\
	$HC=AC.\sin B=BC.\sin B.\sin B=a.\sin^2B$ (C đúng)\\
	Có: $BH.HC=a.\cos^2B.a.\sin^2B=a^2.\sin^2B.\cos^2B$ (D sai)}
\end{ex}
%%==========Câu 50
\begin{ex}%[9H4V2-4]
	Cho tam giác $ABC$ vuông tại $A$ có $AB=c$, $AC=b$, $BC=a$. Tia phân giác của góc $B$ cắt $AC$ tại $D$. Trong các khẳng định sau, khẳng định nào đúng?
	\choice
	{$\tan\dfrac{B}{2}=\dfrac{a+c}{b-c}$}
	{$\tan\dfrac{B}{2}=\dfrac{b}{a-c}$}
	{\True $\tan\dfrac{B}{2}=\dfrac{b}{a+c}$}
	{$\tan\dfrac{B}{2}=\dfrac{a}{a+c}$}
	\loigiai{
% TODO: \usepackage{graphicx} required
\begin{center}
	\includegraphics[scale=0.6]{C:/texstudio-man/Anh/{E7F84591-C1B1-4F91-97A6-2870A32AE185}}
\end{center}
	Xét tam giác $ABD$ vuông tại $A$, ta có: $\tan\dfrac{B}{2}=\tan B_1=\dfrac{AD}{AB}$ \\
	Vì $BD$ là đường phân giác của góc $B$, nên: $\dfrac{AD}{AB}=\dfrac{DC}{BC}=\dfrac{AD+DC}{AB+BC}=\dfrac{AC}{AB+BC}=\dfrac{b}{a+c}$ \\
	Vậy $\tan\dfrac{B}{2}=\dfrac{b}{a+c}$ \\
	}
\end{ex}
\Closesolutionfile{ans}
\indapan{10}{ans/ans9K1-12-TN-D5}
\begin{dang}{Câu thực tế tính chiều cao, khoảng cách}
\end{dang}
\Opensolutionfile{ans}[ans/ans9K1-12-TN-D6]
%%==========Câu 51
\begin{ex}%[9H4H2-5]
	Một con sông rộng $300m$. Một chiếc thuyền đi vuông góc với dòng nước, vì nước chảy nên chiếc thuyền phải đi $420m$ mới sang được tới bờ bên kia. Góc mà dòng nước đã đẩy chiếc thuyền lệch đi là \ldots. (kết quả làm tròn đến phút)
	\choice
	{$50^\circ$}
	{$60^\circ$}
	{\True $44^\circ 25'$}
	{$56^\circ$}
	\loigiai{
	Gọi góc mà dòng nước đã gạt chiếc thuyền đi là $\alpha$ ($\alpha>0$)\\
	Ta có $\cos\alpha=\dfrac{300}{420}=\dfrac{5}{7}$, suy ra $\alpha\approx 44^\circ 25'$}
\end{ex}
%%==========Câu 52
\begin{ex}%[9H4H2-5]
	Một tòa nhà có chiều cao là $AH$. Khi tia sáng mặt trời tạo với mặt đất một góc $65^\circ$ thì bóng của tòa nhà trên mặt đất có độ dài $BH=35m$. Chiều cao $AH$ của tòa nhà (kết quả làm tròn đến mét) là \ldots.
	\choice
	{$16m$}
	{$\text{32m}$}
	{\True $\text{75m}$}
	{$15m$}
	\loigiai{
	Xét tam giác $ABH$ vuông tại $H$, ta có: $AH=BH.\tan B=35.\tan 65^\circ\approx 75m$}
\end{ex}
%%==========Câu 53
\begin{ex}%[9H4V2-5]
	\immini[thm]
	{
	Từ trên tháp quan sát của ngọn hải đăng cao $150m$ so với mực nước biển, người ta nhìn thấy một chiếc tàu ở xa theo phương $BA$ tạo với phương nằm ngang một góc là $20^\circ$ (như hình vẽ). Khoảng cách $BH$ từ chân tháp đến tàu (kết quả làm tròn đến mét) là \ldots.
	\choice
	{$55m$}
	{$\text{51m}$}
	{$\text{141m}$}
	{\True $412m$}
	}
	{
	\includegraphics[scale=0.4]{C:/texstudio-man/Anh/{333788AF-50DC-4A35-8B64-EFA5B4CC362C}}
	}
	\loigiai{
	Xét tam giác $ABH$ vuông tại $H$, ta có: $\tan B=\dfrac{AH}{BH}$ hay $BH=\dfrac{AH}{\tan B}=\dfrac{150}{\tan 20^\circ}\approx 412m$}
\end{ex}
%%==========Câu 54
\begin{ex}%[9H4H2-5]
	Một máy bay từ mặt đất có đường bay lên $EF$ tạo với mặt đất $ED$ một góc $\widehat{DEF}=30^\circ$. Sau khi bay được quãng đường $EF=10\mathrm{~km}$ thì khoảng cách $FD$ của máy bay và mặt đất là
	\choice
	{$5\sqrt{3}\mathrm{~km}$}
	{\True $\text{5km}$}
	{$\text{10}\sqrt{3}\mathrm{~km}$}
	{$6\mathrm{~km}$}
	\loigiai{
	Xét tam giác $DEF$ vuông tại $D$, ta có: $FD=EF\sin DEF=10.\sin 30^\circ=5\left(\mathrm{~km}\right)$}
\end{ex}
%%==========Câu 55
\begin{ex}%[9H4H2-5]
	Nhà bạn An có một chiếc thang dài $6m$ Để an toàn người ta đặt thang sao cho nó tạo với mặt đất một góc \lq\lq an toàn\rq\rq là $65^\circ$ (tức là đảm bảo thang không bị đổ khi sử dụng). Khi đó khoảng cách giữa chân thang và chân tường là.
	\choice
	{$\text{3m}$}
	{$\text{3,5m}$}
	{$\text{2m}$}
	{\True $\text{2,5m}$}
	\loigiai{
	Khoảng cách từ chân thang tới tường là $AB$ \\
	Xét $\triangle ABC$ vuông tại $A$, ta có $AB=BC.\cos B=6.\cos 65^\circ\approx 2,5m$ \\
	Vậy cần đặt chân thang cách chân tường một khoảng $2,5m$}
\end{ex}
%%==========Câu 56
\begin{ex}%[9H4V2-5]
\immini[thm]
{
	Thang xếp chữ A gồm hai thang đơn tựa vào nhau. Để an toàn thì góc tạo bởi hai thang đơn phải đạt được $50^\circ$. Nếu muốn tạo một thang xếp chữ A cao 2m tính từ mặt đất thì độ dài mỗi thang đơn (kết quả làm tròn đến chữ số thập phân thứ nhất) là
	\choice
	{$\text{2m}$}
	{$\text{4,7m}$}
	{\True $\text{2,2m}$}
	{$1,8m$}
}
{
	% TODO: \usepackage{graphicx} required
	\includegraphics[scale=0.5]{C:/texstudio-man/Anh/{4D8BA695-40D7-40A2-87CF-0C42BAC1F4BF}}
}
	\loigiai{
	Do tam giác $ABC$ cân tại $A$ nên $\widehat{B}=\dfrac{180^\circ-\widehat{BAC}}{2}=\dfrac{180^\circ-50^\circ}{2}=65^\circ$ \\
	Vì tam giác $ABC$ cân tại $A$ nên đường cao $AH$ cũng là đường trung tuyến hay $H$ là trung điểm của $BC$.\\
	Xét $\triangle ABH$ vuông tại H, ta có $\sin B=\dfrac{AH}{AB}$ suy ra $AB=\dfrac{AH}{\sin B}=\dfrac{2}{\sin 65^\circ}\approx 2,2m$ \\
	Vậy thang đơn có chiều dài $2,2m$}
\end{ex}
%%==========Câu 57
\begin{ex}%[9H4V2-5]
\immini[thm]
{
	Trong hình vẽ sau biết $CH=20$, khoảng cách giữa $A$ và $B$ là \ldots.
	\choice
	{$20m$}
	{$10\sqrt{3}m$}
	{\True $20\left(\sqrt{3}-1\right)m$}
	{$20\sqrt{3}m$}
}
{
% TODO: \usepackage{graphicx} required
\includegraphics[scale=0.35]{C:/texstudio-man/Anh/{F2D56BB4-8AD9-433D-9C1D-1017BEDD035F}}

}
	\loigiai{
	Xét tam giác $AHC$ vuông tại $H$, ta có: $\widehat{ACH}=60^\circ$, nên $AH=CH.\tan 60^\circ=20\sqrt{3}m$.\\
	Xét tam giác $BCH$ vuông tại $H$, ta có: $BH=CH.\tan BCH=20.\tan 45^\circ=20m$ \\
	Suy ra $AB=AH-BH=20\sqrt{3}-20=20\left(\sqrt{3}-1\right)m$}
\end{ex}
%%==========Câu 58
\begin{ex}%[9H4V2-5]
\immini[thm]
{
	Trong hình vẽ sau, chiều cao $AD$ của cái cây là \ldots.
	\choice
	{\True $24m$}
	{$20m$}
	{$17m$}
	{$22m$}
}
{
	% TODO: \usepackage{graphicx} required
	\includegraphics[scale=0.47]{C:/texstudio-man/Anh/{3F52FE0B-F2C9-4BF0-88B1-16B6893C0CA4}}
}
	\loigiai{
	Xét tam giác $ABC$ vuông tại $C$, ta có: $AC=BC.\tan ABC$ \\
	Mà $BC=ED$ nên $AC=ED.\tan ABC=25.\tan 42^\circ\approx 22,5m$.\\
	Suy ra $AD=AC+CD\approx 22,5+1,7\approx 24m$}
\end{ex}
%%==========Câu 59
\begin{ex}%[9H4C2-5]
\immini[thm]
{
	Một người đứng trên tháp của một ngọn hải đăng cao $60m$ quan sát hai lần một con thuyền đang hướng về ngọn hải đăng. Lần thứ nhất người đó nhìn thấy thuyền với góc hạ là $20^\circ$, lần thứ 2 người đó nhìn thấy thuyền với góc hạ là $30^\circ$. Độ dài quãng đường con thuyền đã đi được giữa hai lần quan sát (làm tròn hai chữ số thập phân) là
	\choice
	{$56,48m$}
	{\True $\text{60,93m}$}
	{$15,27m$}
	{$12,80m$}
}
{
	% TODO: \usepackage{graphicx} required
	\includegraphics[scale=0.38]{C:/texstudio-man/Anh/{0B2B50B5-4A2B-4E13-A60D-4DF79A390A76}}
}
	\loigiai{+) Xét tam giác $ABC$ vuông tại $A$, ta có: $\tan ACB=\dfrac{AB}{AC}$ \\
	Hay $AC=\dfrac{AB}{\tan ACB}=\dfrac{60}{\tan 20^\circ}\approx 164,85m$.\\
	+) Xét tam giác $ABD$ vuông tại $A$, ta có: $\tan ADB=\dfrac{AB}{AD}$ \\
	Hay $AD=\dfrac{AB}{\tan ADB}=\dfrac{60}{\tan 30^\circ}=60\sqrt{3}m$ \\
	Suy ra $CD=AC-AD\approx 164,5-60\sqrt{3}\approx 60,93m$}
\end{ex}
%%==========Câu 60
\begin{ex}%[9H4C2-5]
	\immini[thm]
	{
	Trên quả đồi có một cái tháp cao $100m$. Từ đỉnh $B$ và chân $C$ của tháp nhìn điểm $A$ ở chân đồi dưới các góc tương ứng bằng $60^\circ$ và $30^\circ$ so với phương nằm ngang (như hình vẽ sau). Chiều cao $h$ của quả đồi là \ldots.
	\choice
	{$45m$}
	{\True $\text{50m}$}
	{$47m$}
	{$52m$}
	}
	{
	% TODO: \usepackage{graphicx} required
	\includegraphics[scale=0.42]{C:/texstudio-man/Anh/{23853E00-9B1B-4D45-B072-7ADD949A868A}}
	}
	\loigiai{Xét tam giác $ABD$ vuông tại $D$, ta có: $\tan ABD=\dfrac{AD}{BD}$ \\
	Vì $\widehat{ABD}=90^\circ-\widehat{ABx}=90^\circ-60^\circ=30^\circ$; $BD=BC+CD=100+h$ \\
	Nên $\tan 30^\circ=\dfrac{AD}{100+h}$ hay $\dfrac{\sqrt{3}}{3}=\dfrac{AD}{100+h}$ \\
	Suy ra: $AD=\dfrac{\sqrt{3}\left(100+h\right)}{3}$ 	(1)\\
	Xét tam giác $ACD$ vuông tại $D$, ta có: $\tan CAD=\dfrac{CD}{AD}$ \\
	Vì $\widehat{CAD}=\widehat{ACE}=30^\circ$ (2 góc so le trong); $CD=h$ \\
	Nên $\tan 30^\circ=\dfrac{h}{AD}$ hay $\dfrac{\sqrt{3}}{3}=\dfrac{h}{AD}$ \\
	Suy ra: $AD=\dfrac{3h}{\sqrt{3}}$ 	(2)\\
	Từ (1) và (2), ta có: $\dfrac{3h}{\sqrt{3}}=\dfrac{\sqrt{3}\left(100+h\right)}{3}$ \\
	$3.3h=\sqrt{3}\left(100+h\right).\sqrt{3}$ \\
	$9h-3h=3.100$ \\
	$6h=300$ \\
	$h=300:6=50$ \\
}
\end{ex}
\Closesolutionfile{ans}
\indapan{10}{ans/ans9K1-12-TN-D6}
\subsubsection{Bài tập trả lời ngắn}
\begin{dang}{Nhận biết hệ thức giữa cạnh và góc trong tam giác vuông. Áp dụng hệ thức tính độ dài cạnh trong tam giác vuông}
\end{dang}
\Opensolutionfile{ans}[ans/ans9K1-12-TLN-D1]
%%==========Câu 61
\begin{ex}%[9H4H2-1]
	Cho tam giác $ABC$ vuông tại $A$ có $AB=c$, $AC=b$, $BC=a$. Ta có hệ thức $b=...=...$
	\shortans{$b=a.\sin B=a.\cos C$}
	\loigiai{Xét tam giác $ABC$ vuông tại $A$ ta có $b=a.\sin B=a.\cos C=c.\tan B=c.\cot C$}
\end{ex}
%%==========Câu 62
\begin{ex}%[9H4H2-1]
	Cho tam giác $ABC$ vuông tại $A$ có $AB=c$, $AC=b$, $BC=a$. Ta có hệ thức $a=...=...$
	\shortans{$a=\dfrac{b}{\sin B}=\dfrac{b}{\cos C}$}
	\loigiai{Xét tam giác $ABC$ vuông tại $A$ ta có $a=\dfrac{b}{\sin B}=\dfrac{b}{\cos C}=\dfrac{c}{\sin C}=\dfrac{c}{\cos B}$}
\end{ex}
%%==========Câu 63
\begin{ex}%[9H4H2-1]
	Cho tam giác $ABC$ vuông tại $A$ có $AB=c$, $AC=b$, $BC=a$, $\widehat{C}=30^{\circ}$. Ta có hệ thức $c=...=...$
	\shortans{$c=a.\cos 30=b.\tan 30$}
	\loigiai{Xét tam giác $ABC$ vuông tại $A$ ta có $c=a.\cos 30^{\circ}=b.\tan 30^{\circ}$}
\end{ex}
%%==========Câu 64
\begin{ex}%[9H4H2-1]
	Cho tam giác $ABC$ vuông tại $A$ có $AC=b$, $BC=a$, $\widehat{C}=50^{\circ}$. Ta có hệ thức $c=...$
	\shortans{$c=a.\cos 50$}
	\loigiai{Xét tam giác $ABC$ vuông tại $A$ ta có $c=a.\cos 50^{\circ}$}
\end{ex}
%%==========Câu 65
\begin{ex}%[9H4H2-1]
	Cho tam giác $MNP$ vuông tại $M$ có $NP=8\,\text{cm}$ $\widehat{P}=30^{\circ}$. Số đo cạnh $MN$ là
	\shortans{$4$}
	\loigiai{Xét tam giác $MNP$ vuông tại $M$ có $MN=NP.\sin P=8.\sin 30^{\circ}=4\,\text{cm}.$}
\end{ex}
%%==========Câu 66
\begin{ex}%[9H4H2-1]
	Cho tam giác $MNP$ vuông tại $M$ có $NP=8\,\text{cm}$, $\widehat{P}=30^{\circ}$. Số đo cạnh $MP$ là .....
	\shortans{$6,93$}
	\loigiai{Xét tam giác $MNP$ vuông tại $M$ có $MP=NP.\cos P=8.\cos 30^{\circ}=4\sqrt{3}\,\text{cm}.$}
\end{ex}
%%==========Câu 67
\begin{ex}%[9H4H2-1]
	Cho tam giác $DEF$ vuông tại $D$ có $DE=6\,\text{cm}$, $\widehat{E}=50^{\circ}$. Số đo cạnh $DE$ (làm tròn tới chữ số thập phân thứ nhất là
	\shortans{$7,2$}
	\loigiai{Xét tam giác $DEF$ vuông tại $D$ có $DE=DF.\tan E=6.\tan 50^{\circ}\approx7,2\,\text{cm}.$}
\end{ex}
%%==========Câu 68
\begin{ex}%[9H4H2-1]
	Cho tam giác $DEF$ vuông tại $D$ có $DE=6\,\text{cm}$, $\widehat{E}=50^{\circ}$. Số đo cạnh $EF$ (làm tròn tới chữ số thập phân thứ nhất là
	\shortans{$7,8$}
	\loigiai{Xét tam giác $DEF$ vuông tại $D$ có $DF=EF.\sin E$ suy ra $EF=\dfrac{DF}{\sin E}=\dfrac{6}{\sin 50^{\circ}}\approx7,8$.}
\end{ex}
\Closesolutionfile{ans}
\indapan{4}{ans/ans9K1-12-TLN-D1}
%%==========Câu 9
\begin{dang}{Tính số đo góc khi biết độ dài các cạnh trong tam giác vuông}
\end{dang}
\Opensolutionfile{ans}[ans/ans9K1-12-TLN-D2]
%%==========Câu 69
\begin{ex}%[9H4H2-1]
	Cho $\triangle ABC$ vuông tại $A$. Biết $AB=3\,\text{cm}$, $AC=7\,\text{cm}$. Số đo $\widehat{B}$ làm tròn đến độ là
	\shortans{$67$}
	\loigiai{Xét $\triangle ABC$ vuông tại $A$. Ta có $\tan B=\dfrac{AC}{AB}=\dfrac{7}{3}$ suy ra $\widehat{B}\approx67^{\circ}$.}
\end{ex}
%%==========Câu 70
\begin{ex}%[9H4H2-1]
	Cho $\triangle ABC$ vuông tại $A$. Biết $AB=2\,\text{cm}$, $BC=5\,\text{cm}$. Số đo $\widehat{B}$ làm tròn đến độ là
	\shortans{$66$}
	\loigiai{Xét $\triangle ABC$ vuông tại $A$. Ta có $\cos B=\dfrac{AB}{BC}=\dfrac{2}{5}$ suy ra $\widehat{B}\approx66^{\circ}$.}
\end{ex}
%%==========Câu 71
\begin{ex}%[9H4H2-1]
	Cho $\triangle MNP$ vuông tại $M$. Biết $MP=9\,\text{cm}$, $NP=15\,\text{cm}$. Số đo $\widehat{P}$ làm tròn đến phút là......
	\shortans{$53,13$}
	\loigiai{Xét $\triangle MNP$ vuông tại $M$. Ta có $\cos P=\dfrac{MP}{NP}=\dfrac{9}{15}$ suy ra $\widehat{P}\approx53^{\circ}8'$.}
\end{ex}
%%==========Câu 72
\begin{ex}%[9H4H2-1]
	Cho $\triangle DEF$ vuông tại $D$. Biết $DE=15\,\text{cm}$, $EF=2\,\text{dm}$. Số đo $\widehat{E}$ làm tròn đến phút là.....
	\shortans{$41,4$}
	\loigiai{Đổi $2\,\text{dm}=20\,\text{cm}$.\\ Xét $\triangle DEF$ vuông tại $D$. Ta có $\cos E=\dfrac{DE}{EF}=\dfrac{15}{20}=\dfrac{3}{4}$ suy ra $\widehat{E}\approx41^{\circ}24'$.}
\end{ex}
%%==========Câu 73
\begin{ex}%[9H4H2-1]
	Cho $\triangle ABC$ vuông tại $A$, đường cao $AH$. Biết $BH=2\,\text{cm}$, $CH=8\,\text{cm}$. Số đo $\widehat{B}$ làm tròn đến độ là
	\shortans{$63$}
	\loigiai{Xét $\triangle ABC$ vuông tại $A$, đường cao $AH$ ta có $AH^{2}=BH.HC=2.8=16$ suy ra $AH=4\,\text{cm}$.\\ Xét $\triangle ABH$ vuông tại $H$ ta có $\tan B=\dfrac{AH}{BH}=\dfrac{4}{2}=2.$ Suy ra $\widehat{B}\approx63^{\circ}$.}
\end{ex}
%%==========Câu 74
\begin{ex}%[9H4H2-1]
	Cho $\triangle ABC$ có $AB=20\,\text{cm}$, $BC=21\,\text{cm}$, $AC=29\,\text{cm}$. Số đo $\widehat{C}$ làm tròn đến độ là....
	\shortans{$44$}
	\loigiai{Ta có $BC^{2}=29^{2}=841$; $AB^{2}+AC^{2}=20^{2}+21^{2}=400+441=841$. Suy ra $BC^{2}=AB^{2}+AC^{2}$ nên $\triangle ABC$ vuông tại $A$. Xét $\triangle ABC$ vuông tại $A$, ta có $\sin~C=\dfrac{AB}{AC}=\dfrac{20}{29}.$ Suy ra $\widehat{C}\approx44^{\circ}$.}
\end{ex}
%%==========Câu 75
\begin{ex}%[9H4H2-1]
	Cho $\triangle ABC$ đường cao $AH$. Biết $BH=2\,\text{cm}$, $CH=4\,\text{cm}$, $\widehat{B}=45^{\circ}$. Số đo $\widehat{C}$ làm tròn đến phút là.....
	\shortans{$26,57$}
	\loigiai{Xét $\triangle ABH$ vuông tại $H$, ta có $AH=BH.\tan B=2.\tan 45^{\circ}=2\,\text{cm}.$ \\ Xét $\triangle ACH$ vuông tại $H$, ta có $\tan C=\dfrac{AH}{HC}=\dfrac{2}{4}=\dfrac{1}{2}.$ Suy ra $\widehat{C}\approx26^{\circ}34'$.}
\end{ex}
%%==========Câu 76
\begin{ex}%[9H4H2-1]
	Cho tam giác $ABC$ có $AB=10\,\text{cm}$, $AC=12\,\text{cm}$, $\widehat{A}=40^{\circ}$. Số đo $\widehat{C}$ (làm tròn đến độ) là.....
	\shortans{$56$}
	\loigiai{Kẻ đường cao $BH$ xuống $AC$ ($H$ thuộc $AC$). \\ Xét tam giác $ABH$ vuông tại $H$, ta có \\ $BH=AB.\sin\widehat{A}=10.\sin 40^{\circ}\approx6,428\,\text{cm}$ \\ $AH=AB.\cos\widehat{A}=10.\cos 40^{\circ}\approx7,66\,\text{cm}$ \\ $CH=AC-AH=12-7,66\approx4,34\,\text{cm}$ \\ Xét tam giác $BHC$ vuông tại $H$, ta có: $\tan C=\dfrac{BH}{CH}\approx\dfrac{6,428}{4,34}\approx1,481$ \\ suy ra $\widehat{C}\approx56^{\circ}$.}
\end{ex}
\Closesolutionfile{ans}
\indapan{6}{ans/ans9K1-12-TLN-D2}
\begin{dang}{Tính độ dài cạnh, số đo góc trong tam giác không vuông}
\end{dang}
\Opensolutionfile{ans}[ans/ans9K1-12-TLN-D3]
%%==========Câu 77
\begin{ex}%[9H4H2-1]
	Cho $\triangle ABC$ vuông tại $A$. Biết $AB=6\,\text{cm}$, $AC=8\,\text{cm}$. Số đo $\widehat{B}$ (làm tròn đến độ) là .....
	\shortans{$53$}
	\loigiai{Xét $\triangle ABC$ vuông tại $A$ ta có $BC=\sqrt{AB^{2}+AC^{2}}=\sqrt{6^{2}+8^{2}}=\sqrt{36+64}=\sqrt{100}=10\,\text{cm}.$ \\ Ta có $\sin B=\dfrac{AC}{BC}=\dfrac{8}{10}=\dfrac{4}{5}.$ Suy ra $\widehat{B}\approx53^{\circ}$.}
\end{ex}
%%==========Câu 78
\begin{ex}%[9H4H2-1]
	Cho tam giác $MNP$ vuông tại $M$ có $NP=8\,\text{cm}$, $\widehat{P}=30^{\circ}$. Chu vi tam giác $MNP$ (làm tròn tới hàng phần chục) là
	\shortans{$18,9$}
	\loigiai{Xét tam giác $MNP$ vuông tại $M$ có: \\ $MN=NP.\sin P=8.\sin 30^{\circ}=8.\dfrac{1}{2}=4\,\text{cm}.$ \\ $MP=NP.\cos P=8.\cos 30^{\circ}=8.\dfrac{\sqrt{3}}{2}=4\sqrt{3}\,\text{cm}.$ \\ Chu vi tam giác $MNP$ bằng $MN+MP+NP = 4+4\sqrt{3}+8 = 12+4\sqrt{3} \approx 12+4 \times 1,732 = 12+6,928 = 18,928 \approx 18,9\,\text{cm}.$}
\end{ex}
%%==========Câu 79
\begin{ex}%[9H4H2-1]
	Cho tam giác $IHK$ vuông tại $I$ có $IH=5\,\text{cm}$, $\widehat{K}=42^{\circ}$. Diện tích tam giác $IHK$ (làm tròn tới hàng phần chục) là.....
	\shortans{$9$}
	\loigiai{Xét tam giác $IHK$ vuông tại $I$ có $IK=IH.\tan K=5.\tan 42^{\circ}\approx4,5\,\text{cm}.$ \\ Diện tích tam giác $IHK$ bằng: $\dfrac{1}{2}IH.IK\approx\dfrac{1}{2}.4.4,5\approx9\,\text{cm}^{2}.$}
\end{ex}
%%==========Câu 80
\begin{ex}%[9H4H2-3]
	Cho $\widehat{xOy}=45^{\circ}$. Trên tia $Oy$ lấy 2 điểm $A,B$ sao cho $AB=\sqrt{2}\,\text{dm}$. Độ dài hình chiếu vuông góc của đoạn thẳng $AB$ trên $Ox$ là:
	\shortans{$1$}
	\loigiai{Từ $A$ kẻ $AH\perp BD$ tại $H$. Suy ra $AH//OD$ (vì cùng vuông góc với $BD$). Suy ra $\widehat{BAH}=\widehat{O}=45^{\circ}$. Ta có: $AH=AB.\sin\widehat{BAD}=\sqrt{2}.\sin 45^{\circ}=1\,\text{dm}.$ Từ giác $ACDH$ là hình chữ nhật nên $CD=AH=1\,\text{dm}$.}
\end{ex}
%%==========Câu 81
\begin{ex}%[9H4H2-3]
	Cho $\triangle ABC$, đường cao $AH$. Biết $BH=3\,\text{cm}$, $\widehat{B}=45^{\circ}$, $\widehat{C}=30^{\circ}$. Chu vi $\triangle ABC$ (làm tròn tới hàng phần chục) là.....
	\shortans{$18,4$}
	\loigiai{Xét tam giác $AHB$ vuông tại $H$ ta có: \\ $AH=BH.\tan B=3.\tan 45^{\circ}=3.1=3\,\text{cm}.$ \\ $AB=\dfrac{AH}{\sin B}=\dfrac{3}{\sin 45^{\circ}}=\dfrac{3}{\frac{\sqrt{2}}{2}}=\dfrac{6}{\sqrt{2}}=3\sqrt{2}\,\text{cm}.$ \\ Xét tam giác $AHC$ vuông tại $H$ ta có: \\ $HC=\dfrac{AH}{\tan C}=\dfrac{3}{\tan 30^{\circ}}=\dfrac{3}{\frac{1}{\sqrt{3}}}=3\sqrt{3}\,\text{cm}.$ \\ $AC=\dfrac{AH}{\sin C}=\dfrac{3}{\sin 30^{\circ}}=\dfrac{3}{\frac{1}{2}}=6\,\text{cm}.$ \\ $BC=BH+HC=3+3\sqrt{3}\,\text{cm}.$ \\ Chu vi tam giác $ABC$ bằng $AB+AC+BC = 3\sqrt{2}+6+3+3\sqrt{3} = 9+3\sqrt{2}+3\sqrt{3} \approx 9+3(1,414)+3(1,732) = 9+4,242+5,196 = 18,438 \approx 18,4\,\text{cm}.$}
\end{ex}
%%==========Câu 82
\begin{ex}%[9H4V2-3]
	Cho $\triangle ABC$, đường cao $AH$. Biết $AB=8\,\text{cm}$, $\widehat{B}=42^{\circ}$, $\widehat{C}=37^{\circ}$. Diện tích $\triangle ABC$ (làm tròn tới hàng phần trăm) là.....
	\shortans{$34,93$}
	\loigiai{Xét tam giác $AHB$ vuông tại $H$ ta có: \\ $AH=AB.\sin B=8.\sin 42^{\circ}$ \\ $BH=AB.\cos B=8.\cos 42^{\circ}$ \\ Xét tam giác $AHC$ vuông tại $H$ ta có: \\ $HC=\dfrac{AH}{\tan C}=\dfrac{8.\sin 42^{\circ}}{\tan 37^{\circ}} = 8.\sin 42^{\circ}.\cot 37^{\circ}.$ \\ $BC=BH+HC=8.\cos 42^{\circ}+8.\sin 42^{\circ}.\cot 37^{\circ}.$ \\ Diện tích tam giác $ABC$ bằng $\dfrac{1}{2}AH.BC=\dfrac{1}{2}.8.\sin 42^{\circ}.(8.\cos 42^{\circ}+8.\sin 42^{\circ}.\cot 37^{\circ})\approx34,93\,\text{cm}^{2}$.}
\end{ex}
%%==========Câu 83
\begin{ex}%[9H4V2-3]
	Cho $\triangle ACD$ vuông tại $C$ có $\widehat{DAC}=30^\circ$. Lấy điểm $B$ nằm trên cạnh $AC$ sao cho $\triangle BDC$ vuông cân tại $C$, khoảng cách $AB$ là ....(làm tròn kết quả đến chữ số thập phân thứ nhất)
	\shortans{$14,6$}
	\loigiai{Xét $\triangle ADC$ vuông tại $C$ ta có $AC=DC.\cot (\text{góc CAD})=20.\cot 30^{\circ}=20.\sqrt{3}\,\text{cm}.$ \\ Xét $\triangle BDC$ vuông tại $C$ ta có $BC=DC.\tan (\text{góc CDB})=20.\tan 45^{\circ}=20.1=20\,\text{cm}.$ \\ $AB=AC-BC=20\sqrt{3}-20 \approx 20(1,73205) - 20 = 34,641 - 20 = 14,641 \approx 14,6\,\text{cm}.$}
\end{ex}
%%==========Câu 84
\begin{ex}%[9H4V2-3]
	Cho $\triangle ABC$. Biết $BC=12\,\text{cm}$, $\widehat{B}=60^{\circ}$, $\widehat{C}=40^{\circ}$. Diện tích tam giác $ABC$ là .....
	\shortans{$40,7$}
	\loigiai{Kẻ đường cao $AH$ từ $A$ xuống $BC$. \\ Xét $\triangle ABH$ vuông tại $H$ ta có $BH=AH.\cot B=AH.\cot 60^{\circ}$. \\ Xét $\triangle ACH$ vuông tại $H$ ta có $CH=AH.\cot C=AH.\cot 40^{\circ}$. \\ Suy ra $BC=BH+CH=AH.\cot B+AH.\cot C=AH.(\cot 60^{\circ}+\cot 40^{\circ})$. \\ Suy ra $AH=\dfrac{BC}{\cot 60^{\circ}+\cot 40^{\circ}}=\dfrac{12}{\cot 60^{\circ}+\cot 40^{\circ}}$. \\ $S_{ABC}=\dfrac{1}{2}.BC.AH=\dfrac{1}{2}.12.\dfrac{12}{\cot 60^{\circ}+\cot 40^{\circ}} = \dfrac{72}{\cot 60^{\circ}+\cot 40^{\circ}} \approx \dfrac{72}{0,57735 + 1,19175} = \dfrac{72}{1,7691} \approx 40,70\,\text{cm}^{2}.$}
\end{ex}
\Closesolutionfile{ans}
\indapan{6}{ans/ans9K1-12-TLN-D3}
\begin{dang}{Tính độ dài cạnh, số đo góc trong tam giác không vuông}
\end{dang}
\Opensolutionfile{ans}[ans/ans9K1-12-TLN-D4]
%%==========Câu 85
\begin{ex}%[9H4H2-3]
	Cho tam giác $ABC$, $CH\perp AB$ có $AB=4,5 \text{ cm}$, $AC=3 \text{ cm,}$ $\widehat{A}=40^{\circ}$. Độ dài đoạn thẳng $CH$ (kết quả làm tròn đến chữ số thập phân thứ hai) là
	\shortans{$1,93$}
	\loigiai{Xét tam giác $HAC$ vuông tại $H$, ta có: $CH=AC.\sin\widehat{HAC}=3.\sin 40^{\circ}\approx1,93 \text{ m}$}
\end{ex}
%%==========Câu 86
\begin{ex}%[9H4H2-3]
	Cho tam giác $ABC$, $CH\perp AB$ có $AB=4,5 \text{ cm}$, $AC=3 \text{ cm,}$ $\widehat{A}=40^{\circ}$.	Diện tích tam giác $ABC$ (kết quả làm tròn đến chữ số thập phân thứ hai) là
	\shortans{$4,34$}
	\loigiai{Xét tam giác $HAC$ vuông tại $H$, ta có: $CH=AC \sin\widehat{HAC}=3 \sin 40^{\circ}\approx1,93 \text{ m}.$ \\
	Diện tích của tam giác $ABC$ là: $S=\dfrac{1}{2}AB.CH\approx\dfrac{1}{2}.4,5.1,93\approx4,34 \text{ cm}^{2}.$}
\end{ex}
%%==========Câu 87
\begin{ex}%[9H4H2-3]
	Cho tam giác $ABC$ có $\widehat{B}=78^{\circ},$ đường cao $AH=12 \text{ cm}$. Độ dài của cạnh $AB$ của tam giác $ABC$ (kết quả làm tròn đến hàng phần mười) là
	\shortans{$12,3$}
	\loigiai{Xét tam giác $HAB$ vuông tại $H$, ta có: \\
	$\sin B=\dfrac{AH}{AB}$ \\
	hay $AB=\dfrac{AH}{\sin B}=\dfrac{12}{\sin 78^{\circ}}\approx12,3 \text{ cm}.$}
\end{ex}
%%==========Câu 88
\begin{ex}%[9H4H2-3]
	Cho tam giác $DEF$ có $DE=8 \text{ cm}$, $\widehat{D}=40^{\circ}$, $\widehat{F}=58^{\circ}$. Kẻ đường cao $EI$ của tam giác đó. Độ dài của đường cao $EI$ (kết quả làm tròn đến chữ số thập phân thứ hai) bằng
	\shortans{$5,14$}
	\loigiai{Xét tam giác $DEI$ vuông tại $I$, ta có: $EI=ED.\sin D=8.\sin 40^{\circ}\approx5,14 \text{ cm}.$}
\end{ex}
%%==========Câu 89
\begin{ex}%[9H4H2-3]
	Cho tam giác $ABC$ có $AB=14 \text{ cm}$, $AC=13 \text{ cm}$ và $\widehat{B}=60^{\circ}$. Độ dài của $BC$ (kết quả làm tròn đến chữ số thập phân thứ nhất) là
	\shortans{$11,7$}
	\loigiai{Kẻ đường cao $AH$. \\
	+) Xét tam giác $ABH$ vuông tại $H$, ta có: $BH=AB. \cos B=AB. \cos 60^{\circ}=14 \cdot \dfrac{1}{2}=7 \text{ cm}.$ \\
	$AH=AB.\sin B=AB.\sin 60^{\circ}=14\cdot\dfrac{\sqrt{3}}{2}=7\sqrt{3} \text{ cm}.$ \\
	+) Xét tam giác $ACH$ vuông tại $H$, ta có: $HC^{2}=AC^{2}-AH^{2}=13^{2}-(7\sqrt{3})^{2}=22$ \\
	Suy ra $HC=\sqrt{22} \text{ cm}$ \\
	Vậy $BC=BH+HC=7+\sqrt{22}\approx11,7 \text{ cm}$}
\end{ex}
%%==========Câu 90
\begin{ex}%[9H4V2-3]
	\immini[thm]
	{
		Cho hình vẽ sau, độ dài đoạn thẳng $BC$ (kết quả làm tròn đến chữ số thập phân thứ hai) là:
	\shortans{$6,73$}
	}
	{
		% TODO: \usepackage{graphicx} required
		\includegraphics[scale=0.5]{C:/texstudio-man/Anh/{03A2DA85-1294-4146-8DFC-0FE2F38EAF62}}
	}
	\loigiai{
	+) Xét tam giác $ABH$ vuông tại $H$, ta có: $HB=AH.\tan \widehat{BAH}=4.\tan 28^{\circ}\approx2,13 \text{ cm}.$ \\
	+) Xét tam giác $AHC$ vuông tại $H$, ta có: $HC=AH.\cot C=4.\cot 41^{\circ}\approx4,60 \text{ cm}.$ \\
	Có: $BC=HB+HC\approx2,13+4,60\approx6,73 \text{ cm}$}
\end{ex}
%%==========Câu 91
\begin{ex}%[9H4V2-3]
	Cho tam giác $ABC$ có $AB=12$, $AC=15$ và $\widehat{B}=60^{\circ}$ Độ dài cạnh $BC$ (kết quả làm tròn đến chữ số thập phân thứ nhất) là:
	\shortans{$16,8$}
	\loigiai{Kẻ đường cao $AH$, ta có hình vẽ sau: \\
	+) Xét tam giác $ABH$ vuông tại $H$, ta có: $BH=AB.\cos B=12.\cos 60^{\circ}=6$ \\
	$AH=AB.\sin B=12.\sin 60^{\circ}=12\cdot\dfrac{\sqrt{3}}{2}=6\sqrt{3}$ \\
	+) Xét tam giác $ACH$ vuông tại $H$, ta có: $HC^{2}=AC^{2}-AH^{2}=15^{2}-(6\sqrt{3})^{2}=117$ \\
	Suy ra $HC=3\sqrt{13}$ \\
	Vậy $BC=BH+HC=6+3\sqrt{13}\approx16,8$}
\end{ex}
%%==========Câu 92
\begin{ex}%[9H4V2-3]
	Cho tam giác $ABC$ có $\widehat{B}=70^{\circ}$, $\widehat{C}=35^{\circ}$, $AC=4,5 \text{ cm}$. Diện tích tam giác $ABC$ (làm tròn đến chữ số thập phân thứ hai) là:
	\shortans{$5,97$}
	\loigiai{Kẻ đường cao $AD$, ta có hình vẽ sau: \\
	+) Xét tam giác $ACD$ vuông tại $D$, ta có: $AD=AC.\sin C=4,5.\sin 35^{\circ}\approx2,58 \text{ cm}$ \\
	$CD=AC. \cos C=4,5. \cos 35^{\circ}\approx3,69 \text{ cm}$ \\
	+) Xét tam giác $ABD$ vuông tại $D$, ta có: $\tan B=\dfrac{AD}{BD}$ \\
	hay $BD=\dfrac{AD}{\tan B}\approx\dfrac{2,58}{\tan 70^{\circ}}\approx0,94 \text{ cm}$ \\
	Suy ra $BC=BD+DC\approx0,94+3,69\approx4,63 \text{ cm}$ \\
	Do đó: $S_{ABC}=\dfrac{1}{2} BC. AD\approx\dfrac{1}{2}.4,63.2,58\approx5,97 \text{ cm}^{2}$}
\end{ex}
\Closesolutionfile{ans}
\indapan{6}{ans/ans9K1-12-TLN-D4}
\begin{dang}{Chứng tỏ được hệ thức liên hệ giữa các cạnh}
\end{dang}
\Opensolutionfile{ans}[ans/ans9K1-12-TLN-D5]
%%==========Câu 93
\begin{ex}%[9H4H2-3]
Cho tam giác $ABC$ vuông tại $A$, có $\widehat{C}=30^{\circ}$ và $BC=6 \text{ cm}$, ta có độ dài đoạn thẳng $AB$ là:
\shortans{$3$}
\loigiai{Xét $\Delta ABC$ vuông tại $A$, ta có: $AB=BC. \sin C=6.\sin 30^{\circ}=6\cdot\dfrac{1}{2}=3 \text{ cm}$}
\end{ex}
%%==========Câu 94
\begin{ex}%[9H4H2-3]
Cho tam giác $ABC$ vuông tại $A$, có $\widehat{C}=30^{\circ}$ và $BC=6 \text{ cm}$, ta có độ dài đoạn thẳng $AC$ (kết quả làm tròn đến hàng phần mười) là
\shortans{$5,2$}
\loigiai{Xét $\Delta ABC$ vuông tại $A$, ta có: $AC=BC.\cos C=6.\cos 30^{\circ}=6\cdot\dfrac{\sqrt{3}}{2}\approx5,2 \text{ cm.}$}
\end{ex}
%%==========Câu 95
\begin{ex}%[9H4H2-1]
	Cho tam giác $DEF$ có $DE=7 \text{ cm,}$ $\widehat{D}=40^{\circ}$, $\widehat{F}=58^{\circ}$. Kẻ đường cao $EI$ của tam giác $DEF$. Độ dài cạnh $EF$ (kết quả làm tròn đến hàng phần trăm) là
	\shortans{$5,31$}
	\loigiai{
	+) Xét tam giác $DEI$ vuông tại $I$ , ta có: $EI=ED. \sin D=7. \sin 40^{\circ}\approx4,50 \text{ cm}$ \\
	+) Xét tam giác $EIF$ vuông tại $I$ , ta có: $EI=EF. \sin F$ hay $EF=\dfrac{EI}{\sin F}\approx\dfrac{4,50}{\sin 58^{\circ}}\approx5,31 \text{ cm}$}
\end{ex}
%%==========Câu 96
\begin{ex}%[9H4V2-1]
	\immini[thm]
	{
		Cho hình vẽ, độ dài đoạn thẳng $AB$ (kết quả làm tròn đến hàng phần trăm) là
	\shortans{$5,93$}
	}
	{
		% TODO: \usepackage{graphicx} required
		\includegraphics[scale=0.6]{C:/texstudio-man/Anh/{81FFE273-922A-42DA-874B-626E134DF544}}
		
	}
	\loigiai{
	+) Xét tam giác $BCH$ vuông tại $H$, ta có: $BH=BC.\sin 30^{\circ}=11\cdot\dfrac{1}{2}=5,5 \text{ cm}$ \\
	Có: $\widehat{C}=30^{\circ}$ nên $\widehat{HBC}=60^{\circ},$ suy ra $\widehat{ABH}=22^{\circ}$ \\
	+) Xét tam giác $AHB$ vuông tại $H$, ta có: \\
	$\cos \widehat{HBA}=\dfrac{BH}{AB}$ \\
	hay $AB=\dfrac{BH}{\cos \widehat{ABH}}=\dfrac{5,5}{\cos 22^{\circ}}\approx5,93 \text{ cm}$}
\end{ex}
%%==========Câu 97
\begin{ex}%[9H4H2-3]
Cho tam giác $ABC$ vuông cân tại $A(AB=AC=3cm)$. Phân giác của góc $B$ cắt $AC$ tại $D$. Độ dài đoạn thẳng $DC$ (kết quả làm tròn đến chữ số thập phân thứ hai) là
\shortans{$1,76$}
\loigiai{
% TODO: \usepackage{graphicx} required
\begin{center}
\includegraphics[scale=0.6]{C:/texstudio-man/Anh/{123FDED6-62F4-4EF0-B326-10F141A67C02}}
\end{center}
Xét $\Delta ABC$ vuông cân tại $A$, ta có: $\widehat{B}=\widehat{C}=45^{\circ}$ \\
Vì $BD$ là tia phân giác của góc $B$, nên: $\widehat{ABD}=\widehat{DBC}=\dfrac{1}{2}\widehat{ABC}=\dfrac{45^{\circ}}{2}=22,5^{\circ}$ \\
Xét $\Delta ABD$ vuông tại $A$, ta có: $AD=AB.\tan \widehat{ABD}=3.\tan 22,5^{\circ}$ \\
Ta có: $AD+DC=AC$ suy ra: $DC=AC-AD=3-3.\tan 22,5^{\circ}\approx1,76 \text{ cm}$}
\end{ex}
%%==========Câu 98
\begin{ex}%[9H4V2-3]
	Cho tam giác $ABC$ có $\widehat{B}=60^{\circ},$ $\widehat{C}=50^{\circ}$, $AC=4 \text{ cm}$. Diện tích tam giác $ABC$ (làm tròn đến hàng phần trăm) là
	\shortans{$6,64$}
	\loigiai{
		% TODO: \usepackage{graphicx} required
		\begin{center}
			\includegraphics[scale=0.6]{C:/texstudio-man/Anh/{44D44F3B-50D6-4C35-8ADF-ECC454687066}}
		\end{center}
	Kẻ đường cao $AD$ \\
	+) Xét tam giác $ACD$ vuông tại $D$, ta có: $AD=AC. \sin C=4.\sin 50^{\circ}\approx3,06 \text{ cm}$ \\
	$CD=AC.\cos C=4.\cos 50^{\circ}\approx2,57 \text{ cm.}$ \\
	+) Xét tam giác $ABD$ vuông tại $D$, ta có: \\
	$\tan B=\dfrac{AD}{BD}$ hay $BD=\dfrac{AD}{\tan B}\approx\dfrac{3,06}{\tan 60^{\circ}}\approx1,77 \text{ cm.}$ \\
	Suy ra $BC=BD+DC\approx1,77+2,57\approx4,34 \text{ cm}$ \\
	Vậy $S_{ABC}=\dfrac{1}{2}\cdot AD\cdot BC\approx\dfrac{1}{2}\cdot3,06\cdot4,34\approx6,64 \text{ cm}^{2}.$}
\end{ex}
%%==========Câu 99
\begin{ex}%[9H4V2-3]
	Cho tam giác $ABC$ có $BC=11 \text{ cm.}$, $\widehat{B}=40^{\circ}$, $\widehat{C}=30^{\circ}$. Gọi $N$ là chân đường vuông góc hạ từ $A$ xuống cạnh $BC$. Độ dài cạnh $AN$ (kết quả làm tròn đến chữ số thập phân thứ hai) là
	\shortans{$3,76$}
	\loigiai{
	% TODO: \usepackage{graphicx} required
	\begin{center}
		\includegraphics[scale=0.6]{C:/texstudio-man/Anh/{AF73AD74-2C38-4400-88DC-DE7B0BB21583}}
	\end{center}
	Đặt $BN=x(0<x<11)$ có $NC=11-x$ \\
	+) Xét tam giác $ABN$ vuông tại $N$, ta có: $AN=BN.\tan B=x.\tan 40^{\circ}$ \\
	+) Xét tam giác $ACN$ vuông tại $N$, ta có: $AN=CN.\tan C=(11-x).\tan 30^{\circ}$ \\
	Nên $x.\tan 40^{\circ}=(11-x).\tan 30^{\circ}$ \\
	$x .\tan 40^{\circ}=11.\tan 30^{\circ}-x.\tan 30^{\circ}$ \\
	$x .\tan 40^{\circ}+x.\tan 30^{\circ}=11.\tan 30^{\circ}$ \\
	$x.(\tan 40^{\circ}+\tan 30^{\circ})=11.\tan 30^{\circ}$ \\
	$x=\dfrac{11.\tan 30^{\circ}}{\tan 40^{\circ}+\tan 30^{\circ}}$ \\
	$x\approx4,48$ \\
	Khi đó, $AN=BN. \tan B\approx4,48. \tan 40^{\circ}\approx3,76 \text{ cm}$}
\end{ex}
%%==========Câu 100
\begin{ex}%[9H4V2-3]
	Cho tứ giác $ABCD$ có $\widehat{A}=\widehat{D}=90^{\circ}$, $\widehat{C}=40^{\circ}$, $AB=4 \text{ cm}$, $AD=3 \text{ cm}$. Diện tích tứ giác $ABCD$ (kết quả làm tròn đến chữ số thập phân thứ nhất) là
	\shortans{$17,3$}
	\loigiai{
		% TODO: \usepackage{graphicx} required
		\begin{center}
			\includegraphics[scale=0.6]{C:/texstudio-man/Anh/{8491DF10-023F-4DED-88DA-AB889CF2849B}}
		\end{center}
	Kẻ đường cao $BE$. \\
	+) Xét tứ giác $ABED$ có $\widehat{A}=\widehat{D}=\widehat{E}=90^{\circ}$ nên $ABED$ là hình chữ nhật. \\
	Suy ra $DE=AB=4 \text{ cm.}$ $BE=AD=3 \text{ cm}$ \\
	+) Xét tam giác $BEC$ vuông tại $E$, ta có: $EC=BE.\cot C=3.\cot 40^{\circ}$ \\
	Suy ra: $DC=DE+EC=4+3.\cot 40^{\circ}$ \\
	Do đó: $S_{ABCD}=\dfrac{(AB+CD).AD}{2}=\dfrac{(4+4+3.\cot 40^{\circ}).3}{2}\approx17,3 \text{ cm}^{2}.$}
\end{ex}
\Closesolutionfile{ans}
\indapan{6}{ans/ans9K1-12-TLN-D5}
\begin{dang}{Bài toán thực tế tính chiều cao, khoảng cách}
\end{dang}
\Opensolutionfile{ans}[ans/ans9K1-12-TLN-D6]
%%==========Câu 101
\begin{ex}%[9H4H2-5]
	Các tia nắng mặt trời tạo với mặt đất một góc bằng $34^{\circ}$ và bóng của một tháp trên mặt đất dài $86 \text{ m}$. Chiều cao của tháp (kết quả làm tròn đến mét) là
	\shortans{$58$}
	\loigiai{
	+) Xét tam giác $ABH$ vuông tại $H$, ta có: $BH=AH.\tan A=86.\tan 34^{\circ}\approx58 \text{ m}$}
\end{ex}
%%==========Câu 102
\begin{ex}%[9H4H2-5]
	\immini[thm]
	{
	Tam giác $ABC$ vuông tại $A$ ở hình sau mô tả cột cờ $AB$ và bóng nắng của cột cờ trên mặt đất là $AC$. Người ta đo được độ dài $AC=13 \text{ m}$ và $\widehat{C}=35^{\circ}$. Chiều cao $AB$ của cột cờ (làm tròn kết quả đến phần trăm của mét) là
	\shortans{$9,10$}
	}
	{
		% TODO: \usepackage{graphicx} required
		\includegraphics[scale=0.6]{C:/texstudio-man/Anh/{AA0AA41B-5694-4D9A-AAB7-E4F17125C907}}
	}
	\loigiai{
	+) Xét tam giác $ABC$ vuông tại $A$, ta có: $AB=AC. \tan C=13. \tan 35^{\circ}\approx9,10 \text{ m}.$}
\end{ex}
%%==========Câu 103
\begin{ex}%[9H4H2-5]
	\immini[thm]
	{
		Hình vẽ sau mô tả tia nắng mặt trời dọc theo $AB$ tạo với phương nằm ngang trên mặt đất một góc $\alpha=\widehat{ABH}$. Biết $AH=2 \text{ m}$, $BH=5 \text{ m}$, số đo góc $\alpha$ (kết quả làm tròn đến độ), là
	\shortans{$22$}
	}
	{
			\includegraphics[scale=0.6]{C:/texstudio-man/Anh/{1B1CC3EA-DC56-4FB3-9F88-232323232323244436444A1C8419}}
		
	}
	\loigiai{
	+) Xét tam giác $ABH$ vuông tại $H$, ta có: $\tan B=\dfrac{AH}{BH}=\dfrac{2}{5}$ \\
	Suy ra $\widehat{B}\approx22^{\circ}$ hay $\alpha\approx22^{\circ}$}
\end{ex}
%%==========Câu 104
\begin{ex}%[9H4H2-5]
	\immini[thm]
	{
		Một cột đèn cao $7 \text{ m}$ có bóng trên mặt đất dài $4 \text{ m}$. Góc (làm tròn đến độ) mà tia sáng Mặt Trời tạo với mặt đất góc $\alpha$ trong hình sau là
	\shortans{$60$}
	}
	{
		% TODO: \usepackage{graphicx} required
		\includegraphics[scale=0.6]{C:/texstudio-man/Anh/{1B1CC3EA-DC56-4FB3-9F88-232323232323244436444A1C8419}1}
	}
	\loigiai{
	+) Xét tam giác $ABH$ vuông tại $H$, ta có: \\
	$\tan A=\tan \alpha=\dfrac{BH}{AH}=\dfrac{7}{4}$ \\
	Suy ra $\widehat{A}\approx60^{\circ}$ hay $\alpha\approx60^{\circ}$}
\end{ex}
%%==========Câu 105
\begin{ex}%[9H4H2-5]
	Một con thuyền với vận tốc $2 \text{ km/h}$ vượt qua một khúc sông nước chảy mạnh mất $5$ phút. Biết rằng đường đi của con thuyền tạo với bờ một góc $70^{\circ}$. Chiều rộng của khúc sông (làm tròn đến mét) là
	\shortans{$157$}
	\loigiai{% TODO: \usepackage{graphicx} required
		\begin{center}
			\includegraphics[scale=0.6]{C:/texstudio-man/Anh/{36265A5C-0362-43F9-9D6C-8DAED83A47C8}}
		\end{center}
	Ta có: $5 \text{ phút} = \dfrac{1}{12} \text{h}$ \\
	Giả sử $AB$ là đoạn đường con thuyền đi được trong $5$ phút; $BH$ là chiều rộng của khúc sông. \\
	Ta có, quãng đường $AB$ thuyền đi trong $5$ phút là: $AB=2\cdot\dfrac{1}{12}=\dfrac{1}{6} \text{km}$ \\
	+) Xét tam giác $ABH$ vuông tại $H$, ta có: $BH=AB.\sin A=\dfrac{1}{6}\cdot \sin 70^{\circ}\approx0,157 \text{ km}\approx157 \text{ m}$ \\
	Vậy chiều rộng của khúc sông (làm tròn đến mét) là $157 \text{ m}$.}
\end{ex}
%%==========Câu 106
\begin{ex}%[9H4V2-5]
	Thang xếp chữ $A$ gồm hai thang đơn tựa vào nhau. Để an toàn thì góc tạo bởi hai thang đơn phải đạt được $50^{\circ}$. Nếu muốn tạo một thang xếp chữ $A$ cao $3 \text{ m}$ tính từ mặt đất thì chiều dài mỗi thang đơn là (kết quả làm tròn đến chữ số thập phân thứ nhất )?
	\shortans{$3,3$}
	\loigiai{% TODO: \usepackage{graphicx} required
		\begin{center}
			\includegraphics[scale=0.6]{C:/texstudio-man/Anh/{76892714-37B0-462D-B38C-6347145159CB}}
		\end{center}
	Do tam giác $ABC$ cân tại $A$ nên $\widehat{B}=\dfrac{180^{\circ}-\widehat{BAC}}{2}=\dfrac{180^{\circ}-50^{\circ}}{2}=65^{\circ}$ \\
	Vì tam giác $ABC$ cân tại $A$ nên đường cao $AH$ cũng là đường trung tuyến hay $H$ là trung điểm của $BC$. \\
	Xét $\Delta ABH$ vuông tại $H$, ta có \\
	$\sin B=\dfrac{AH}{AB}$ \\
	suy ra $AB=\dfrac{AH}{\sin B}=\dfrac{3}{\sin 65^{\circ}}\approx3,3 \text{ m}$ \\
	Vậy thang đơn có chiều dài $3,3 \text{ m}$.}
\end{ex}
%%==========Câu 107
\begin{ex}%[9H4C2-5]
	\immini[thm]
	{
		Độ cao của ăng – ten $CH$ trong hình vẽ (làm tròn đến mét) là
	\shortans{$17$}
	}
	{
		% TODO: \usepackage{graphicx} required
		\includegraphics[scale=0.5]{C:/texstudio-man/Anh/{4B7C5FF3-36DA-45E2-83CA-81A2D22257E4}}
	}
	\loigiai{
	+) Xét tam giác $A C H$ vuông tại $H$, ta có: $\tan A=\dfrac{C H}{A H}$ hay $A H=\dfrac{C H}{\tan A}$\\
	Tương tự $B H=\dfrac{C H}{\tan \widehat{C B H}}$\\
	+) Có: $A B=A H-B H=\dfrac{C H}{\tan A}-\dfrac{C H}{\tan \widehat{C B H}}=C H\left(\dfrac{1}{\tan A}-\dfrac{1}{\tan \widehat{C B H}}\right)$\\
	Hay $8,5=C H\left(\frac{1}{\tan 20^{\circ}}-\frac{1}{\tan 24^{\circ}}\right)$\\
	$C H=\frac{8,5}{\left(\frac{1}{\tan 20^{\circ}}-\frac{1}{\tan 24^{\circ}}\right)} \approx 17(\mathrm{~m})$\\
	Vậy độ cao của ăng-ten $CH$ trong hình vẽ (làm tròn đến mét) là $17 \text{ m}$.}
\end{ex}
%%==========Câu 108
\begin{ex}%[9H4C2-5]
	\immini[thm]
	{
		Từ vị trí $A$ ở phía trên một tòa nhà có chiều cao $AD=68 \text{ m}$, người ta nhìn thấy vị trí $C$ cao nhất của một tháp truyền hình, góc tạo bởi tia $AC$ và tia $AH$ theo phương nằm ngang là $\widehat{CAH}=43^{\circ}$. Người ta cũng nhìn thấy chân tháp tại vị trí $B$ mà góc tạo bởi tia $AB$ và tia $AH$ là $\widehat{BAH}=28^{\circ}$, điểm $H$ thuộc đoạn thẳng $BC$ (như hình vẽ). Tổng khoảng cách $BD$ từ chân tháp đến chân tòa nhà và chiều cao $BC$ của tháp truyền hình (làm tròn kết quả đến hàng đơn vị) là
	\shortans{$315$}
	}
	{
		% TODO: \usepackage{graphicx} required
		\includegraphics[scale=0.6]{C:/texstudio-man/Anh/{8F1FDD38-629C-49A6-AF42-7E38173F5BC5}}
		
	}
	\loigiai{
	+) Xét tứ giác $ADBH$ có: $\widehat{DAH}=\widehat{ADB}=\widehat{AHB}=90^{\circ}$ \\
	Suy ra tứ giác $ADBH$ là hình chữ nhật. \\
	Do đó: $AH=DB$ và $BH=AD$ (tính chất) \\
	+) Có: $\widehat{BAD}=90^{\circ}-\widehat{BAH}=90^{\circ}-28^{\circ}=62^{\circ}$ \\
	Xét tam giác $ADB$ vuông tại $D$, ta có: $BD=AD. \tan \widehat{BAD}=68.\tan 62^{\circ}\approx128 \text{ m.}$ \\
	Suy ra: $AH=BD\approx128 \text{ m.}$ \\
	+) Xét tam giác $ACH$ vuông tại $H$, ta có: $HC=AH. \tan \widehat{CAH}\approx128. \tan 43^{\circ} \approx 119 \text{ m.}$ \\
	+) Có: $BC=BH+HC$ \\
	Mà $BH=AD$ nên $BC\approx68+119\approx187 \text{ m.}$.\\
	Suy ra tổng khoảng cách $128+187=315$}
\end{ex}
\Closesolutionfile{ans}
\indapan{6}{ans/ans9K1-12-TLN-D6}
\subsubsection{Bài tập đúng sai}
\Opensolutionfile{ans}[ans/ans9K1-12-DS-1112]
%%==========Câu 109
\begin{ex}%[9H4H2-4]
Cho $\triangle ABC$ vuông tại $A$. Điền Đ cho câu trả lời đúng và S cho câu trả lời sai:
\choiceTF
{\True $AB=BC \cdot \sin C$}
{$AC=AB \cdot \tan C$}
{$BC=\dfrac{AB}{\sin B}$}
{\True $\sin B=\dfrac{AC}{BC}$}
\loigiai{
\begin{itemchoice}
\itemch Xét $\triangle ABC$ vuông tại $A$ ta có $AB=BC \cdot \sin C$. Vậy mệnh đề đúng.
\itemch Xét $\triangle ABC$ vuông tại $A$ ta có $AC=AB \cdot \cot C$. Vậy mệnh đề sai.
\itemch Xét $\triangle ABC$ vuông tại $A$ ta có $AB=BC \cdot \sin C$ suy ra $BC=\dfrac{AB}{\sin C}$. Vậy mệnh đề sai.
\itemch Xét $\triangle ABC$ vuông tại $A$ ta có $\sin B=\dfrac{AC}{BC}$. Vậy mệnh đề đúng.
\end{itemchoice}
}
\end{ex}
%%==========Câu 110
\begin{ex}%[9H4H2-4]
Cho $\triangle ABC$ vuông cân tại $A$. Điền Đ cho câu trả lời đúng và S cho câu trả lời sai:
\choiceTF
{\True $\sin B=\dfrac{\sqrt{2}}{2}$}
{\True $AB=\dfrac{\sqrt{2}}{2}BC$}
{\True $BC=\sqrt{2}\cdot AB$}
{$\tan B=\sqrt{2}$}
\loigiai{
\begin{itemchoice}
\itemch Vì $\triangle ABC$ vuông cân tại $A$ nên $\widehat{B}=\widehat{C}=45^{\circ}$ suy ra $\sin B=\sin 45^{\circ}=\dfrac{\sqrt{2}}{2}.$ Vậy mệnh đề đúng.
\itemch Xét $\triangle ABC$ vuông tại $A$ ta có $AB=BC \cdot \sin C=BC \cdot \sin 45^{\circ}=\dfrac{\sqrt{2}}{2}BC.$ Vậy mệnh đề đúng.
\itemch Xét $\triangle ABC$ vuông tại $A$ ta có $AB=BC \cdot \sin C$ suy ra $BC=\dfrac{AB}{\sin C}=\dfrac{AB}{\sin 45^{\circ}}=AB\sqrt{2}.$ Vậy mệnh đề đúng.
\itemch Vì $\triangle ABC$ vuông cân tại $A$ nên $\widehat{B}=\widehat{C}=45^{\circ}$ suy ra $\tan B=\tan 45^{\circ}=1$. Vậy mệnh đề sai.
\end{itemchoice}
}
\end{ex}
%%==========Câu 111
\begin{ex}
	Cho $\triangle ABC$ vuông tại $A$ đường cao $AH$. Điền Đ cho câu trả lời đúng và S cho câu trả lời sai:
	\choiceTF
	{\True $AB \cdot \sin B=AC \cdot \sin C$}
	{\True $AH=AC \cdot \sin C$}
	{$BH=AH \cdot \tan C$}
	{\True $BC=AB \cdot \cos B+AC \cdot \cos C$}
	\loigiai{
	\begin{itemchoice}
	\itemch Xét $\triangle ABC$ vuông tại $A$ ta có $AB \cdot \sin B=AB \cdot \dfrac{AC}{BC}=\dfrac{AB \cdot AC}{BC}$. $AC \cdot \sin C=AC \cdot \dfrac{AB}{BC}=\dfrac{AB \cdot AC}{BC}$. Vậy mệnh đề đúng.
	\itemch Xét $\triangle AHC$ vuông tại $H$ ta có $AH=AC \cdot \sin C$. Vậy mệnh đề đúng. \\
	\itemch Xét $\triangle AHB$ vuông tại $H$ ta có $BH=AH \cdot \tan B$. Vậy mệnh đề sai. \\
	\itemch Xét $\triangle AHB$ vuông tại $H$ ta có $BH=AB \cdot \cos B$.\\
	 Xét $\triangle AHC$ vuông tại $H$ ta có $CH=AC \cdot \cos C$. Suy ra $BC=BH+CH=AB \cdot \cos B+AC \cdot \cos C$. Vậy mệnh đề đúng.
	\end{itemchoice}
	}
\end{ex}
%%==========Câu 112
\begin{ex}
	Cho $\triangle MNP$ vuông tại $P$. Điền Đ cho câu trả lời đúng và S cho câu trả lời sai:
	\choiceTF
	{$MP=NP \cdot \tan M$}
	{\True $PN=MN \cdot \sin M$}
	{\True $MN=\dfrac{NP}{\cos N}$}
	{$\cot P=\dfrac{MN}{NP}$}
	\loigiai{
	\begin{itemchoice}
	\itemch Xét $\triangle MNP$ vuông tại $P$ ta có $MP=NP \cdot \tan N$. Vậy mệnh đề sai.
	\itemch Xét $\triangle MNP$ vuông tại $P$ ta có $PN=MN \cdot \sin M$. Vậy mệnh đề đúng.
	\itemch Xét $\triangle MNP$ vuông tại $P$ ta có $PN=MN \cdot \cos N$ suy ra $MN=\dfrac{NP}{\cos N}$. Vậy mệnh đề đúng.
	\itemch Xét $\triangle MNP$ vuông tại $P$ . Vậy mệnh đề sai.\\
	\end{itemchoice}
	}
\end{ex}
%%==========Câu 113
\begin{ex}
	Cho $\triangle ABC$ vuông tại $A$ có $\widehat{C}=30^{\circ}$. Điền Đ cho câu trả lời đúng và S cho câu trả lời sai:
	\choiceTF
	{\True $AB=\dfrac{1}{2}BC$}
	{\True $\widehat{B}=60^{\circ}$}
	{\True $AC=\sqrt{3}\cdot AB$}
	{$BC=\sqrt{3}\cdot AC$}
	\loigiai{
	\begin{itemchoice}
	\itemch $\triangle ABC$ vuông tại $A$ có $AB=BC \cdot \sin C=BC \cdot \sin 30^{\circ}=\dfrac{1}{2}BC.$ Vậy mệnh đề đúng.
	\itemch $\triangle ABC$ vuông tại $A$ có $\widehat{B}=90^{\circ}-\widehat{C}=90^{\circ}-30^{\circ}=60^{\circ}.$ Vậy mệnh đề đúng.
	\itemch $\triangle ABC$ vuông tại $A$ có $AC=AB \cdot \cot C=AB \cdot \cot 30^{\circ}=\sqrt{3} \cdot AB$. Vậy mệnh đề đúng.
	\itemch $\triangle ABC$ vuông tại $A$ có $AC=BC \cdot \cos C$ suy ra $BC=\dfrac{AC}{\cos C}=\dfrac{AC}{\cos 30^{\circ}}=\dfrac{2AC}{\sqrt{3}}$. Vậy mệnh đề sai.
	\end{itemchoice}
	}
\end{ex}
%%==========Câu 114
\begin{ex}
	Cho $\triangle ABC$ vuông tại $A$ đường cao $AH$. Điền Đ cho câu trả lời đúng và S cho câu trả lời sai:
	\choiceTF
	{\True $AC=AB \cdot \tan B=AB \cdot \cot C$}
	{$AB=AC \cdot \cot B=AC \cdot \sin C$}
	{\True $BC=AH \cdot (\cot B+\cot C)$}
	{\True $S_{ABC}=\dfrac{1}{2} AC \cdot BC \cdot \sin C$}
	\loigiai{
	\begin{itemchoice}
	\itemch Xét $\triangle ABC$ vuông tại $A$ ta có $AC=AB \cdot \tan B$. Vì $\widehat{B}$ và $\widehat{C}$ phụ nhau nên $\tan B = \cot C$. Do đó $AC=AB \cdot \cot C$. Vậy mệnh đề đúng.
	\itemch Xét $\triangle A B C$ vuông tại $A$ ta có $A B=A C \cdot \tan B=B C \cdot \sin C$. 
	\itemch Xét $\triangle AHC$ vuông tại $H$ ta có $CH=AH \cdot \cot C$.\\
	Xét $\triangle AHB$ vuông tại $H$ ta có $BH=AH \cdot \cot B$.\\
	Suy ra $BC=BH+CH=AH \cdot \cot B+AH \cdot \cot C=AH \cdot (\cot B+\cot C)$. Vậy mệnh đề đúng.
	\itemch Xét $\triangle AHC$ vuông tại $H$ ta có $AH=AC \cdot \sin C$.\\
	Suy ra $S_{ABC}=\dfrac{1}{2}AH \cdot BC = \dfrac{1}{2} (AC \cdot \sin C) \cdot BC = \dfrac{1}{2} AC \cdot BC \cdot \sin C$. Vậy mệnh đề đúng.
	\end{itemchoice}
	}
\end{ex}
%%==========Câu 115
\begin{ex}
	Cho $\triangle ABC$ vuông tại $A$ đường cao $AH$. Điền Đ cho câu trả lời đúng và S cho câu trả lời sai:
	\choiceTF
	{$BC=AB \cdot \cos B$}
	{$AB=AH \cdot \sin B$}
	{\True $AH=AC \cdot \sin C$}
	{\True $\dfrac{AB}{AC}=\dfrac{\sin C}{\sin B}$}
	\loigiai{
	\begin{itemchoice}
	\itemch Xét $\triangle ABC$ vuông tại $A$ ta có $AB=BC \cdot \cos B$. Vậy mệnh đề sai.
	\itemch Xét $\triangle AHB$ vuông tại $H$ ta có $AH=AB \cdot \sin B$. Vậy mệnh đề sai.
	\itemch Xét $\triangle AHC$ vuông tại $H$ ta có $AH=AC \cdot \sin C$. Vậy mệnh đề đúng.
	\itemch Xét $\triangle ABC$ vuông tại $A$ ta có $\dfrac{\sin C}{\sin B}=\dfrac{AB/BC}{AC/BC}=\dfrac{AB}{AC}.$ Vậy mệnh đề đúng.
	\end{itemchoice}
	}
\end{ex}
%%==========Câu 116
\begin{ex}
	Cho $\triangle ABC$ vuông tại $A$ đường cao $AH$. $\widehat{B}=60^{\circ}$. Điền Đ cho câu trả lời đúng và S cho câu trả lời sai:
	\choiceTF
	{\True $\widehat{C}=30^{\circ}$}
	{\True $AB=\dfrac{\sqrt{3}}{3}\cdot AC$}
	{$AH=\dfrac{AB}{2}$}
	{\True $BC=2 \cdot AB$}
	\loigiai{
	\begin{itemchoice}
	\itemch $\triangle ABC$ vuông tại $A$ có $\widehat{C}=90^{\circ}-\widehat{B}=90^{\circ}-60^{\circ}=30^{\circ}$. Vậy mệnh đề đúng.
	\itemch $\triangle ABC$ vuông tại $A$ có $AB = AC \cdot \tan C = AC \cdot \tan 30^{\circ}=\dfrac{\sqrt{3}}{3} AC$. Vậy mệnh đề đúng.
	\itemch $\triangle ABH$ vuông tại $H$ có $AH=AB \cdot \sin B=AB \cdot \sin 60^{\circ}=\dfrac{\sqrt{3}}{2}AB$. Vậy mệnh đề sai.
	\itemch $\triangle ABC$ vuông tại $A$ có $AB=BC \cdot \sin C=BC \cdot \sin 30^{\circ}=\dfrac{1}{2}BC$ suy ra $BC=2 \cdot AB$. Vậy mệnh đề đúng.
	\end{itemchoice}
	}
\end{ex}
%%==========Câu 117
\begin{ex}
	Cho $\triangle ABC$ vuông tại $A$ có $AB=3\,\text{cm}$, $AC=6\,\text{cm}$. Điền Đ cho câu trả lời đúng và S cho câu trả lời sai:
	\choiceTF
	{\True $\widehat{C}\approx27^{\circ}$}
	{\True $\sin B=\dfrac{2\sqrt{5}}{5}$}
	{$\widehat{B}\approx53^{\circ}$}
	{$\dfrac{AB}{BC}=\dfrac{1}{5}$}
	\loigiai{
	\begin{itemchoice}
	\itemch Xét $\triangle ABC$ vuông tại $A$ có: $\tan C=\dfrac{AB}{AC}=\dfrac{3}{6}=\dfrac{1}{2}$ suy ra $\widehat{C}\approx27^{\circ}$. Vậy mệnh đề đúng.
	\itemch Xét $\triangle ABC$ vuông tại $A$ có: $BC=\sqrt{AB^2+AC^2}=\sqrt{3^2+6^2}=3\sqrt{5}\,\text{cm}$. $\sin B=\dfrac{AC}{BC}=\dfrac{6}{3\sqrt{5}}=\dfrac{2\sqrt{5}}{5}.$ Vậy mệnh đề đúng.
	\itemch Xét $\triangle ABC$ vuông tại $A$ có: $\widehat{B}=90^{\circ}-\widehat{C}\approx90^{\circ}-27^{\circ}=63^{\circ}$. Vậy mệnh đề sai.
	\itemch Xét $\triangle ABC$ vuông tại $A$ có: $BC=3\sqrt{5}\,\text{cm}$ suy ra $\dfrac{AB}{BC}=\dfrac{3}{3\sqrt{5}}=\dfrac{\sqrt{5}}{5}.$ Vậy mệnh đề sai.
	\end{itemchoice}
	}
\end{ex}
%%==========Câu 118
\begin{ex}
	Cho $\triangle ABC$ có $AB=5\,\text{cm}$, $AC=13\,\text{cm}$, $BC=12\,\text{cm}$. Điền Đ cho câu trả lời đúng và S cho câu trả lời sai:
	\choiceTF
	{$\triangle ABC$ vuông tại $A$}
	{$\sin B=\dfrac{13}{12}$}
	{\True $\widehat{C}\approx23^{\circ}$}
	{$\widehat{B}\approx67^{\circ}22^{\prime}$}
	\loigiai{
	\begin{itemchoice}
	\itemch Ta có $AC^{2}=13^{2}=169$; $AB^{2}+BC^{2}=5^{2}+12^{2}=25+144=169$. Suy ra $AC^{2}=AB^{2}+BC^{2}$. Suy ra $\triangle ABC$ vuông tại $B$. Vậy mệnh đề sai.
	\itemch Xét $\triangle ABC$ vuông tại $B$ Vậy mệnh đề sai.\\
	\itemch Xét $\triangle ABC$ vuông tại $B$ (theo lời giải gốc) ta có $\sin C=\dfrac{AB}{AC}=\dfrac{5}{13}.$ Suy ra $\widehat{C}\approx23^{\circ}$. Vậy mệnh đề đúng.\\
	\itemch Xét $\triangle ABC$ vuông tại $B$ Suy ra $\widehat{B}=90^\circ$. Vậy mệnh đề sai.
	\end{itemchoice}
	}
\end{ex}
%%==========Câu 119
\begin{ex}
	Cho $\triangle ABC$ vuông tại $A$, đường cao $AH$. Biết $BH=4\,\text{cm}$, $CH=9\,\text{cm}$. Điền Đ cho câu trả lời đúng và S cho câu trả lời sai:
	\choiceTF
	{\True $AH=6\,\text{cm}$}
	{$\widehat{B}\approx56^{\circ}18^{\prime}$}
	{\True $\widehat{C}\approx33^{\circ}41^{\prime}$}
	{\True $AC=3\sqrt{13}\,\text{cm}$}
	\loigiai{
	\begin{itemchoice}
	\itemch Xét $\triangle ABC$ vuông tại $A$, đường cao $AH$: $AH^{2}=BH \cdot HC=4 \cdot 9=36$. Suy ra $AH=6\,\text{cm}$. Vậy mệnh đề đúng.
	\itemch Xét $\triangle ABH$ vuông tại $H$ ta có $\tan B=\dfrac{AH}{BH}=\dfrac{6}{4}=1,5$. Suy ra $\widehat{B}\approx56^{\circ}19^{\prime}$. Vậy mệnh đề sai.
	\itemch Xét $\triangle ABC$ vuông tại $A$ (hoặc $\triangle AHC$ vuông tại $H$ và $\widehat{C}=90^\circ-\widehat{B}$): $\widehat{C}=90^{\circ}-\widehat{B}\approx90^{\circ}-56^{\circ}19^{\prime}=33^{\circ}41^{\prime}$. Vậy mệnh đề đúng.
	\itemch Xét $\triangle ABC$ vuông tại $A$, đường cao $AH$: $AC^{2}=CH \cdot BC=CH \cdot (BH+CH) = 9 \cdot (4+9)=9 \cdot 13=117$. Suy ra $AC=\sqrt{117}=3\sqrt{13}\,\text{cm}$. Vậy mệnh đề đúng.
	\end{itemchoice}
	}
\end{ex}
%%==========Câu 120
\begin{ex}
	Cho $\triangle ABC$ vuông tại $A$, đường cao $AH$. Biết $AB=3\,\text{cm}$, $AH=2,4\,\text{cm}$. Điền Đ cho câu trả lời đúng và S cho câu trả lời sai:
	\choiceTF
	{\True $\widehat{BAH}\approx37^{\circ}$}
	{\True $\widehat{B}\approx53^{\circ}$}
	{$BC=4\,\text{cm}$}
	{$S_{ABC}=4,8\,\text{cm}^{2}$}
	\loigiai{
	\begin{itemchoice}
	\itemch Xét $\triangle ABH$ vuông tại $H$ ta có $\cos\widehat{BAH}=\dfrac{AH}{AB}=\dfrac{2,4}{3}=\dfrac{4}{5}.$ Suy ra $\widehat{BAH}\approx37^{\circ}$. Vậy mệnh đề đúng.
	\itemch Xét $\triangle ABH$ vuông tại $H$ ta có $\sin\widehat{B}=\dfrac{AH}{AB}=\dfrac{2,4}{3}=\dfrac{4}{5}.$ Suy ra $\widehat{B}\approx53^{\circ}$. Vậy mệnh đề đúng.
	\itemch Xét $\triangle ABH$ vuông tại $H$ ta có $BH=\sqrt{AB^{2}-AH^{2}}=\sqrt{3^{2}-2,4^{2}}=\sqrt{9-5,76}=\sqrt{3,24}=1,8\,\text{cm}.$ Xét $\triangle ABC$ vuông tại $A$ ta có $AB^{2}=BH \cdot BC$ suy ra $BC=\dfrac{AB^{2}}{BH}=\dfrac{3^{2}}{1,8}=\dfrac{9}{1,8}=5\,\text{cm}.$ Vậy mệnh đề sai.
	\itemch $S_{ABC}=\dfrac{1}{2}AH \cdot BC=\dfrac{1}{2} \cdot 2,4 \cdot 5 = \dfrac{12}{2} = 6\,\text{cm}^{2}.$ Vậy mệnh đề sai.
	\end{itemchoice}
	}
\end{ex}
%%==========Câu 121
\begin{ex}
	Cho $\triangle ABC$ có đường cao $AH$. Biết $BH=2\,\text{cm}$, $CH=3\,\text{cm}$, $\widehat{B}=45^{\circ}$. Điền Đ cho câu trả lời đúng và S cho câu trả lời sai:
	\choiceTF
	{$AH=\sqrt{6}\,\text{cm}$}
	{\True $\widehat{BAH}=45^{\circ}$}
	{$\widehat{C}\approx33^{\circ}$}
	{\True $\widehat{BAC}\approx101^{\circ}$}
	\loigiai{
	\begin{itemchoice}
	\itemch Xét $\triangle ABH$ có $AH$ là đường cao nên $\triangle ABH$ vuông tại $H$. $AH=BH \cdot \tan B=2 \cdot \tan 45^{\circ}=2 \cdot 1 = 2\,\text{cm}.$ Vậy mệnh đề sai.
	\itemch Xét $\triangle ABH$ vuông tại $H$ ta có $\widehat{BAH}=90^{\circ}-\widehat{B}=90^{\circ}-45^{\circ}=45^{\circ}.$ Vậy mệnh đề đúng.
	\itemch Xét $\triangle ACH$ vuông tại $H$ ta có $\tan\widehat{C}=\dfrac{AH}{CH}=\dfrac{2}{3}.$ Suy ra $\widehat{C}\approx33,69^\circ \approx 34^{\circ}$. Vậy mệnh đề ($\widehat{C}\approx33^{\circ}$) sai.
	\itemch Ta có $\widehat{BAC}= \widehat{BAH} + \widehat{CAH}$. $\widehat{CAH} = 90^\circ - \widehat{C} \approx 90^\circ - 34^\circ = 56^\circ$. $\widehat{BAC} \approx 45^\circ + 56^\circ = 101^\circ$. (Hoặc $\widehat{BAC}=180^{\circ}-\widehat{B}-\widehat{C}\approx180^{\circ}-45^{\circ}-34^{\circ}=101^{\circ}$). Vậy mệnh đề đúng.
	\end{itemchoice}
	}
\end{ex}
%%==========Câu 122
\begin{ex}
	Cho $\triangle ABC$ vuông tại $A$, phân giác góc $B$ cắt $AC$ tại $D$. Biết $AD=2,5\,\text{cm}$, $DC=5\,\text{cm}$. Điền Đ cho câu trả lời đúng và S cho câu trả lời sai:
	\choiceTF
	{$\tan C=\dfrac{1}{2}$}
	{$\widehat{C}\approx27^{\circ}$}
	{\True $\widehat{ABD}=30^{\circ}$}
	{\True $\widehat{BDC}=120^{\circ}$}
	\loigiai{
	\begin{itemchoice}
	\itemch Xét $\triangle ABC$ vuông tại $C$ (thực tế là vuông tại $A$) ta có $\sin C=\dfrac{AC}{BC}.$ Vì $BD$ là phân giác $\widehat{ABC}$ nên ta có $\dfrac{DA}{DC}=\dfrac{AB}{BC}=\dfrac{2,5}{5}=\dfrac{1}{2}$ Suy ra $\sin C=\dfrac{1}{2}.$ 
	\itemch Ta có $\sin C=\dfrac{1}{2}$ suy ra $\widehat{C}=30^{\circ}$. Vậy mệnh đề sai (vì $27^\circ \neq 30^\circ$). 
	\itemch Xét $\triangle ABC$ vuông tại $A$ ta có $\widehat{ABC}=90^{\circ}-\widehat{C}=90^{\circ}-30^{\circ}=60^{\circ}$ (sử dụng $\widehat{C}=30^\circ$ từ lời giải trước). Vì $BD$ là phân giác $\widehat{ABC}$ nên ta có $\widehat{ABD}=\dfrac{\widehat{ABC}}{2}=\dfrac{60^{\circ}}{2}=30^{\circ}.$ Vậy mệnh đề đúng.
	\itemch Ta có $\widehat{BDC}=180^{\circ}-\widehat{DBC}-\widehat{C}\approx180^{\circ}-30^{\circ}-30^{\circ}=120^{\circ}$ (sử dụng $\widehat{C}=30^\circ$ và $\widehat{DBC}=30^\circ$ từ lời giải trước). Vậy mệnh đề đúng.
	\end{itemchoice}
	}
\end{ex}
%%==========Câu 123
\begin{ex}
	\immini[thm]
	{
	Một chiếc cầu trượt cho trẻ em được thiết kế bao gồm một cầu thang đi lên dài $3\,\text{m}$, một máng trượt xuống dài $5\,\text{m}$, biết độ cao của cầu trượt là $2\,\text{m}$. Gọi $AH$ là độ cao ($AH=2\,\text{m}$), $AC$ là chiều dài cầu thang ($AC=3\,\text{m}$), $AB$ là chiều dài máng trượt ($AB=5\,\text{m}$). Điền Đ cho câu trả lời đúng và S cho câu trả lời sai:
	}
	{
	% TODO: \usepackage{graphicx} required
	\includegraphics[scale=0.35]{C:/texstudio-man/Anh/{990083B7-60BE-429D-919D-F2A3A0CE3142}}
	}
	\choiceTF
	{\True Góc tạo bởi cầu thang và mặt đất làm tròn đến độ là $42^{\circ}$}
	{Góc tạo bởi máng trượt và mặt đất làm tròn đến độ là $23^{\circ}$}
	{\True Góc tạo bởi máng trượt và cầu thang làm tròn đến độ là $114^{\circ}$}
	{\True Khoảng cách giữa chân cầu thang và máng trượt khoảng $6,8\,\text{m}$}
	\loigiai{
	\begin{itemchoice}
	\itemch Góc tạo bởi cầu thang và mặt đất là $\widehat{C}$. Xét $\triangle AHC$ vuông tại $H$ ta có $\sin\widehat{C}=\dfrac{AH}{AC}=\dfrac{2}{3}.$ Suy ra $\widehat{C}\approx41,81^\circ \approx 42^{\circ}$. Vậy mệnh đề đúng.
	\itemch Góc tạo bởi máng trượt và mặt đất là $\widehat{B}$. Xét $\triangle ABH$ vuông tại $H$ ta có $\sin\widehat{B}=\dfrac{AH}{AB}=\dfrac{2}{5}.$ Suy ra $\widehat{B}\approx23,58^\circ \approx 24^{\circ}$. Mệnh đề nói $\approx 23^\circ$. Vậy mệnh đề sai. 
	\itemch Góc tạo bởi máng trượt và cầu thang là $\widehat{BAC}$. Ta có $\widehat{BAC}=180^{\circ}-\widehat{B}-\widehat{C}\approx180^{\circ}-23,58^{\circ}-41,81^{\circ} \approx 180^\circ - 65,39^\circ \approx 114,61^\circ \approx 115^{\circ}$. 
	\itemch
	 Khoảng cách giữa chân cầu thang và máng trượt là $BC = BH+CH$.\\
	 Xét $\triangle ABH$ vuông tại $H$:\\ $BH=\sqrt{AB^{2}-AH^{2}}=\sqrt{5^{2}-2^{2}}=\sqrt{25-4}=\sqrt{21} \approx 4,583\,\text{m}.$\\
	 Xét $\triangle ACH$ vuông tại $H$:\\ $CH=\sqrt{AC^{2}-AH^{2}}=\sqrt{3^{2}-2^{2}}=\sqrt{9-4}=\sqrt{5} \approx 2,236\,\text{m}.$\\
	 $BC = BH+CH \approx 4,583 + 2,236 = 6,819\,\text{m} \approx 6,8\,\text{m}$. Vậy mệnh đề đúng.
	\end{itemchoice}
	}
\end{ex}
%%==========Câu 124
\begin{ex}
	Cho $\triangle ABC$ vuông tại $A$, phân giác góc $B$ cắt $AC$ tại $D$, phân giác góc ngoài tại đỉnh $B$ cắt $AC$ tại $E$. Biết $AB=6\,\text{cm}$, $BC=10\,\text{cm}$. Điền Đ cho câu trả lời đúng và S cho câu trả lời sai:
	\choiceTF
	{\True $\widehat{C}\approx37^{\circ}$}
	{\True $AD=3\,\text{cm}$}
	{$\widehat{ADB}=53^{\circ}$}
	{$\widehat{EBC}=120^{\circ}$}
	\loigiai{
	\begin{itemchoice}
	\itemch Xét $\triangle ABC$ vuông tại $A$ ta có $\sin\widehat{C}=\dfrac{AB}{BC}=\dfrac{6}{10}=\dfrac{3}{5}.$ Suy ra $\widehat{C}\approx36,87^\circ \approx 37^{\circ}$. Vậy mệnh đề đúng.
	\itemch Xét $\triangle ABC$ vuông tại $A$, nên ta có $AC=\sqrt{BC^{2}-AB^{2}}=\sqrt{10^{2}-6^{2}}=\sqrt{64}=8\,\text{cm}.$
	%% Sửa lại phần áp dụng tính chất đường phân giác từ lời giải gốc:
	Vì $BD$ là tia phân giác trong của $\widehat{B}$ nên $\dfrac{AD}{DC} = \dfrac{AB}{BC} = \dfrac{6}{10} = \dfrac{3}{5}$. Mà $AD+DC=AC=8$.
	$\dfrac{AD}{3} = \dfrac{DC}{5} = \dfrac{AD+DC}{3+5} = \dfrac{8}{8}=1$. Suy ra $AD=3 \cdot 1 = 3\,\text{cm}$. Vậy mệnh đề đúng.
	\itemch Xét $\triangle ABD$ vuông tại $A$ ta có $\tan\widehat{ADB}=\dfrac{AB}{AD}=\dfrac{6}{3}=2.$ Suy ra $\widehat{ADB}\approx63,43^\circ \approx 63^{\circ}$. Vậy mệnh đề ($\widehat{ADB}=53^{\circ}$) sai.
	\itemch Vì $BD$ và $BE$ là phân giác trong và phân giác ngoài của góc $B$ nên $BD\perp BE$. Xét $\triangle BED$ vuông tại $B$ có $\widehat{E}=90^{\circ}-\widehat{BDA}\approx90^{\circ}-63^{\circ}=27^{\circ}$. $\widehat{EBC}=180^{\circ}-\widehat{E}-\widehat{C}\approx180^{\circ}-27^{\circ}-37^{\circ}=116^{\circ}$. Vậy mệnh đề sai. \\
	\end{itemchoice}
	}
\end{ex}
%%==========Câu 125
\begin{ex}
	Cho $\triangle ABC$ vuông tại $A$ có $AB=5\,\text{cm}$, $\widehat{B}=30^{\circ}$. Điền Đ vào các câu trả lời đúng, S vào các câu trả lời sai.
	\choiceTF
	{\True $\widehat{C}=60^{\circ}$}
	{$AC=2,5\,\text{cm}$}
	{$BC=10\,\text{cm}$}
	{\True $S_{ABC}=\dfrac{25\sqrt{3}}{6}\,\text{cm}^{2}$}
	\loigiai{
	\begin{itemchoice}
	\itemch %% Lời giải đúng (theo Đáp án Đ và đề bài $\widehat{B}=30^\circ$):
	Xét $\triangle ABC$ vuông tại $A$ có: $\widehat{C}=90^{\circ}-\widehat{B}=90^{\circ}-30^{\circ}=60^{\circ}.$ Vậy mệnh đề đúng.
	\itemch %% Lời giải đúng (theo Đáp án S và đề bài $\widehat{B}=30^\circ$):
	Xét $\triangle ABC$ vuông tại $A$ có: $AC=AB \cdot \tan B=5 \cdot \tan 30^{\circ}=5 \cdot \dfrac{1}{\sqrt{3}}=\dfrac{5\sqrt{3}}{3}\,\text{cm} \approx 2,887\,\text{cm}$.
	Mệnh đề $AC=2,5\,\text{cm}$ là sai.
	\itemch %% Lời giải đúng (theo Đáp án S và đề bài $\widehat{B}=30^\circ$):
	Xét $\triangle ABC$ vuông tại $A$ có: $BC=\dfrac{AB}{\cos B}=\dfrac{5}{\cos 30^{\circ}}=\dfrac{5}{\sqrt{3}/2}=\dfrac{10}{\sqrt{3}}=\dfrac{10\sqrt{3}}{3}\,\text{cm} \approx 5,77\,\text{cm}$.
	Mệnh đề $BC=10\,\text{cm}$ là sai.
	\itemch %% Lời giải đúng (theo Đáp án Đ và đề bài $\widehat{B}=30^\circ$):
	$AC = \dfrac{5\sqrt{3}}{3}\,\text{cm}$.
	$S_{ABC}=\dfrac{1}{2} AB \cdot AC = \dfrac{1}{2} \cdot 5 \cdot \dfrac{5\sqrt{3}}{3} = \dfrac{25\sqrt{3}}{6}\,\text{cm}^{2}.$ Vậy mệnh đề đúng.
	\end{itemchoice}
	}
\end{ex}
%%==========Câu 126
\begin{ex}
	Cho $\triangle IPK$ có $IP=8\,\text{cm}$, $\widehat{P}=30^{\circ},$ $\widehat{K}=45^{\circ}$, kẻ đường cao $IH$. Điền Đ vào các câu trả lời đúng, S vào các câu trả lời sai.
	\choiceTF
	{\True $IH = 4\,\text{cm}$}
	{$PH=4\sqrt{2}\,\text{cm}$}
	{\True $IK=4\sqrt{2}\,\text{cm}$}
	{$PK=12\,\text{cm}$}
	\loigiai{
	\begin{itemchoice}
	\itemch Xét $\triangle IPH$ vuông tại $H$ ta có: $IH=IP \cdot \sin P=8 \cdot \sin 30^{\circ}=4\,\text{cm}.$ Vậy mệnh đề đúng.
	\itemch Xét $\triangle IPH$ vuông tại $H$ ta có: $PH=IP \cdot \cos P=8 \cdot \cos 30^{\circ}=4\sqrt{3}\,\text{cm}$. Vậy mệnh đề sai.
	\itemch Xét $\triangle IKH$ vuông tại $H$ ta có: $IH=IK \cdot \sin K$ suy ra $IK=\dfrac{IH}{\sin K}=\dfrac{4}{\sin 45^{\circ}}=4\sqrt{2}\,\text{cm}.$ Vậy mệnh đề đúng.
	\itemch Xét $\triangle IKH$ vuông tại $H$ ta có: $HK=IH \cdot \cot K=4 \cdot \cot 45^{\circ}=4\,\text{cm}.$ $PK=PH+HK=4\sqrt{3}+4 \approx 4(1,732) + 4 = 6,928+4 = 10,928\,\text{cm}$. Vậy mệnh đề sai.
	\end{itemchoice}
	}
\end{ex}
%%==========Câu 127
\begin{ex}
	Cho tam giác $ABC$ vuông tại $A$, đường cao $AH$. Biết $AB=2,5\,\text{cm}$, $BH=1,5\,\text{cm}$. Điền Đ (Đúng), S (Sai) cho các phát biểu.
	\choiceTF
	{$\widehat{B}\approx53^{\circ}7^{\prime}$} % Theo lời giải gốc là Sai, Đáp án gốc Đ
	{\True $AC\approx3,3\,\text{cm}$} % Theo lời giải gốc là Đúng, Đáp án gốc S
	{$\text{BC}=\dfrac{AB}{\sin B}$} % Theo lời giải gốc là Sai, Đáp án gốc Đ
	{\True $AH=2\,\text{cm}$} % Theo lời giải gốc là Đúng, Đáp án gốc Đ
	\loigiai{
	\begin{itemchoice}
	\itemch Xét tam giác $ABH$ vuông tại $H$ ta có $\cos\widehat{B}=\dfrac{BH}{AB}=\dfrac{1,5}{2,5}=0,6$. Suy ra $\widehat{B}\approx53^{\circ}8^{\prime}$. Vậy mệnh đề ($\widehat{B}\approx53^{\circ}7^{\prime}$) sai.
	\itemch Ta có $\widehat{C}=90^{\circ}-\widehat{B} \approx 90^{\circ}-53^{\circ}8^{\prime}=36^{\circ}52^{\prime}$. Xét $\triangle ABC$ vuông tại $A$, ta có $AC=AB \cdot \tan\widehat{B}=2,5 \cdot \tan 53^{\circ}8^{\prime}\approx2,5 \cdot \dfrac{4}{3} \approx 3,33\,\text{cm}.$ Vậy mệnh đề ($AC\approx3,3\,\text{cm}$) đúng (làm tròn).
	\itemch Xét $\triangle ABC$ vuông tại $A$, ta có $BC=\dfrac{AB}{\cos B}$. Vậy mệnh đề ($BC=\dfrac{AB}{\sin B}$) sai.
	\itemch Xét tam giác $ABH$ vuông tại $H$ ta có $AH=\sqrt{AB^{2}-BH^{2}}=\sqrt{2,5^{2}-1,5^{2}}=\sqrt{6,25-2,25}=\sqrt{4}=2\,\text{cm}.$ Vậy mệnh đề đúng.
	\end{itemchoice}
	}
\end{ex}
%%==========Câu 128
\begin{ex}
	\immini[thm]
	{
	Cho hình vẽ, biết $\widehat{A}=35^{\circ}$, $BC=2\,\text{cm}$, $AB=3\,\text{cm}$. (Giả sử $A, B, C$ thẳng hàng theo thứ tự đó để $AC=AB+BC=5\,\text{cm}$, và $BH \perp AK, CK \perp AK$ với $K$ là điểm trên đường thẳng chứa $A$). Điền Đ vào các câu trả lời đúng, S vào các câu trả lời sai.
	\choiceTF
	{\True $HB\approx1,7\,\text{cm}$}
	{$AH\approx2,45\,\text{cm}$}
	{\True $CK\approx2,9\,\text{cm}$}
	{\True $AK\approx4,1\,\text{cm}$}
	}
	{
	% TODO: \usepackage{graphicx} required
	\includegraphics[scale=0.5]{C:/texstudio-man/Anh/{7CE18D1A-707C-4E1C-9BED-7A48B87B139F}}
	}
	\loigiai{
	\begin{itemchoice}
	\itemch Xét tam giác $ABH$ vuông tại $H$ ta có $HB=AB \cdot \sin A = 3 \cdot \sin 35^{\circ}\approx1,72\,\text{cm}.$ Vậy mệnh đề ($HB\approx1,7\,\text{cm}$) đúng.
	\itemch Xét tam giác $ABH$ vuông tại $H$ ta có $AH=AB \cdot \cos A = 3 \cdot \cos 35^{\circ}\approx2,457\,\text{cm}.$ Vậy mệnh đề ($AH\approx2,45\,\text{cm}$) sai (vì $2,457 \approx 2,46$).
	\itemch Gọi $AC = AB+BC = 3+2=5\,\text{cm}$. Xét tam giác $ACK$ vuông tại $K$ ta có $CK=AC \cdot \sin A = 5 \cdot \sin 35^{\circ}\approx2,868\,\text{cm}.$ Vậy mệnh đề ($CK\approx2,9\,\text{cm}$) đúng.
	\itemch Xét tam giác $ACK$ vuông tại $K$ ta có $AK=AC \cdot \cos A = 5 \cdot \cos 35^{\circ}\approx4,096\,\text{cm}.$ Vậy mệnh đề ($AK\approx4,1\,\text{cm}$) đúng.
	\end{itemchoice}
	}
\end{ex}
%%==========Câu 129
\begin{ex}
	Cho tam giác $ABC$ cân tại $A$, đường cao $BH$ ($BH \perp AC$). Biết $\widehat{A}=50^{\circ}$, $BH=2,3\,\text{cm}$. Điền Đ vào các câu trả lời đúng, S vào các câu trả lời sai.
	\choiceTF
	{\True $\widehat{C}=65^{\circ}$}
	{$AH\approx3,32\,\text{cm}$}
	{\True $BC\approx2,54\,\text{cm}$}
	{ Chu vi tam giác $ABC$ làm tròn đến hàng phần trăm là $8,52\,\text{cm}$}
	\loigiai{
	\begin{itemchoice}
	\itemch Ta có tam giác $ABC$ cân tại $A$ suy ra $\widehat{B}=\widehat{C}$ nên $\widehat{C}=\dfrac{1}{2}(180^{\circ}-\widehat{A})=\dfrac{1}{2}(180^{\circ}-50^{\circ})=65^{\circ}.$ Vậy mệnh đề đúng.
	\itemch Xét tam giác $AHB$ vuông tại $H$ ta có $AH=BH \cdot \cot\widehat{A}=2,3 \cdot \cot 50^{\circ}\approx2,3 \cdot 0,8391 \approx 1,93\,\text{cm}.$ Vậy mệnh đề ($AH\approx3,32\,\text{cm}$) sai.
	\itemch Xét $\triangle BHC$ vuông tại $H$ ta có $BC=\dfrac{BH}{\sin\widehat{C}}=\dfrac{2,3}{\sin 65^{\circ}}\approx\dfrac{2,3}{0,9063}\approx2,537\,\text{cm} \approx 2,54\,\text{cm}.$ Vậy mệnh đề đúng.
	\itemch Ta có $AC = AB$. Trong $\triangle ABH$ vuông tại $H$, $AB = \dfrac{BH}{\sin A} = \dfrac{2,3}{\sin 50^\circ} \approx \dfrac{2,3}{0,766} \approx 3,003\,\text{cm}$.
	Chu vi tam giác $ABC$ bằng $AB+BC+AC \approx 3,003+2,537+3,002 \approx 8,542\,\text{cm} \approx 8,54\,\text{cm}$.
	\end{itemchoice}
	}
\end{ex}
%%==========Câu 130
\begin{ex}
	\immini[thm]
	{
	Cho hình vẽ. Biết rằng các khoảng cách từ một điểm $C$ đến $A$ và đến $B$ là $CA=90\,\text{m}$, $CB=150\,\text{m}$ và $\widehat{ACB}=120^{\circ}$. Kẻ đường cao $AH$ từ $A$ xuống đường thẳng $BC$ ($H$ nằm ngoài đoạn $BC$ sao cho $C$ nằm giữa $B$ và $H$). Điền Đ vào các câu trả lời đúng, S vào các câu trả lời sai.
	\choiceTF
	{$AH=45\,\text{m}$}
	{$CH=45\sqrt{3}\,\text{m}$}
	{\True $BH=195\,\text{m}$}
	{\True $AB=210\,\text{m}$}
	}
	{
	% TODO: \usepackage{graphicx} required
	\includegraphics[scale=0.6]{C:/texstudio-man/Anh/{D3EBFA71-98B9-4B3A-A1D4-0641546BE1D4}}
	}
	\loigiai{
	\begin{itemchoice}
	\itemch Ta có: $\widehat{ACH}=180^{\circ}-\widehat{ACB}=180^{\circ}-120^{\circ}=60^{\circ}$. Xét tam giác $AHC$ vuông tại $H$ ta có: $AH=AC \cdot \sin\widehat{ACH} = 90 \cdot \sin 60^{\circ}=90 \cdot \dfrac{\sqrt{3}}{2}=45\sqrt{3}\,\text{m}.$ Vậy mệnh đề ($AH=45\,\text{m}$) sai.
	\itemch Xét tam giác $AHC$ vuông tại $H$ ta có: $CH=AC \cdot \cos\widehat{ACH} = 90 \cdot \cos 60^{\circ}=90 \cdot \dfrac{1}{2}=45\,\text{m}.$ Vậy mệnh đề ($CH=45\sqrt{3}\,\text{m}$) sai.
	\itemch Ta có $BH=BC+CH=150+45=195\,\text{m}.$ Vậy mệnh đề đúng.
	\itemch Theo Pythagore trong $\triangle ABH$ vuông tại $H$: $AB=\sqrt{BH^{2}+AH^{2}}=\sqrt{195^{2}+(45\sqrt{3})^{2}}=\sqrt{38025+6075}=\sqrt{44100}=210\,\text{m}.$ Vậy mệnh đề đúng.
	\end{itemchoice}
	}
\end{ex}
%%==========Câu 131
\begin{ex}
	\immini[thm]
	{
	Hai máy bay cùng cất cánh từ một sân bay $O$ nhưng bay theo hai hướng khác nhau. Một chiếc ($A$) di chuyển với tốc độ $450\,\text{km/h}$ theo hướng tây và chiếc còn lại ($B$) di chuyển theo hướng lệch so với hướng bắc $25^{\circ}$ về phía tây với tốc độ $630\,\text{km/h}$. Giả sử chúng đang ở cùng độ cao, sau $90$ phút ($1,5$ giờ) hãy Điền Đ vào các câu trả lời đúng, S vào các câu trả lời sai.
	\choiceTF
	{Góc giữa hai máy bay là $90^{\circ}$}
	{\True Máy bay thứ nhất đi được quãng đường là $675\,\text{km}$}
	{\True Máy bay thứ hai đi được quãng đường là $945\,\text{km}$}
	{\True Khoảng cách giữa hai máy bay gần bằng $900\,\text{km}$}
	}
	{
	% TODO: \usepackage{graphicx} required
	\includegraphics[scale=0.45]{C:/texstudio-man/Anh/{52C3A525-425D-4EFC-939D-50439CB3E59B}}
	}
	\loigiai{
	\begin{itemchoice}
	\itemch Hướng tây vuông góc với hướng bắc. Máy bay $B$ bay theo hướng lệch $25^{\circ}$ về phía tây so với hướng bắc. Vậy góc giữa hướng bay của $A$ (tây) và hướng bay của $B$ là $90^\circ - 25^\circ = 65^\circ$. Vậy mệnh đề (góc giữa hai máy bay là $90^{\circ}$) sai.
	\itemch Sau $90\,\text{phút} = 1,5$ giờ: Máy bay thứ nhất đi được quãng đường $OA$ là: $450 \cdot 1,5=675\,\text{km}.$ Vậy mệnh đề đúng.
	\itemch Máy bay thứ hai đi được quãng đường $OB$ là: $630 \cdot 1,5=945\,\text{km}.$ Vậy mệnh đề đúng.
	\itemch Kẻ $B H \perp O A$\\
	$\triangle O B H$ vuông tại $H$ có $O H=O B \cdot \cos 65^{\circ}=945 \cdot \cos 65^{\circ}$ và $B H=O B \cdot \sin 65^{\circ}=945 \cdot \sin 65^{\circ}$\\
	$	A H=O A-O H=675-945 \cdot \cos 65^{\circ}	$\\
	$\triangle A B H$ vuông tại $H$, theo định lí Pytago ta có $A B^2=A H^2+B H^2$\\
	Suy ra $A B=\sqrt{\left(945 \cdot \sin 65^{\circ}\right)^2+\left(675-945 \cdot \cos 65^{\circ}\right)^2} \approx 900 \mathrm{~km}$. 
	\end{itemchoice}
	}
\end{ex}
%%==========Câu 132
\begin{ex}
\immini[thm]
{
	Hải đăng Kê Gà có độ cao $HK=66\,\text{m}$. Tại hai điểm $A$ và $B$ trên đất liền người ta nhìn thấy ngọn đèn $H$ với góc nâng lần lượt là $\widehat{HAK}=30^{\circ}$ và $\widehat{HBK}=40^{\circ}$ ($K$ là chân hải đăng, $A$ và $B$ cùng phía so với $K$, $A$ xa hơn $B$). Hãy điền Đ vào câu trả lời đúng và S vào câu trả lời sai (các kết quả làm tròn đến hàng đơn vị):
	\choiceTF
	{Khoảng cách từ điểm $A$ đến chân ngọn hải đăng là $132\,\text{m}$}
	{\True Khoảng cách từ điểm $B$ đến chân ngọn hải đăng là $79\,\text{m}$}
	{\True Khoảng cách từ điểm $A$ đến ngọn đèn là $132\,\text{m}$}
	{Khoảng cách giữa hai người là $39\,\text{m}$}
}
{
	% TODO: \usepackage{graphicx} required
	\includegraphics[scale=0.45]{C:/texstudio-man/Anh/{3C243C8A-7FC9-4828-934B-DDFFF4B93C4C}}
}
	\loigiai{
	\begin{itemchoice}
	\itemch Khoảng cách từ điểm $A$ đến chân ngọn hải đăng là $AK$. Xét tam giác $AHK$ vuông tại $K$ ta có $AK=HK \cdot \cot A=66 \cdot \cot 30^{\circ} = 66\sqrt{3} \approx 66 \cdot 1,73205 \approx 114,315 \approx 114\,\text{m}.$ Vậy mệnh đề ($AK=132\,\text{m}$) sai.
	\itemch Khoảng cách từ điểm $B$ đến chân ngọn hải đăng là $BK$. Xét tam giác $BHK$ vuông tại $K$ ta có $BK=HK \cdot \cot B = 66 \cdot \cot 40^{\circ} \approx 66 \cdot 1,19175 \approx 78,655 \approx 79\,\text{m}.$ Vậy mệnh đề đúng.
	\itemch Khoảng cách từ điểm $A$ đến ngọn đèn là $AH$. Xét tam giác $AHK$ vuông tại $K$ ta có $AH=\dfrac{HK}{\sin\widehat{A}}=\dfrac{66}{\sin 30^{\circ}}=\dfrac{66}{1/2}=132\,\text{m}.$ Vậy mệnh đề đúng.
	\itemch Khoảng cách giữa hai người là $AB = AK-BK \approx 114,315 - 78,655 \approx 35,66 \approx 36\,\text{m}.$ Vậy mệnh đề ($AB=39\,\text{m}$) sai. (Lời giải PDF: $114-79=35\,\text{m}$).
	\end{itemchoice}
	}
\end{ex}
%%==========Câu 133
\begin{ex}
	Cho tam giác $ABC$ cân tại $A$, đường cao $AH$ có $\widehat{BAC}=120^{\circ}$, $BC=12\,\text{cm}$. Điền (Đ) cho câu trả lời đúng, (S) cho câu trả lời sai.
	\choiceTF
	{ $\widehat{B}=\widehat{C}=60^{\circ}$}
	{\True $BH=CH=6\,\text{cm}$}
	{ $AH=4\sqrt{3}\,\text{cm}$}
	{$AB=4\sqrt{3}\,\text{cm}$}
	\loigiai{
	\begin{itemchoice}
	\itemch Ta có: $\Delta ABC$ cân tại $A$, nên ta có $\widehat{B}+\widehat{C}=180^{\circ} - 120^{\circ} = 60^{\circ}\Rightarrow\widehat{B}=\widehat{C}=30^{\circ}$. \\ Vậy mệnh đề sai.
	\itemch Lại có $H$ là trung điểm của $BC \Rightarrow BH=CH=\dfrac{BC}{2} = \dfrac{12}{2}=6\,\text{cm}$. \\ Vậy mệnh đề đúng.
	\itemch Xét $\Delta ABH$ vuông tại $H$, $\tan \widehat{ABH}=\dfrac{AH}{BH}$ hay $AH=BH \cdot \tan\widehat{ABH}$. $AH=6 \cdot \tan 30^{\circ}=6\cdot\dfrac{\sqrt{3}}{3}=2\sqrt{3}\,\text{cm}$. \\ Vậy mệnh đề sai.
	\itemch Xét $\Delta ABH$ vuông tại $H$, $\sin \widehat{ABH}=\dfrac{AH}{AB}$. \\
	Suy ra $AB=\dfrac{AH}{\sin \widehat{ABH}}=\dfrac{2\sqrt{3}}{\sin 60^\circ}=4$
	\end{itemchoice}
	}
\end{ex}
%%==========Câu 134
\begin{ex}
\immini[thm]
{
	Cho hình vẽ sau, biết $\widehat{B}=60^{\circ}$, $\widehat{C}=30^{\circ}$, $AB=4\,\text{cm}$. Điền (Đ) cho câu trả lời đúng, (S) cho câu trả lời sai.
	\choiceTF
	{ $AH=2\,\text{cm}$}
	{\True $BH=2\,\text{cm}$}
	{\True $AC=4\sqrt{3}\,\text{cm}$}
	{ $BC=4\,\text{cm}$}
}
{
	% TODO: \usepackage{graphicx} required
	\includegraphics[scale=0.5]{C:/texstudio-man/Anh/{995B2A78-368C-435A-B052-60D9CD197FB7}}
	
}
	\loigiai{
	\begin{itemchoice}
	\itemch Xét tam giác $ABH$ vuông tại $H$, ta có: $AH=AB \cdot \sin B=4 \cdot \sin 60^{\circ}=4 \cdot \dfrac{\sqrt{3}}{2}=2\sqrt{3}\,\text{cm}$. \\ Vậy mệnh đề sai.
	\itemch Xét tam giác $ABH$ vuông tại $H$, ta có: $BH=AB \cdot \cos B=4 \cdot \cos 60^{\circ}=4 \cdot \dfrac{1}{2}=2\,\text{cm}$. \\ Vậy mệnh đề đúng.
	\itemch Xét tam giác $ACH$ vuông tại $H$, ta có: $AH=2\sqrt{3}\,\text{cm}$ (từ câu a). $\sin C=\dfrac{AH}{AC}$ hay $AC=\dfrac{AH}{\sin C}=\dfrac{2\sqrt{3}}{\sin 30^{\circ}}=\dfrac{2\sqrt{3}}{1/2}=4\sqrt{3}\,\text{cm}$. \\ Vậy mệnh đề đúng.
	\itemch Xét tam giác $ACH$ vuông tại $H$, ta có: $HC=AC \cdot \cos C=4\sqrt{3} \cdot \cos 30^{\circ}=4\sqrt{3} \cdot \dfrac{\sqrt{3}}{2}=6\,\text{cm}$. $BC=BH+HC=2+6=8\,\text{cm}$. \\ Vậy mệnh đề sai.
	\end{itemchoice}
	}
\end{ex}
%%==========Câu 135
\begin{ex}
	Tam giác $ABC$ có $\widehat{B}=60^{\circ}$, $\widehat{C}=35^{\circ}$ và $AC=35\,\text{cm}.$ Kẻ đường cao $AH$. (Các kết quả làm tròn đến hàng phần mười). Điền (Đ) cho câu trả lời đúng, (S) cho câu trả lời sai.
	\choiceTF
	{\True $AH \approx 20,1\,\text{cm}$}
	{ $HC \approx 24,5\,\text{cm}$}
	{\True $BH \approx 11,6\,\text{cm}$}
	{ $BC=BH+HC \approx 36,1\,\text{cm}$}
	\loigiai{
	\begin{itemchoice}
	\itemch Xét tam giác $ACH$ vuông tại $H$, ta có: $AH=AC \cdot \sin C=35 \cdot \sin 35^{\circ} \approx 35 \cdot 0,5736 \approx 20,076 \approx 20,1\,\text{cm}$. \\ Vậy mệnh đề đúng.
	\itemch Xét tam giác $ACH$ vuông tại $H$, ta có: $HC=AC \cdot \cos C=35 \cdot \cos 35^{\circ} \approx 35 \cdot 0,8192 \approx 28,672 \approx 28,7\,\text{cm}$. \\ Vậy mệnh đề sai.
	\itemch Xét tam giác $ABH$ vuông tại $H$, ta có: $\tan B=\dfrac{AH}{BH}$ hay $BH=\dfrac{AH}{\tan B} \approx \dfrac{20,1}{\tan 60^{\circ}} \approx \dfrac{20,1}{1,732} \approx 11,605 \approx 11,6\,\text{cm}$. \\ Vậy mệnh đề đúng.
	\itemch Có: $BC=BH+HC \approx 11,6 + 28,7 \approx 40,3\,\text{cm}$. \\ Vậy mệnh đề sai.
	\end{itemchoice}
	}
\end{ex}
%%==========Câu 136
\begin{ex}
	Cho hình thoi $ABCD$ có hai đường chéo $AC=2\sqrt{3}\,\text{cm}$ và $BD=2\,\text{cm}$. Điền (Đ) cho câu trả lời đúng, (S) cho câu trả lời sai.
	\choiceTF
	{\True $BO=1\,\text{cm}$}
	{ $OC=2\sqrt{3}\,\text{cm}$}
	{\True $\widehat{BCO}=30^{\circ}$}
	{\True $\widehat{BCD}=60^{\circ}$}
	\loigiai{
	\begin{itemchoice}
	\itemch Gọi giao điểm của hai đường chéo hình thoi là điểm $O$. Vì $ABCD$ là hình thoi nên $OA=OC=\dfrac{AC}{2}=\dfrac{2\sqrt{3}}{2}=\sqrt{3}\,\text{cm}$ và $OB=OD=\dfrac{BD}{2}=\dfrac{2}{2}=1\,\text{cm}$. Do đó $BO=1\,\text{cm}$. \\ Vậy mệnh đề đúng.
	\itemch $OA=OC=\sqrt{3}\,\text{cm}$. \\ Vậy mệnh đề sai.
	\itemch Có $AC \perp BD$ tại $O$ (tính chất hình thoi), nên xét $\triangle BCO$ vuông tại $O$, ta có: $\tan \widehat{BCO}=\dfrac{BO}{OC}=\dfrac{1}{\sqrt{3}}$ hay $\widehat{BCO}=30^{\circ}$. \\ Vậy mệnh đề đúng.
	\itemch Suy ra $\widehat{BCD}=2 \cdot \widehat{BCO}=2 \cdot 30^{\circ}=60^{\circ}$. \\ Vậy mệnh đề đúng.
	\end{itemchoice}
	}
\end{ex}
%%==========Câu 137
\begin{ex}
	Cho tam giác $ABC$ có $BC=12\,\text{cm},$ $\widehat{B}=40^{\circ}$, $\widehat{C}=30^{\circ}$. Gọi $N$ là chân đường vuông góc hạ từ $A$ xuống cạnh $BC$. (Các kết quả làm tròn đến hàng phần mười). Điền (Đ) cho câu trả lời đúng, (S) cho câu trả lời sai.
	\choiceTF
	{\True $AN \approx 4,1\,\text{cm}$}
	{\True $AC \approx 8,2\,\text{cm}$}
	{ $NC \approx 4,1\,\text{cm}$}
	{\True $S_{ABC} \approx 24,6\,\text{cm}$}
	\loigiai{% TODO: \usepackage{graphicx} required
		\begin{center}
			\includegraphics[scale=0.6]{C:/texstudio-man/Anh/{EF65A6B4-6C67-4BF8-8C90-3E60D7E37895}}
		\end{center}
	Đặt $BN=x\,\text{cm}$ ($0<x<12$) có $NC=12-x\,\text{cm}$.
	\begin{itemchoice}
	\itemch Xét tam giác $ABN$ vuông tại $N$, ta có: $AN=BN \cdot \tan B=x \cdot \tan 40^{\circ}$.\\
	 Xét tam giác $ACN$ vuông tại $N$, ta có: $AN=CN \cdot \tan C=(12-x) \cdot \tan 30^{\circ}$.\\ Nên $x \cdot \tan 40^{\circ}=(12-x) \cdot \tan 30^{\circ}$\\
	 $ \Rightarrow x \cdot \tan 40^{\circ}=12 \cdot \tan 30^{\circ}-x \cdot \tan 30^{\circ}$\\ $\Rightarrow x \cdot (\tan 40^{\circ}+\tan 30^{\circ})=12 \cdot \tan 30^{\circ}$\\
	 $\Rightarrow x=\dfrac{12 \cdot \tan 30^{\circ}}{\tan 40^{\circ}+\tan 30^{\circ}} \approx \dfrac{12 \cdot 0,577}{0,839+0,577} \approx \dfrac{6,924}{1,416} \approx 4,89\,\text{cm}$.\\
	 Khi đó, $AN=BN \cdot \tan B \approx 4,89 \cdot \tan 40^{\circ} \approx 4,89 \cdot 0,839 \approx 4,103 \approx 4,1\,\text{cm}$. \\ Vậy mệnh đề đúng.
	\itemch Xét tam giác $ACN$ vuông tại $N$, ta có:\\
	$\sin C=\dfrac{AN}{AC}$ hay $AC=\dfrac{AN}{\sin C} \approx \dfrac{4,1}{\sin 30^{\circ}} \approx \dfrac{4,1}{0,5} \approx 8,2\,\text{cm}$. \\ Vậy mệnh đề đúng.
	\itemch Xét tam giác $ACN$ vuông tại $N$, ta có: \\
	$NC=AC \cdot \cos C \approx 8,2 \cdot \cos 30^{\circ} \approx 8,2 \cdot \dfrac{\sqrt{3}}{2} \approx 8,2 \cdot 0,866 \approx 7,1012 \approx 7,1\,\text{cm}$. \\ Vậy mệnh đề sai.
	\itemch $S_{ABC}=\dfrac{1}{2} \cdot AN \cdot BC \approx \dfrac{1}{2} \cdot 4,1 \cdot 12 \approx 24,6\,\text{cm}$. \\ Vậy mệnh đề đúng.
	\end{itemchoice}
	}
\end{ex}
%%==========Câu 138
\begin{ex}
	Cho tam giác $ABC$ cân tại $A$, có $\widehat{B}=65^{\circ}$ đường cao $CH=3,6\,\text{cm}.$ Điền (Đ) cho câu trả lời đúng, (S) cho câu trả lời sai. (Các kết quả làm tròn đến hàng phần trăm):
	\choiceTF
	{\True $\widehat{A}=50^{\circ}$}
	{\True $\widehat{C}=65^{\circ}$}
	{ $AB=AC \approx 5,60\,\text{cm}$}
	{ $BC \approx 8,52\,\text{cm}$}
	\loigiai{% TODO: \usepackage{graphicx} required
		\begin{center}
			\includegraphics[scale=0.6]{C:/texstudio-man/Anh/{11C1BFD6-1878-49A2-8C1E-A368F2B107E5}}
		\end{center}
	\begin{itemchoice}
	\itemch Ta có: $\widehat{A}+\widehat{B}+\widehat{C}=180^{\circ}$ (định lý tổng ba góc trong một tam giác). Suy ra $\widehat{A}=180^{\circ}-(\widehat{B}+\widehat{C})=180^{\circ}-(65^{\circ}+65^{\circ})=50^{\circ}$. \\ Vậy mệnh đề đúng.
	\itemch Vì tam giác $ABC$ cân tại $A$ nên $\widehat{B}=\widehat{C}=65^{\circ}$. \\ Vậy mệnh đề đúng.
	\itemch Xét tam giác $ACH$ vuông tại $H$, ta có: $\sin A=\dfrac{CH}{AC}$ hay $AC=\dfrac{CH}{\sin A}=\dfrac{3,6}{\sin 50^{\circ}} \approx \dfrac{3,6}{0,7660} \approx 4,70\,\text{cm}$. Vì $ABC$ cân tại $A$ nên $AB=AC \approx 4,70\,\text{cm}$. \\ Vậy mệnh đề sai.
	\itemch Xét tam giác $BCH$ vuông tại $H$, ta có: $\sin B=\dfrac{CH}{BC}$ hay $BC=\dfrac{CH}{\sin B}=\dfrac{3,6}{\sin 65^{\circ}} \approx \dfrac{3,6}{0,9063} \approx 3,97\,\text{cm}$. \\ Vậy mệnh đề sai.
	\end{itemchoice}
	}
\end{ex}
%%==========Câu 139
\begin{ex}
	Cho tam giác $ABC$ cân tại $A$, đường cao $AH$. Biết $AH=BC$. Điền (Đ) cho câu trả lời đúng, (S) cho câu trả lời sai. (Các kết quả về số đo góc làm tròn đến phút)
	\choiceTF
	{\True $\tan B=2$}
	{ $\tan C=\dfrac{1}{2}$}
	{ $\widehat{B}=\widehat{C} \approx 65^{\circ}57'$}
	{\True $\widehat{A} \approx 53^{\circ}8'$}
	\loigiai{% TODO: \usepackage{graphicx} required
		\begin{center}
			\includegraphics[scale=0.6]{C:/texstudio-man/Anh/{E4B55124-EAC2-4471-ACEF-9015876E5BCF}}
		\end{center}
	Giả sử $AH=BC=a$.
	\begin{itemchoice}
	\itemch Vì $ABC$ cân tại $A$ nên $AH$ là đường cao đồng thời là đường trung tuyến. Do đó, $H$ là trung điểm của $BC$, hay $HB=HC=\dfrac{BC}{2}=\dfrac{a}{2}$.\\
	Xét tam giác $ABH$ vuông tại $H$, ta có: \\
	$\tan B=\dfrac{AH}{BH}=\dfrac{a}{\frac{a}{2}}=2$. \\ Vậy mệnh đề đúng.
	\itemch Vì $ABC$ cân tại $A$ nên $\tan B=\tan C=2$. \\ Vậy mệnh đề sai.
	\itemch Từ $\tan B=2$, suy ra: $\widehat{B} \approx 63^{\circ}26'$. Vì $ABC$ cân tại $A$ nên $\widehat{B}=\widehat{C} \approx 63^{\circ}26'$. \\ Vậy mệnh đề sai.
	\itemch Ta có $\widehat{A}+\widehat{B}+\widehat{C}=180^{\circ}$ (định lý tổng ba góc trong một tam giác).\\
	Suy ra $\widehat{A}=180^{\circ}-(\widehat{B}+\widehat{C}) \approx 180^{\circ}-(63^{\circ}26'+63^{\circ}26') = 180^{\circ} - 126^{\circ}52' \approx 53^{\circ}8'$. \\ Vậy mệnh đề đúng.
	\end{itemchoice}
	}
\end{ex}
%%==========Câu 140
\begin{ex}
	Cho hình thang $ABCD$ vuông tại $A$ và $D$; $\widehat{C}=50^{\circ}$. Biết $AB=2\,\text{cm}$, $AD=1,2\,\text{cm}.$ Kẻ $BE \perp DC$, $E \in CD$. Điền (Đ) cho câu trả lời đúng, (S) cho câu trả lời sai.
	\choiceTF
	{ $ED=1,2\,\text{cm}$}
	{\True $EC=\dfrac{1,2}{\tan 50^{\circ}}$}
	{\True $DC=2+\dfrac{1,2}{\tan 50^{\circ}}$}
	{ $S_{ABCD} \approx 4\,\text{cm}^2$}
	\loigiai{% TODO: \usepackage{graphicx} required
		\begin{center}
			\includegraphics[scale=0.6]{C:/texstudio-man/Anh/{3771ECC8-E0DC-4229-8779-9EE095B84A24}}
		\end{center}
		
	\begin{itemchoice}
	\itemch Xét tứ giác $ABED$ có $\widehat{A}=\widehat{D}=\widehat{E}=90^{\circ}$. Suy ra $ABED$ là hình chữ nhật.\\
	Do đó, $AB=ED=2\,\text{cm}$ và $AD=BE=1,2\,\text{cm}$. \\ Vậy mệnh đề sai.
	\itemch Xét tam giác $BCE$ vuông tại $E$, ta có: $\tan C=\dfrac{BE}{CE}$ hay $CE=\dfrac{BE}{\tan C}=\dfrac{1,2}{\tan 50^{\circ}}$. \\ Vậy mệnh đề đúng.
	\itemch Ta có: $DC=DE+CE=2+\dfrac{1,2}{\tan 50^{\circ}}$. \\ Vậy mệnh đề đúng.
	\itemch $S_{ABCD}=\dfrac{(AB+CD) \cdot AD}{2}=\dfrac{(2+2+\frac{1,2}{\tan 50^{\circ}}) \cdot 1,2}{2} = \approx \dfrac{6,0072}{2} \approx 3\,\text{cm}^2$. \\ Vậy mệnh đề sai.
	\end{itemchoice}
	}
\end{ex}
%%==========Câu 141
\begin{ex}
	\immini[thm]
	{
		Cho hình vẽ. Điền (Đ) cho câu trả lời đúng, (S) cho câu trả lời sai.
	\choiceTF
	{\True $b=a \cdot \sin B$}
	{ $a=b \cdot \sin B$}
	{ $c=\dfrac{a}{\sin C}$}
	{\True $a=\dfrac{c}{\sin C}$}
	}
	{
		% TODO: \usepackage{graphicx} required
		\includegraphics[scale=0.6]{C:/texstudio-man/Anh/{25E7F14A-CB3F-4913-978F-6014E41DBAAD}}
		
	}
	\loigiai{
	\begin{itemchoice}
	\itemch Trong tam giác vuông tại $A$, $\sin B = \dfrac{b}{a}$, suy ra $b=a \cdot \sin B$. \\ Vậy mệnh đề đúng.
	\itemch Từ $b=a \cdot \sin B$, thì $a = \dfrac{b}{\sin B}$. \\ Vậy mệnh đề sai.
	\itemch Trong tam giác vuông tại $A$, $\sin C = \dfrac{c}{a}$, suy ra $c=a \cdot \sin C$ hoặc $a = \dfrac{c}{\sin C}$. Do đó $c=\dfrac{a}{\sin C}$ là sai (trừ khi $\sin^2 C=1$). \\ Vậy mệnh đề sai.
	\itemch Trong tam giác vuông tại $A$, $\sin C = \dfrac{c}{a}$, suy ra $a=\dfrac{c}{\sin C}$. \\ Vậy mệnh đề đúng.
	\end{itemchoice}
	}
\end{ex}
%%==========Câu 142
\begin{ex}
	\immini[thm]
	{
	Cho hình vẽ. Điền (Đ) cho câu trả lời đúng, (S) cho câu trả lời sai.
	\choiceTF
	{\True $c=a \cdot \cos B$}
	{ $a=c \cdot \cos B$}
	{\True $b=a \cdot \cos C$}
	{ $a=b \cdot \cos C$}
	}
	{
		% TODO: \usepackage{graphicx} required
		\includegraphics[scale=0.6]{C:/texstudio-man/Anh/{F672B879-3FBB-4335-97AD-C933C4175B69}}
		
	}
	\loigiai{
	\begin{itemchoice}
	\itemch Trong tam giác vuông tại $A$, $\cos B = \dfrac{c}{a}$, suy ra $c=a \cdot \cos B$.\\ Vậy mệnh đề đúng.
	\itemch Từ $c=a \cdot \cos B$, thì $a = \dfrac{c}{\cos B}$. \\ Vậy mệnh đề sai.
	\itemch Trong tam giác vuông tại $A$, $\cos C = \dfrac{b}{a}$, suy ra $b=a \cdot \cos C$. \\ Vậy mệnh đề đúng.
	\itemch Từ $b=a \cdot \cos C$, thì $a = \dfrac{b}{\cos C}$. \\ Vậy mệnh đề sai.
	\end{itemchoice}
	}
\end{ex}
%%==========Câu 143
\begin{ex}
\immini[thm]
{
	Cho hình vẽ. Điền (Đ) cho câu trả lời đúng, (S) cho câu trả lời sai.
	\choiceTF
	{\True $b=c \cdot \tan B$}
	{\True $c=\dfrac{b}{\tan B}$}
	{$c=\dfrac{b}{\tan C}$}
	{$b=c \cdot \tan C$}
}
{
	% TODO: \usepackage{graphicx} required
	\includegraphics[scale=0.6]{C:/texstudio-man/Anh/{A53D5EDE-988E-4512-A083-3356249BDE0B}}
	
}
	\loigiai{
	Sử dụng tỉ số lượng giác trong tam giác vuông
	}
\end{ex}
%%==========Câu 144
\begin{ex}
	Cho tam giác $ABC$ nhọn, đường cao $AH$. Điền (Đ) cho câu trả lời đúng, (S) cho câu trả lời sai.
	\choiceTF
	{\True $AH=AB \cdot \sin B$}
	{\True $AH=AC \cdot \sin C$}
	{ $BC=AB \cdot \cos B+AH \cdot \cos C$}
	{ $BC=AB \cdot \cos B-AC \cdot \cos C$}
	\loigiai{
	\begin{itemchoice}
	\itemch Xét tam giác $ABH$ vuông tại $H$, ta có: $AH=AB \cdot \sin B$. \\ Vậy mệnh đề đúng.
	\itemch Xét tam giác $ACH$ vuông tại $H$, ta có: $AH=AC \cdot \sin C$. \\ Vậy mệnh đề đúng.
	\itemch Xét tam giác $ABH$ vuông tại $H$, $BH=AB \cdot \cos B$. Xét tam giác $ACH$ vuông tại $H$, $HC=AC \cdot \cos C$. Có $BC = BH+HC = AB \cdot \cos B+AC \cdot \cos C$. Mệnh đề đã cho là $BC=AB \cdot \cos B+AH \cdot \cos C$. \\ Vậy mệnh đề sai.
	\itemch Có $BC = BH+HC = AB \cdot \cos B+AC \cdot \cos C$. \\ Vậy mệnh đề sai.
	\end{itemchoice}
	}
\end{ex}
%%==========Câu 145
\begin{ex}
	Cho tam giác $ABC$ có trực tâm $H$ là trung điểm của đường cao $AD$. Điền (Đ) cho câu trả lời đúng, (S) cho câu trả lời sai.
	\choiceTF
	{\True $\cos B \cdot \cos C=\dfrac{BD}{AB}\cdot\dfrac{DC}{AC}$}
	{\True $DB \cdot DC=DA \cdot DH$}
	{\True $\cos A=\cos B \cdot \cos C$}
	{ $\cos A=\cos A \cdot \cos C$}
	\loigiai{
		\begin{center}
			\includegraphics[scale=0.55]{C:/texstudio-man/Anh/{8045531E-F061-48F5-AA90-CAADFC0BAE05}}
		\end{center}
	\begin{itemchoice}
	\itemch Xét tam giác $ABD$ vuông tại $D$, ta có: $\cos \widehat{ABD}=\dfrac{BD}{AB}$.\\
	Xét tam giác $ACD$ vuông tại $D$, ta có: $\cos C=\dfrac{CD}{AC}$.\\
	Suy ra: $\cos \widehat{ABD} \cdot \cos C=\dfrac{BD}{AB} \cdot \dfrac{DC}{AC}$. (Trong $\triangle ABC$, $\widehat{ABD}$ là $\widehat{B}$). \\ Vậy mệnh đề đúng.
	\itemch Chứng minh được: $\Delta DBH \sim \Delta DAC$ (g.g).\\
	Suy ra: $\dfrac{DB}{DA}=\dfrac{DH}{DC}$ hay $DB \cdot DC=DA \cdot DH$. \\ 
	Vậy mệnh đề đúng.
	\itemch Xét tam giác $ABE$ vuông tại $E$, ta có: $\cos \widehat{BAE}=\dfrac{AE}{AB}$. (1)\\
	Chứng minh được: $\Delta DBH \sim \Delta DAC$ (g.g) $\Rightarrow DB \cdot DC=DA \cdot DH$.\\
	Mà $AD=2DH$ (do $H$ là trung điểm $AD$).\\
	 Nên $DB \cdot DC=2DH \cdot DH = 2DH^2$. (2)\\
	Chứng minh được: $\Delta AEH \sim \Delta ADC$ (g.g). Suy ra: $\dfrac{AE}{AD}=\dfrac{AH}{AC}$ hay $AE \cdot AC=AD \cdot AH = AD \cdot DH$. (3)\\
	Từ (2) và (3), ta có: $DB \cdot DC=AE \cdot AC$. (4)\\
	Thay (4) vào biểu thức của mệnh đề (a): $\cos B \cdot \cos C = \dfrac{DB \cdot DC}{AB \cdot AC} = \dfrac{AE \cdot AC}{AB \cdot AC} = \dfrac{AE}{AB}$.\\
	Từ (1), $\dfrac{AE}{AB} = \cos \widehat{BAE} = \cos A$.\\
	Vậy $\cos A=\cos B \cdot \cos C$. \\ Vậy mệnh đề đúng.
	\itemch Từ chứng minh trên, $\cos A=\cos B \cdot \cos C$. \\ Vậy mệnh đề sai.
	\end{itemchoice}
	}
\end{ex}
%%==========Câu 146
\begin{ex}
	Cho tam giác đều $ABC$ có $AB=a$. Kẻ đường cao $AH$. Điền (Đ) cho câu trả lời đúng, (S) cho câu trả lời sai.
	\choiceTF
	{\True $\widehat{B}=60^{\circ}$}
	{\True $AH=\dfrac{a\sqrt{3}}{2}$}
	{ $BH=a$}
	{\True $S_{ABC}=\dfrac{a^{2}\sqrt{3}}{4}$}
	\loigiai{
	\begin{itemchoice}
	\itemch Ta có $\Delta ABC$ đều, có $AH$ là đường cao nên $\widehat{B}=60^{\circ}$. \\ Vậy mệnh đề đúng.
	\itemch Ta có $\Delta ABH$ vuông tại $H$ nên $AH=AB \cdot \sin B=a \cdot \sin 60^{\circ}=a \cdot \dfrac{\sqrt{3}}{2}$.\\ Vậy mệnh đề đúng.
	\itemch Ta có $\Delta ABH$ vuông tại $H$ nên $BH=AB \cdot \cos B=a \cdot \cos 60^{\circ}=a \cdot \dfrac{1}{2}$.\\ Vậy mệnh đề sai.
	\itemch $S_{ABC}=\dfrac{1}{2}AH \cdot BC = \dfrac{1}{2} \cdot \dfrac{a\sqrt{3}}{2} \cdot a = \dfrac{a^{2}\sqrt{3}}{4}$. \\ Vậy mệnh đề đúng.
	\end{itemchoice}
	}
\end{ex}
%%==========Câu 147
\begin{ex}
	Cho tam giác nhọn $ABC$. $AH$ là đường cao, $D, E$ lần lượt là hình chiếu của $H$ trên $AB, AC$. Điền (Đ) cho câu trả lời đúng, (S) cho câu trả lời sai.
	\choiceTF
	{\True $AH^{2}=AD \cdot AB$}
	{ $AH^{2}=AE \cdot AB$}
	{\True $AD \cdot AB=AE \cdot AC$}
	{ $\widehat{AED}=\widehat{ACB}$}
	\loigiai{
	\vspace*{.3cm}
		\begin{center}
			\includegraphics[scale=0.6]{C:/texstudio-man/Anh/{D1CB3BA3-D456-428F-BDB7-E063F8921C20}}
		\end{center}
	\begin{itemchoice}
	\itemch Xét $\Delta AHB$ vuông tại $H$ có $HD$ là đường cao, ta có: $AH^{2}=AD \cdot AB$. \\ Vậy mệnh đề đúng.
	\itemch Xét $\Delta AHC$ vuông tại $H$ có $HE$ là đường cao, ta có: $AH^{2}=AE \cdot AC$. \\ Vậy mệnh đề sai.
	\itemch Từ hai kết quả trên, $AD \cdot AB = AH^2$ và $AE \cdot AC = AH^2$. Do đó, $AD \cdot AB=AE \cdot AC$. \\ Vậy mệnh đề đúng.
	\itemch Xét $\Delta AED$ và $\Delta ABC$ có: $\widehat{EAD}$ chung. Từ $AD \cdot AB=AE \cdot AC \Rightarrow \dfrac{AE}{AB}=\dfrac{AD}{AC}$. Suy ra $\Delta AED \sim \Delta ABC$ (c.g.c). Nên $\widehat{AED}=\widehat{ABC}$. \\ Vậy mệnh đề sai.
	\end{itemchoice}
	}
\end{ex}
%%==========Câu 148
\begin{ex}
	Cho tam giác nhọn $ABC$, $BD$ và $CE$ là hai đường cao. Các điểm $N, M$ trên các đường thẳng $BD, CE$ sao cho $\widehat{AMB}=\widehat{ANC}=90^{\circ}$. Điền (Đ) cho câu trả lời đúng, (S) cho câu trả lời sai.
	\choiceTF
	{\True $\Delta ABD \sim \Delta ACE$}
	{ $\dfrac{AB}{AE}=\dfrac{AD}{AC}$}
	{\True $AM^{2}=AE \cdot AC$}
	{\True $\Delta AMN$ cân tại A}
	\loigiai{
		\begin{center}
			\includegraphics[scale=0.6]{C:/texstudio-man/Anh/{73DD0BD7-041B-4BCE-A872-78270ACAEE8B}}
		\end{center}
	\begin{itemchoice}
	\itemch Xét $\Delta ABD$ và $\Delta ACE$ có: $\widehat{BAD}$ chung và $\widehat{ADB}=\widehat{AEC}(=90^{\circ})$. Do đó $\Delta ABD \sim \Delta ACE$ (g.g)
	\itemch Từ $\Delta ABD \sim \Delta ACE$ (g.g), suy ra $\dfrac{AB}{AC}=\dfrac{AD}{AE}\Rightarrow AB.AE=AC.AD$\hfill (1).\\
	Mệnh đề là $\dfrac{AB}{AE}=\dfrac{AD}{AC}$. \\ Vậy mệnh đề sai.
	\itemch Xét $\triangle AMB$ vuông tại $M$ (gt), $ME$ là đường cao (do $CE$ là đường cao của $\triangle ABC$, $M \in CE$, $E \in AB$).\\
	Ta có: $AM^{2}=AE \cdot AB$.\hfill(2)\\
	Mà mệnh đề là $AM^{2}=AE \cdot AC$
	\\ Vậy mệnh đề sai.
	\itemch Tương tự, $AN^{2}=AD \cdot AC$.\hfill (3)\\
	Từ (1), (2), (3) suy ra $A M^2=A N^2$ hay $A M=A N$  .Vậy mệnh đề đúng.
	\end{itemchoice}
	}
\end{ex}
%%==========Câu 149
\begin{ex}
	Một cột đèn $AB$ cao $7,5\,\text{m}$ có bóng in trên mặt đất là $AC$ dài $4,5\,\text{m}$. Điền (Đ) cho câu trả lời đúng, (S) cho câu trả lời sai (số đo góc làm tròn đến đơn vị độ).
	\choiceTF
	{ $\sin C=\dfrac{3}{5}$}
	{ $\cos B=\dfrac{7,5}{3,5}$}
	{\True $\tan C=\dfrac{5}{3}$}
	{\True $\widehat{C} \approx 59^{\circ}$}
	\loigiai{Xét $\triangle ABC$ vuông tại $A$.
	\begin{itemchoice}
	\itemch $BC = \sqrt{AB^2+AC^2} = \sqrt{7,5^2+4,5^2} = \sqrt{56,25+20,25} = \sqrt{76,5} \approx 8,746\,\text{m}$. $\sin C = \dfrac{AB}{BC} = \dfrac{7,5}{\sqrt{76,5}} \approx \dfrac{7,5}{8,746} \approx 0,857 \neq \dfrac{3}{5}=0,6$. \\ Vậy mệnh đề sai.
	\itemch $\cos B = \dfrac{AB}{BC} = \dfrac{7,5}{\sqrt{76,5}} \approx 0,857$. Mệnh đề $\cos B=\dfrac{7,5}{3,5}$ là sai. \\ Vậy mệnh đề sai.
	\itemch $\tan C=\dfrac{AB}{AC}=\dfrac{7,5}{4,5}=\dfrac{75}{45}=\dfrac{5}{3}$. \\ Vậy mệnh đề đúng.
	\itemch Từ $\tan C=\dfrac{5}{3} \approx 1,6667$, suy ra: $\widehat{C} \approx \arctan(5/3) \approx 59,036^{\circ} \approx 59^{\circ}$. \\ Vậy mệnh đề đúng.
	\end{itemchoice}
	}
\end{ex}
%%==========Câu 150
\begin{ex}
	\immini[thm]
	{
	Một bức tường đang xây dở có dạng hình thang vuông $ABCD$, vuông góc ở $A$ và $D$, $AB=1\,\text{m}.$ $CD=5\,\text{m}.$ $AD=7\,\text{m}.$ Giả sử góc $\alpha$ là góc tạo bởi đường thẳng $BC$ và mặt đất $AD$, $I$ là giao điểm của hai đường thẳng $AD$ và $BC$. Qua $B$ kẻ đường thẳng song song với $AD$ (mặt đất) cắt $CD$ ở $E$. Điền (Đ) cho câu trả lời đúng, (S) cho câu trả lời sai. (Biết số đo góc làm tròn đến độ và độ dài đoạn thẳng làm tròn đến mét).
	\choiceTF
	{\True $BE=7\,\text{m}$}
	{ $EC=6\,\text{m}$}
	{ $\alpha \approx 35^{\circ}$}
	{\True $BC=8\,\text{m}$}
	}
	{
		% TODO: \usepackage{graphicx} required
		\includegraphics[scale=0.45]{C:/texstudio-man/Anh/{C3724FD8-EEEB-4AB9-BDE9-2AFD094EEA83}}
		
	}
	\loigiai{Ta có: $\widehat{EBC}=\widehat{DIC}=\alpha$ (hai góc đồng vị). Tứ giác $ABED$ có $BE \parallel AD$, $AB \parallel DE$, $\widehat{BAD}=90^{\circ}$ nên $ABED$ là hình chữ nhật.
	\begin{itemchoice}
	\itemch Do $ABED$ là hình chữ nhật, $BE=AD=7\,\text{m}$. \\ Vậy mệnh đề đúng.
	\itemch $AB=ED=1\,\text{m}$. Suy ra $EC=CD-ED=5-1=4\,\text{m}$. \\ Vậy mệnh đề sai.
	\itemch Xét $\triangle BCE$ vuông tại $E$, ta có: $\tan \alpha=\tan \widehat{EBC}=\dfrac{EC}{BE}=\dfrac{4}{7} \approx 0,5714$. Suy ra $\alpha \approx 30^{\circ}$. \\ Vậy mệnh đề sai.
	\itemch Xét $\triangle BCE$ vuông tại $E$, ta có: $BC = \sqrt{BE^2+EC^2} = \sqrt{7^2+4^2} = \sqrt{49+16} = \sqrt{65} \approx 8,06\,\text{m} \approx 8\,\text{m}$. Hoặc $BC=\dfrac{EC}{\sin \alpha} \approx \dfrac{4}{\sin 30^{\circ}} = \dfrac{4}{0,5} = 8\,\text{m}$. \\ Vậy mệnh đề đúng.
	\end{itemchoice}
	}
\end{ex}
%%==========Câu 151
\begin{ex}
\immini[thm]
{
	Để đo chiều cao của một cây thông đỉnh $O$, người ta lấy hai điểm $B$ và $C$ trên mặt đất với $BC=2\,\text{m}$. Góc nhìn đỉnh $O$ từ $B$ là $47^{\circ}$ từ $C$ là $38^{\circ}$ ( $A$ là chân cây, $C,B,A$ thẳng hàng). Điền (Đ) cho câu trả lời đúng, (S) cho câu trả lời sai. (Độ dài làm tròn đến hàng phần trăm).
	\choiceTF
	{ $OA=AB \cdot \sin 47^{\circ}$}
	{\True $OA=AC \cdot \tan 38^{\circ}$}
	{\True $AB \approx 5,37\,\text{m}$}
	{\True Chiều cao của cây thông xấp xỉ bằng $5,76\,\text{m}$}
}
{
	% TODO: \usepackage{graphicx} required
	\includegraphics[scale=0.5]{C:/texstudio-man/Anh/{9F9CCC39-204D-4EC6-8ED4-D3478F2C6571}}
}
	\loigiai{
	\begin{itemchoice}
	\itemch Xét tam giác $OAB$ vuông tại $A$ có $OA=AB \cdot \tan \widehat{OBA}=AB \cdot \tan 47^{\circ}$. \\ Vậy mệnh đề sai.
	\itemch Xét tam giác $OAC$ vuông tại $A$ có $OA=AC \cdot \tan \widehat{OCA}=AC \cdot \tan 38^{\circ}$. \\ Vậy mệnh đề đúng.
	\itemch Ta có $AC = AB+BC = AB+2$.\\
	Do đó: $AB \cdot \tan 47^{\circ}=AC \cdot \tan 38^{\circ} \Rightarrow AB \cdot \tan 47^{\circ}=(AB+2) \cdot \tan 38^{\circ} \Rightarrow AB \cdot \tan 47^{\circ}=AB \cdot \tan 38^{\circ} + 2 \cdot \tan 38^{\circ} \Rightarrow AB \cdot (\tan 47^{\circ}-\tan 38^{\circ})=2 \cdot \tan 38^{\circ} \Rightarrow AB=\dfrac{2 \cdot \tan 38^{\circ}}{\tan 47^{\circ}-\tan 38^{\circ}} \approx \dfrac{2 \cdot 0,7813}{1,0724-0,7813} \approx \dfrac{1,5626}{0,2911} \approx 5,3679\,\text{m} \approx 5,37\,\text{m}$. \\ Vậy mệnh đề đúng.
	\itemch Suy ra $OA = AB \cdot \tan 47^{\circ} \approx 5,37 \cdot \tan 47^{\circ} \approx 5,37 \cdot 1,0724 \approx 5,7585\,\text{m} \approx 5,76\,\text{m}$. Vậy chiều cao của cây thông xấp xỉ bằng $5,76\,\text{m}$. \\ Vậy mệnh đề đúng.
	\end{itemchoice}
	}
\end{ex}
%%==========Câu 152
\begin{ex}
\immini[thm]
{
	Cho hình vẽ sau ($CH=30\,\text{m}$, $\widehat{BCH}=45^\circ$, $\widehat{CAH}=30^\circ$, $\triangle AHC$ và $\triangle BCH$ vuông tại $H$). Điền (Đ) cho câu trả lời đúng, (S) cho câu trả lời sai.
	\choiceTF
	{\True $\widehat{ACH}=60^{\circ}$}
	{\True $AH=30\sqrt{3}\,\text{m}$}
	{ $BH=15\sqrt{2}\,\text{m}$}
	{\True Khoảng cách giữa hai điểm $A$ và $B$ xấp xỉ $21,96\,\text{m}$}
}
{
	% TODO: \usepackage{graphicx} required
	\includegraphics[scale=0.5]{C:/texstudio-man/Anh/{7AA425F8-AB5A-4DB4-BEA2-DC1A6CA69FA1}}
	
}
	\loigiai{
	\begin{itemchoice}
	\itemch Xét tam giác $AHC$ vuông tại $H$, ta có: $\widehat{ACH} = 90^{\circ} - \widehat{CAH} = 90^{\circ} - 30^{\circ} = 60^{\circ}$. \\ Vậy mệnh đề đúng.
	\itemch Nên $AH=CH \cdot \tan \widehat{ACH} = CH \cdot \tan 60^{\circ}=30\sqrt{3}\,\text{m}$. \\ Vậy mệnh đề đúng.
	\itemch Xét tam giác $BCH$ vuông tại $H$, ta có: $BH=CH \cdot \tan \widehat{BCH}=30 \cdot \tan 45^{\circ}=30 \cdot 1 = 30\,\text{m}$. \\ Vậy mệnh đề sai.
	\itemch Suy ra $AB=AH-BH=30\sqrt{3}-30 \approx 30 \cdot 1,732 - 30 = 51,96 - 30 = 21,96\,\text{m}$. Vậy khoảng cách giữa hai điểm $A$ và $B$ xấp xỉ $21,96\,\text{m}$. \\ Vậy mệnh đề đúng.
	\end{itemchoice}
	}
\end{ex}
%%==========Câu 153
\begin{ex}
	\immini[thm]
	{
		Hai địa điểm $A$ và $B$ ở hai bên hồ nước. Biết rằng các khoảng cách từ một điểm $C$ đến $A$ và đến $B$ là $CA=80\,\text{m},$ $CB=120\,\text{m}$ và $\widehat{ACB}=120^{\circ}$, $AH$ là đường cao của tam giác $ABC$ ($H$ thuộc tia đối của tia $CB$). Điền (Đ) cho câu trả lời đúng, (S) cho câu trả lời sai.
	\choiceTF
	{ $HC=50\,\text{m}$}
	{ $AH=50\sqrt{3}\,\text{m}$}
	{\True $BH=160\,\text{m}$}
	{\True $AB=40\sqrt{19}\,\text{m}$}
	}
	{
		% TODO: \usepackage{graphicx} required
		\includegraphics[scale=0.5]{C:/texstudio-man/Anh/{525E6E9B-0452-46F3-8FDB-C13407EAB8EA}}
		
	}
	\loigiai{
	\begin{itemchoice}
	\itemch Xét tam giác $ACH$ vuông tại $H$ có $\widehat{ACH}=180^{\circ} - \widehat{ACB} = 180^{\circ}-120^{\circ}=60^{\circ}$ nên $HC=AC \cdot \cos \widehat{ACH} = 80 \cdot \cos 60^{\circ} = 80 \cdot \dfrac{1}{2}=40\,\text{m}$. \\ Vậy mệnh đề sai.
	\itemch $AH=AC \cdot \sin \widehat{ACH} = 80 \cdot \sin 60^{\circ} = 80 \cdot \dfrac{\sqrt{3}}{2} = 40\sqrt{3}\,\text{m}$. \\ Vậy mệnh đề sai.
	\itemch Từ đó $BH=BC+CH=120+40=160\,\text{m}$. \\ Vậy mệnh đề đúng.
	\itemch Xét tam giác $AHB$ vuông tại $H$ có $AB^{2}=AH^{2}+BH^{2}=(40\sqrt{3})^{2}+160^{2}=1600 \cdot 3 + 25600 = 4800+25600=30400$. Hay $AB=\sqrt{30400} = \sqrt{1600 \cdot 19} = 40\sqrt{19}\,\text{m}$. \\ Vậy mệnh đề đúng.
	\end{itemchoice}
	}
\end{ex}
%%==========Câu 154
\begin{ex}
	\immini[thm]
	{
	Kèo của một mái nhà dạng tam giác cân $ABC$ (cân tại $A$). Biết $BC=4,2\,\text{m}$, chiều cao $AH=1,7\,\text{m}.$ Góc $\alpha = \widehat{B}$. Điền (Đ) cho câu trả lời đúng, (S) cho câu trả lời sai.
	\choiceTF
	{\True $BH=2,1\,\text{m}$}
	{\True $\alpha \approx 39^{\circ}$}
	{ $AB \approx 2,3\,\text{m}$}
	{\True $HD \approx 1,3\,\text{m}$}
	}
	{
			\includegraphics[scale=0.5]{C:/texstudio-man/Anh/{D695FEEE-C26B-428B-A126-9E8AD4519EDA}}
		
	}
	\loigiai{
	\begin{itemchoice}
	\itemch $AH$ là đường cao ứng với cạnh đáy của tam giác cân $ABC$ nên: $HB=HC=\dfrac{BC}{2}=\dfrac{4,2}{2}=2,1\,\text{m}$. \\ Vậy mệnh đề đúng.
	\itemch Xét tam giác $ABH$ vuông tại $H$ ta có $\tan \alpha=\dfrac{AH}{BH}=\dfrac{1,7}{2,1} \approx 0,8095$. Suy ra $\alpha\approx 39^{\circ}$. \\ Vậy mệnh đề đúng.
	\itemch Xét tam giác $ABH$ vuông tại $H$ ta có $AB=\sqrt{AH^{2}+HB^{2}}=\sqrt{1,7^{2}+2,1^{2}}=\sqrt{2,89+4,41}=\sqrt{7,3} \approx 2,70\,\text{m}$. \\ Vậy mệnh đề sai.
	\itemch Xét tam giác $DBH$ vuông tại $D$ (với $D$ là hình chiếu của $H$ trên $AB$), ta có $HD=HB \cdot \sin B \approx 2,1 \cdot \sin 39^{\circ} \approx 2,1 \cdot 0,6293 \approx 1,32\,\text{m} \approx 1,3\,\text{m}$. \\ Vậy mệnh đề đúng.
	\end{itemchoice}
	}
\end{ex}
%%==========Câu 155
\begin{ex}
\immini[thm]
{
	Hai con thuyền $P$ và $Q$ cách nhau $300\,\text{m}$ và thẳng hàng với chân $B$ của tháp hải đăng $AB$ trên bờ biển. Từ $P$ và $Q$, người ta nhìn thấy tháp hải đăng dưới các góc $\widehat{BPA}=15^{\circ}$ và $\widehat{BQA}=40^{\circ}$. Đặt $h=AB$ là chiều cao của tháp hải đăng. Điền (Đ) cho câu trả lời đúng, (S) cho câu trả lời sai.
	\choiceTF
	{\True $BQ=\dfrac{h}{\tan 40^{\circ}}$}
	{ $BP=\dfrac{h}{\sin 15^{\circ}}$}
	{ $h=\dfrac{300}{\dfrac{1}{\sin 15^{\circ}}-\dfrac{1}{\tan 40^{\circ}}}$}
	{\True $h \approx 118,1\,\text{m}$}
}
{
	% TODO: \usepackage{graphicx} required
	\includegraphics[scale=0.45]{C:/texstudio-man/Anh/{ED6F3BB9-957C-4591-A036-36F834045C89}}
	
}
	\loigiai{
	\begin{itemchoice}
	\itemch Xét tam giác $BQA$ vuông tại $B$, ta có: $\tan \widehat{BQA}=\dfrac{AB}{BQ}$ hay $BQ=\dfrac{AB}{\tan \widehat{BQA}}=\dfrac{h}{\tan 40^{\circ}}$. \\ Vậy mệnh đề đúng.
	\itemch Xét tam giác $BPA$ vuông tại $B$, ta có: $\tan \widehat{BPA}=\dfrac{AB}{BP}$ hay $BP=\dfrac{AB}{\tan \widehat{BPA}}=\dfrac{h}{\tan 15^{\circ}}$. \\ Vậy mệnh đề sai.
	\itemch Ta có: $BP-BQ=PQ$. Suy ra $\dfrac{h}{\tan 15^{\circ}}-\dfrac{h}{\tan 40^{\circ}}=300 \Rightarrow h \left(\dfrac{1}{\tan 15^{\circ}}-\dfrac{1}{\tan 40^{\circ}}\right)=300 \Rightarrow h=\dfrac{300}{\dfrac{1}{\tan 15^{\circ}}-\dfrac{1}{\tan 40^{\circ}}}$. \\ Vậy mệnh đề sai.
	\itemch $h=\dfrac{300}{\dfrac{1}{\tan 15^{\circ}}-\dfrac{1}{\tan 40^{\circ}}} \approx \dfrac{300}{\dfrac{1}{0,2679}-\dfrac{1}{0,8391}} \approx \dfrac{300}{3,7320-1,1918} \approx \dfrac{300}{2,5402} \approx 118,10\,\text{m}$. \\ Vậy mệnh đề đúng.
	\end{itemchoice}
	}
\end{ex}
%%==========Câu 156
\begin{ex}
	\immini[thm]
	{
	Mặt cắt ngang của một chân đập ngăn nước có dạng hình thang $ABCD$. Chiều rộng của mặt trên $AB$ của đập là $4\,\text{m}.$ Độ dốc của sườn $AD$, tức là $\tan D=1,25$, độ dốc của sườn $BC$, tức là $\tan C=1,5$. Chiều cao $AH$ (và $BK$) của đập là $4,5\,\text{m}.$ ($AH \perp DC, BK \perp DC$). Điền (Đ) cho câu trả lời đúng, (S) cho câu trả lời sai.
	\choiceTF
	{\True $DH=3,6\,\text{m}$}
	{ $CD=6,6\,\text{m}$}
	{ $AD=33,21\,\text{m}$}
	{\True $BC=\dfrac{3\sqrt{13}}{2}\,\text{m}$}
	}
	{
		% TODO: \usepackage{graphicx} required
		\includegraphics[scale=0.5]{C:/texstudio-man/Anh/{E34A29B8-D48E-41B8-9E13-77D88F5745BC}}
		
	}
	\loigiai{
	\begin{itemchoice}
	\itemch Xét tam giác $ADH$ vuông tại $H$, ta có: $\tan D=\dfrac{AH}{DH}$ hay $DH=\dfrac{AH}{\tan D}=\dfrac{4,5}{1,25}=3,6\,\text{m}$. \\ Vậy mệnh đề đúng.
	\itemch Kẻ $BK \perp DC$ tại $K$. $ABKH$ là hình chữ nhật nên $HK=AB=4\,\text{m}$ và $BK=AH=4,5\,\text{m}$. Xét tam giác $BCK$ vuông tại $K$, ta có:\\
	$\tan C=\dfrac{BK}{CK}$ hay $CK=\dfrac{BK}{\tan C}=\dfrac{4,5}{1,5}=3\,\text{m}$.\\
	Có: $DC=DH+HK+CK=3,6+4+3=10,6\,\text{m}$. \\ Vậy mệnh đề sai.
	\itemch Xét tam giác $ADH$ vuông tại $H$, ta có:\\ $AD^{2}=AH^{2}+DH^{2}=4,5^{2}+3,6^{2}=20,25+12,96=33,21$. Hay $AD=\sqrt{33,21}\,\text{m}$. (Lưu ý: PDF ghi (c sai) cho mệnh đề $AD=\sqrt{33,21}\,\text{m}$). \\ Vậy mệnh đề sai.
	\itemch Xét tam giác $BCK$ vuông tại $K$, ta có: $BC^{2}=BK^{2}+CK^{2}=4,5^{2}+3^{2}=20,25+9=29,25$. Hay $BC=\sqrt{29,25}=\sqrt{\dfrac{117}{4}}=\dfrac{\sqrt{9 \cdot 13}}{2}=\dfrac{3\sqrt{13}}{2}\,\text{m}$. \\ Vậy mệnh đề đúng.
	\end{itemchoice}
	}
\end{ex}
\Closesolutionfile{ans}
\indapanTF{3}{ans/ans9K1-12-DS-1112}
